% =========================================================================
\section{Introducción: el conductor perfecto}\label{sec:intro}
% =========================================================================
Hasta ahora todo nuestro estudio se ha desarrollado en el vacío: dadas las cargas, la ley de Coulomb, la ley de Gauss y el potencial
permiten calcular el campo~\cite{TEMA1}, y el desarrollo multipolar permite aproximarlo a grandes distancias~\cite{TEMA2}. En este tema
damos el primer paso hacia la materia. Los materiales se dividen, a efectos electrostáticos, en \emph{conductores} y \emph{dieléctricos};
aquí estudiamos los primeros. La novedad es profunda: en un conductor la distribución de carga \emph{no es un dato}, sino una incógnita que
se ajusta sola hasta alcanzar el equilibrio. Veremos que esa autorregulación impone condiciones tan fuertes que, paradójicamente, simplifica
el problema: basta conocer la geometría y unas pocas cargas o potenciales para determinarlo todo.

Seguimos la estructura de los apuntes manuscritos y de teoría de la asignatura~\cite{APUNTES-T3}, cuyas figuras~\cite{FIGURAS-T3} hemos
recalculado y redibujado; para ampliar, pueden consultarse los textos clásicos~\cite{GRIFFITHS,WANGSNESS,VICTORIANO,PURCELL,FEYNMAN}.

% -------------------------------------------------------------------------
\subsection{Conductor perfecto y equilibrio electrostático}\label{sec:conductor-perfecto}
% -------------------------------------------------------------------------
De forma genérica, un material conductor es un cuerpo en el que las cargas eléctricas pueden moverse libremente bajo la acción de un campo
eléctrico. El ejemplo más común son los metales (cobre, plata, oro, aluminio\ldots): en el cobre hay aproximadamente un electrón libre por
átomo. Pero no son los únicos: en el agua salada, por ejemplo, lo que se desplaza son los iones.

\begin{definicion}[Conductor perfecto]\label{def:conductor}
Un \emph{conductor perfecto} es un cuerpo material, que ocupa una región cerrada y acotada $\Vc\subset\R^{3}$ de frontera $\Sc=\partial\Vc$,
que cumple:
\begin{enumerate}[label=(\roman*)]
  \item las cargas libres se mueven por su interior sin oposición alguna (se desprecia la agitación térmica);
  \item admite un número ilimitado de cargas libres: dispone de una reserva inagotable de carga de ambos signos.
\end{enumerate}
\end{definicion}

\begin{observacion}[Una idealización excelente]\label{obs:idealizacion}
Ningún conductor real es perfecto: los metales tienen una resistividad no nula y una densidad de electrones libres finita
($n\simeq8{,}5\cdot10^{28}\,\mathrm{m^{-3}}$ en el cobre). Sin embargo, la aproximación funciona extraordinariamente bien en electrostática,
donde solo interesa el estado final de equilibrio. Las cargas que hay que redistribuir para anular un campo externo típico son una fracción
insignificante de las disponibles, y la redistribución es prácticamente instantánea: en un conductor óhmico de conductividad $\sigma_c$ cualquier
exceso de carga en el interior decae como $\eul^{-t/\tau}$, con $\tau=\eo/\sigma_c$, que para el cobre ($\sigma_c\simeq5{,}96\cdot10^{7}\,\mathrm{S/m}$)
vale $\tau\simeq1{,}5\cdot10^{-19}\,\mathrm{s}$~\cite{GRIFFITHS}.
\end{observacion}

\begin{definicion}[Equilibrio electrostático]\label{def:equilibrio}
Un conductor está en \emph{equilibrio electrostático} cuando sus cargas libres están en reposo a escala macroscópica, es decir, elemento de
volumen a elemento de volumen: no hay ningún flujo neto de carga en su interior ni sobre su superficie.
\end{definicion}

\begin{proposicion}[Campo nulo en el interior]\label{prop:campo-nulo}
En un conductor perfecto en equilibrio electrostático el campo eléctrico es nulo en todo punto interior:
\begin{equation}\label{eq:campo-nulo}
  \resalta{\vE(\vc{r})=\vc{0}\qquad\forall\,\vc{r}\in\Vc\setminus\Sc .}
\end{equation}
\end{proposicion}
\begin{proof}
Supongamos que existe un punto interior $\vc{r}_0$ con $\vE(\vc{r}_0)\neq\vc{0}$. Las cargas libres situadas en un entorno de $\vc{r}_0$ (que existen
por la hipótesis~(ii) de la \cref{def:conductor}) experimentan la fuerza $\dif\vF=\dif q\,\vE$, que por la hipótesis~(i) no encuentra
oposición: se pondrían en movimiento, lo que contradice la \cref{def:equilibrio}. Luego $\vE=\vc{0}$ en todo el interior.
\end{proof}

La demostración dice \emph{qué} ocurre en el equilibrio, pero no \emph{cómo} se llega a él. Supongamos que introducimos un conductor
descargado en un campo externo uniforme $\vE_0$ (\cref{fig:placa}). En los primeros instantes el campo penetra en el material y empuja a los
electrones libres en sentido opuesto a $\vE_0$: se acumulan en la cara por la que entran las líneas, que queda con carga negativa, mientras
que la cara opuesta, que pierde electrones, queda con carga positiva. Esta \emph{carga inducida} crea, en el interior, un campo $\vE_{\mathrm{ind}}$
dirigido de la carga positiva a la negativa, es decir, \emph{opuesto} a $\vE_0$. La redistribución continúa mientras quede campo neto en el
interior, y se detiene exactamente cuando
\[
  \vE=\vE_0+\vE_{\mathrm{ind}}=\vc{0}\qquad\text{en el interior.}
\]
La clave es, pues, que el campo creado por la redistribución de carga inducida por el campo externo se opone a dicho campo y acaba
cancelándolo. Fuera del conductor, en cambio, el campo inducido \emph{no} cancela al externo: lo deforma, de modo que las líneas de campo
terminan en las cargas negativas y vuelven a nacer en las positivas.

% =========================================================================
%  Figura: placa conductora en un campo externo uniforme.
%  (a) instante inicial: el campo penetra y desplaza a los electrones;
%  (b) equilibrio: la carga inducida crea E_ind = -E_0 en el interior.
% =========================================================================
\begin{figure}[t]
\centering
\begin{tikzpicture}[x=1cm, y=1cm, font=\small]
  % ---------------------------------------------------------------- (a)
  \begin{scope}
    \fill[conductorplano, fill=metal!45!white, draw=azulprof] (-0.8,-1.7) rectangle (0.8,1.7);
    \foreach \y in {-1.3,-0.65,0,0.65,1.3}{
      \draw[azulacero, line width=0.6pt] (-2.5,\y) -- (-0.8,\y);
      \draw[azulacero!55, line width=0.6pt, densely dashed] (-0.8,\y) -- (0.8,\y);
      \draw[azulacero, line width=0.6pt, -{Stealth[length=4pt]}] (0.8,\y) -- (2.5,\y);
      \draw[azulacero, line width=0.6pt, -{Stealth[length=4pt]}] (-2.5,\y) -- (-1.4,\y);}
    \foreach \y in {-0.98,-0.33,0.33,0.98}{
      \draw[-{Stealth[length=3.5pt]}, carganeg, line width=0.7pt] (0.3,\y) -- (-0.2,\y);
      \node[font=\scriptsize, text=carganeg, anchor=west, inner sep=1pt] at (0.32,\y) {$e^{-}$};}
    \node[above, text=azulprof] at (-1.6,1.75) {$\vE_0$};
    \node[text=azulprof, font=\footnotesize] at (0,-2.05) {conductor};
    \node[font=\small] at (0,-2.65) {(a) Instante inicial};
  \end{scope}
  % ---------------------------------------------------------------- (b)
  \begin{scope}[xshift=6.6cm]
    \fill[metal!45!white, draw=azulprof, line width=0.6pt] (-0.8,-1.7) rectangle (0.8,1.7);
    \foreach \y in {-1.3,-0.65,0,0.65,1.3}{
      \draw[azulacero, line width=0.6pt, -{Stealth[length=4pt]}] (-2.5,\y) -- (-1.4,\y);
      \draw[azulacero, line width=0.6pt] (-1.4,\y) -- (-0.87,\y);
      \draw[azulacero, line width=0.6pt, -{Stealth[length=4pt]}] (0.87,\y) -- (2.5,\y);}
    \foreach \y in {-1.5,-1.1,...,1.55}{
      \node[signoneg] at (-0.66,\y) {$-$};
      \node[signopos] at (0.66,\y) {$+$};}
    \draw[-{Stealth[length=3.5pt]}, cargapos, line width=0.8pt] (0.32,0.5) -- (-0.32,0.5);
    \node[font=\scriptsize, text=cargapos] at (0,0.74) {$\vE_{\mathrm{ind}}$};
    \draw[-{Stealth[length=3.5pt]}, azulacero, line width=0.8pt] (-0.32,-0.5) -- (0.32,-0.5);
    \node[font=\scriptsize, text=azulacero] at (0,-0.74) {$\vE_{0}$};
    \node[font=\footnotesize, text=azulprof, fill=white, inner sep=1.5pt, rounded corners=1pt] at (0,0.02) {$\vE=\vc{0}$};
    \node[above, text=azulprof] at (-1.6,1.75) {$\vE_0$};
    \node[text=carganeg, font=\footnotesize] at (-0.95,-2.05) {$\sigma<0$};
    \node[text=cargapos, font=\footnotesize] at (0.95,-2.05) {$\sigma>0$};
    \node[font=\small] at (0,-2.65) {(b) Equilibrio electrostático};
  \end{scope}
\end{tikzpicture}
\caption{Conductor neutro en un campo externo uniforme $\vE_0$. (a) Al principio el campo penetra en el metal (trazos) y arrastra a los
electrones libres en sentido opuesto a $\vE_0$. (b) En el equilibrio, la carga inducida en las caras crea en el interior un campo
$\vE_{\mathrm{ind}}=-\vE_0$ que anula al externo: las líneas terminan en la cara negativa y renacen en la positiva.}
\label{fig:placa}
\end{figure}


La \cref{fig:uniforme} muestra el resultado de ese proceso para dos conductores neutros: una esfera, cuya solución es analítica, y un
conductor de forma irregular, que hemos resuelto numéricamente. En ambos casos el campo exterior se deforma para llegar
\emph{perpendicularmente} a la superficie, la carga inducida se concentra en las zonas más curvadas y el interior queda libre de campo.
Las secciones siguientes demuestran cada uno de estos rasgos.

% =========================================================================
%  Figura «tipo Mathematica»: conductores neutros en un campo uniforme.
%  (a) esfera (solución analítica, plano meridiano): líneas de campo como
%      curvas de nivel de la función de Stokes Psi = (E0/2) y^2 (1 + 2a^3/r^3);
%  (b) conductor irregular (ecuación de Laplace 2D por diferencias finitas).
%  Datos: figuras/datos/esfera_uniforme_* e irregular_uniforme_*.
% =========================================================================
\begin{figure}[t]
\centering
\begin{tikzpicture}[font=\small]
  % ---------------------------------------------------------------- (a)
  \begin{scope}[x=1.04cm, y=1.04cm]
    \fill[azulfondo!60!white] (-2.6,-2.6) rectangle (2.6,2.6);
    \begin{scope}
      \clip (-2.6,-2.6) rectangle (2.6,2.6);
      \tikzset{numEQ/.style={azulcielo, line width=0.35pt, densely dashed}}
      % Esfera en campo uniforme - equipotenciales
% Generado por scripts/generar_datos.py (no editar a mano).
\EQ{(2.454,-2.600) (2.455,-2.567) (2.456,-2.535) (2.457,-2.502) (2.458,-2.470) (2.459,-2.438) (2.460,-2.405) (2.462,-2.373) (2.463,-2.340) (2.464,-2.308) (2.465,-2.275) (2.467,-2.243) (2.468,-2.210) (2.469,-2.178) (2.471,-2.145) (2.472,-2.112) (2.473,-2.080) (2.475,-2.047) (2.476,-2.015) (2.478,-1.982) (2.479,-1.950) (2.480,-1.917) (2.482,-1.885) (2.484,-1.853) (2.485,-1.820) (2.487,-1.788) (2.488,-1.755) (2.490,-1.723) (2.491,-1.690) (2.493,-1.657) (2.495,-1.625) (2.496,-1.593) (2.498,-1.560) (2.499,-1.528) (2.501,-1.495) (2.503,-1.462) (2.504,-1.430) (2.506,-1.397) (2.508,-1.365) (2.509,-1.333) (2.511,-1.300) (2.513,-1.268) (2.514,-1.235) (2.516,-1.202) (2.518,-1.170) (2.519,-1.137) (2.521,-1.105) (2.522,-1.073) (2.524,-1.040) (2.526,-1.007) (2.527,-0.975) (2.529,-0.942) (2.530,-0.910) (2.532,-0.877) (2.533,-0.845) (2.534,-0.812) (2.536,-0.780) (2.537,-0.747) (2.538,-0.715) (2.540,-0.682) (2.541,-0.650) (2.542,-0.617) (2.543,-0.585) (2.544,-0.552) (2.545,-0.520) (2.546,-0.487) (2.547,-0.455) (2.548,-0.422) (2.549,-0.390) (2.549,-0.357) (2.550,-0.325) (2.551,-0.292) (2.551,-0.260) (2.552,-0.228) (2.552,-0.195) (2.553,-0.163) (2.553,-0.130) (2.553,-0.097) (2.553,-0.065) (2.553,-0.032) (2.553,0.000) (2.553,0.033) (2.553,0.065) (2.553,0.098) (2.553,0.130) (2.553,0.163) (2.552,0.195) (2.552,0.228) (2.551,0.260) (2.551,0.292) (2.550,0.325) (2.549,0.358) (2.549,0.390) (2.548,0.423) (2.547,0.455) (2.546,0.488) (2.545,0.520) (2.544,0.553) (2.543,0.585) (2.542,0.618) (2.541,0.650) (2.540,0.683) (2.538,0.715) (2.537,0.748) (2.536,0.780) (2.534,0.812) (2.533,0.845) (2.532,0.878) (2.530,0.910) (2.529,0.943) (2.527,0.975) (2.526,1.008) (2.524,1.040) (2.522,1.073) (2.521,1.105) (2.519,1.138) (2.518,1.170) (2.516,1.203) (2.514,1.235) (2.513,1.268) (2.511,1.300) (2.509,1.333) (2.508,1.365) (2.506,1.398) (2.504,1.430) (2.503,1.462) (2.501,1.495) (2.499,1.528) (2.498,1.560) (2.496,1.593) (2.495,1.625) (2.493,1.658) (2.491,1.690) (2.490,1.723) (2.488,1.755) (2.487,1.788) (2.485,1.820) (2.484,1.853) (2.482,1.885) (2.480,1.917) (2.479,1.950) (2.478,1.983) (2.476,2.015) (2.475,2.047) (2.473,2.080) (2.472,2.113) (2.471,2.145) (2.469,2.178) (2.468,2.210) (2.467,2.243) (2.465,2.275) (2.464,2.308) (2.463,2.340) (2.462,2.373) (2.460,2.405) (2.459,2.438) (2.458,2.470) (2.457,2.501) (2.456,2.535) (2.455,2.568) (2.454,2.600)}
\EQ{(2.056,-2.600) (2.058,-2.567) (2.059,-2.535) (2.060,-2.502) (2.062,-2.470) (2.063,-2.438) (2.065,-2.405) (2.066,-2.373) (2.068,-2.340) (2.069,-2.308) (2.071,-2.275) (2.073,-2.243) (2.074,-2.210) (2.076,-2.178) (2.078,-2.145) (2.080,-2.112) (2.082,-2.080) (2.084,-2.047) (2.086,-2.015) (2.087,-1.982) (2.090,-1.950) (2.092,-1.917) (2.094,-1.885) (2.096,-1.853) (2.098,-1.820) (2.100,-1.788) (2.102,-1.755) (2.105,-1.723) (2.107,-1.690) (2.109,-1.657) (2.112,-1.625) (2.114,-1.593) (2.116,-1.560) (2.119,-1.528) (2.121,-1.495) (2.124,-1.462) (2.126,-1.430) (2.129,-1.397) (2.131,-1.365) (2.134,-1.333) (2.137,-1.300) (2.139,-1.268) (2.142,-1.235) (2.144,-1.202) (2.147,-1.170) (2.149,-1.137) (2.152,-1.105) (2.155,-1.073) (2.157,-1.040) (2.160,-1.007) (2.162,-0.975) (2.165,-0.943) (2.167,-0.910) (2.169,-0.877) (2.172,-0.845) (2.174,-0.812) (2.176,-0.780) (2.178,-0.747) (2.180,-0.715) (2.183,-0.682) (2.185,-0.650) (2.186,-0.617) (2.188,-0.585) (2.190,-0.552) (2.192,-0.520) (2.193,-0.487) (2.195,-0.455) (2.196,-0.422) (2.198,-0.390) (2.199,-0.357) (2.200,-0.325) (2.201,-0.292) (2.202,-0.260) (2.203,-0.228) (2.204,-0.197) (2.204,-0.163) (2.205,-0.130) (2.205,-0.097) (2.205,-0.065) (2.206,-0.032) (2.206,0.000) (2.206,0.033) (2.205,0.065) (2.205,0.098) (2.205,0.130) (2.204,0.163) (2.204,0.195) (2.203,0.228) (2.202,0.260) (2.201,0.292) (2.200,0.325) (2.199,0.358) (2.198,0.390) (2.196,0.423) (2.195,0.455) (2.193,0.488) (2.192,0.520) (2.190,0.553) (2.188,0.585) (2.186,0.618) (2.185,0.650) (2.183,0.683) (2.180,0.715) (2.178,0.748) (2.176,0.780) (2.174,0.812) (2.172,0.845) (2.169,0.878) (2.167,0.910) (2.165,0.943) (2.162,0.975) (2.160,1.008) (2.157,1.040) (2.155,1.073) (2.152,1.105) (2.149,1.138) (2.147,1.170) (2.144,1.203) (2.142,1.235) (2.139,1.268) (2.137,1.300) (2.134,1.333) (2.131,1.365) (2.129,1.398) (2.126,1.430) (2.124,1.462) (2.121,1.495) (2.119,1.526) (2.116,1.560) (2.114,1.593) (2.112,1.625) (2.109,1.658) (2.107,1.690) (2.105,1.723) (2.102,1.755) (2.100,1.788) (2.098,1.820) (2.096,1.853) (2.094,1.885) (2.092,1.917) (2.090,1.950) (2.087,1.983) (2.086,2.015) (2.084,2.047) (2.082,2.080) (2.080,2.113) (2.078,2.145) (2.076,2.178) (2.074,2.210) (2.073,2.243) (2.071,2.275) (2.069,2.308) (2.068,2.340) (2.066,2.373) (2.065,2.405) (2.063,2.438) (2.062,2.470) (2.061,2.502) (2.059,2.535) (2.058,2.568) (2.056,2.600)}
\EQ{(1.657,-2.600) (1.658,-2.567) (1.660,-2.535) (1.661,-2.502) (1.663,-2.470) (1.665,-2.438) (1.667,-2.405) (1.668,-2.373) (1.670,-2.340) (1.672,-2.308) (1.674,-2.275) (1.676,-2.243) (1.679,-2.210) (1.681,-2.178) (1.683,-2.145) (1.685,-2.112) (1.688,-2.080) (1.690,-2.047) (1.693,-2.015) (1.696,-1.982) (1.698,-1.950) (1.701,-1.917) (1.704,-1.885) (1.707,-1.853) (1.710,-1.820) (1.713,-1.788) (1.716,-1.755) (1.719,-1.723) (1.723,-1.691) (1.726,-1.657) (1.729,-1.625) (1.733,-1.593) (1.737,-1.560) (1.740,-1.528) (1.744,-1.495) (1.748,-1.462) (1.752,-1.430) (1.755,-1.397) (1.759,-1.365) (1.763,-1.333) (1.767,-1.300) (1.771,-1.268) (1.776,-1.235) (1.780,-1.202) (1.784,-1.170) (1.788,-1.137) (1.792,-1.105) (1.796,-1.073) (1.800,-1.040) (1.804,-1.007) (1.809,-0.975) (1.813,-0.942) (1.817,-0.910) (1.821,-0.877) (1.824,-0.845) (1.828,-0.812) (1.832,-0.780) (1.836,-0.747) (1.839,-0.715) (1.843,-0.682) (1.846,-0.652) (1.849,-0.617) (1.853,-0.586) (1.856,-0.552) (1.859,-0.520) (1.861,-0.487) (1.864,-0.455) (1.866,-0.422) (1.869,-0.390) (1.871,-0.357) (1.873,-0.325) (1.875,-0.292) (1.876,-0.260) (1.878,-0.228) (1.879,-0.195) (1.880,-0.163) (1.881,-0.130) (1.881,-0.097) (1.882,-0.065) (1.882,-0.032) (1.882,0.000) (1.882,0.033) (1.882,0.065) (1.881,0.098) (1.881,0.130) (1.880,0.163) (1.879,0.195) (1.878,0.228) (1.876,0.260) (1.875,0.292) (1.873,0.325) (1.871,0.358) (1.869,0.390) (1.866,0.423) (1.864,0.455) (1.861,0.488) (1.859,0.520) (1.856,0.553) (1.853,0.585) (1.849,0.618) (1.846,0.650) (1.843,0.683) (1.840,0.713) (1.836,0.748) (1.832,0.780) (1.828,0.812) (1.824,0.845) (1.821,0.878) (1.817,0.910) (1.813,0.943) (1.809,0.975) (1.804,1.008) (1.800,1.039) (1.796,1.073) (1.792,1.105) (1.788,1.138) (1.784,1.170) (1.780,1.203) (1.776,1.235) (1.771,1.268) (1.767,1.300) (1.763,1.333) (1.759,1.365) (1.755,1.398) (1.752,1.430) (1.748,1.462) (1.744,1.495) (1.740,1.528) (1.737,1.560) (1.733,1.593) (1.729,1.625) (1.726,1.658) (1.723,1.690) (1.719,1.723) (1.716,1.755) (1.713,1.788) (1.710,1.820) (1.707,1.853) (1.704,1.885) (1.701,1.917) (1.698,1.950) (1.696,1.983) (1.693,2.015) (1.690,2.047) (1.688,2.080) (1.685,2.112) (1.683,2.145) (1.681,2.178) (1.679,2.210) (1.676,2.243) (1.674,2.275) (1.672,2.308) (1.670,2.340) (1.668,2.373) (1.667,2.405) (1.665,2.438) (1.663,2.470) (1.661,2.502) (1.660,2.535) (1.658,2.568) (1.657,2.600)}
\EQ{(1.252,-2.600) (1.254,-2.567) (1.255,-2.535) (1.257,-2.502) (1.259,-2.470) (1.261,-2.438) (1.263,-2.405) (1.265,-2.373) (1.267,-2.340) (1.269,-2.308) (1.272,-2.275) (1.274,-2.243) (1.277,-2.210) (1.279,-2.178) (1.282,-2.145) (1.285,-2.112) (1.288,-2.080) (1.291,-2.047) (1.294,-2.015) (1.298,-1.982) (1.301,-1.950) (1.305,-1.917) (1.308,-1.885) (1.312,-1.853) (1.316,-1.820) (1.320,-1.788) (1.325,-1.755) (1.329,-1.723) (1.334,-1.690) (1.338,-1.657) (1.343,-1.625) (1.348,-1.593) (1.354,-1.560) (1.359,-1.527) (1.365,-1.495) (1.370,-1.462) (1.376,-1.430) (1.382,-1.397) (1.388,-1.365) (1.394,-1.333) (1.401,-1.300) (1.407,-1.268) (1.414,-1.235) (1.420,-1.202) (1.427,-1.170) (1.434,-1.137) (1.441,-1.105) (1.448,-1.073) (1.454,-1.040) (1.461,-1.007) (1.468,-0.975) (1.475,-0.942) (1.482,-0.910) (1.489,-0.878) (1.495,-0.846) (1.502,-0.813) (1.508,-0.781) (1.514,-0.747) (1.521,-0.715) (1.527,-0.682) (1.532,-0.650) (1.538,-0.617) (1.543,-0.585) (1.548,-0.552) (1.553,-0.520) (1.558,-0.487) (1.563,-0.455) (1.567,-0.424) (1.571,-0.390) (1.574,-0.357) (1.578,-0.325) (1.581,-0.292) (1.583,-0.260) (1.586,-0.228) (1.588,-0.195) (1.590,-0.163) (1.591,-0.130) (1.592,-0.097) (1.593,-0.065) (1.594,-0.032) (1.594,0.000) (1.594,0.033) (1.593,0.065) (1.592,0.098) (1.591,0.130) (1.590,0.163) (1.588,0.195) (1.586,0.228) (1.583,0.260) (1.581,0.292) (1.578,0.325) (1.574,0.358) (1.571,0.390) (1.567,0.423) (1.563,0.455) (1.558,0.488) (1.554,0.519) (1.548,0.553) (1.543,0.585) (1.538,0.618) (1.532,0.650) (1.527,0.683) (1.521,0.712) (1.515,0.747) (1.508,0.780) (1.502,0.812) (1.495,0.845) (1.489,0.878) (1.482,0.909) (1.476,0.940) (1.469,0.971) (1.463,1.001) (1.456,1.033) (1.450,1.063) (1.443,1.094) (1.437,1.125) (1.430,1.156) (1.424,1.187) (1.417,1.219) (1.411,1.251) (1.405,1.281) (1.398,1.313) (1.392,1.346) (1.386,1.378) (1.380,1.411) (1.374,1.443) (1.368,1.476) (1.362,1.508) (1.357,1.541) (1.352,1.573) (1.346,1.606) (1.341,1.638) (1.336,1.671) (1.332,1.703) (1.327,1.736) (1.323,1.768) (1.319,1.800) (1.315,1.833) (1.311,1.866) (1.307,1.898) (1.303,1.930) (1.300,1.963) (1.296,1.996) (1.293,2.028) (1.290,2.061) (1.287,2.090) (1.284,2.126) (1.281,2.158) (1.278,2.191) (1.276,2.223) (1.273,2.256) (1.271,2.288) (1.269,2.321) (1.266,2.353) (1.264,2.386) (1.262,2.418) (1.260,2.451) (1.258,2.483) (1.257,2.516) (1.255,2.548) (1.253,2.581) (1.252,2.600)}
\EQ{(0.841,-2.600) (0.843,-2.567) (0.844,-2.535) (0.846,-2.502) (0.848,-2.470) (0.849,-2.438) (0.851,-2.405) (0.853,-2.373) (0.855,-2.340) (0.857,-2.308) (0.860,-2.275) (0.862,-2.243) (0.864,-2.212) (0.867,-2.178) (0.870,-2.145) (0.873,-2.112) (0.876,-2.080) (0.879,-2.047) (0.883,-2.015) (0.887,-1.982) (0.890,-1.950) (0.894,-1.917) (0.899,-1.885) (0.903,-1.853) (0.908,-1.820) (0.913,-1.788) (0.918,-1.755) (0.924,-1.723) (0.930,-1.690) (0.936,-1.657) (0.942,-1.625) (0.949,-1.593) (0.956,-1.563) (0.964,-1.527) (0.971,-1.495) (0.980,-1.462) (0.988,-1.431) (0.997,-1.397) (1.006,-1.365) (1.016,-1.333) (1.026,-1.300) (1.036,-1.268) (1.047,-1.236) (1.058,-1.202) (1.069,-1.170) (1.079,-1.140) (1.089,-1.111) (1.101,-1.079) (1.112,-1.049) (1.124,-1.014) (1.136,-0.982) (1.148,-0.949) (1.159,-0.916) (1.170,-0.886) (1.182,-0.851) (1.194,-0.819) (1.205,-0.786) (1.215,-0.754) (1.226,-0.722) (1.235,-0.693) (1.244,-0.663) (1.254,-0.630) (1.263,-0.598) (1.272,-0.566) (1.280,-0.533) (1.288,-0.500) (1.296,-0.468) (1.303,-0.435) (1.309,-0.403) (1.315,-0.370) (1.321,-0.338) (1.326,-0.307) (1.331,-0.273) (1.335,-0.240) (1.339,-0.208) (1.342,-0.175) (1.344,-0.143) (1.346,-0.111) (1.348,-0.078) (1.349,-0.045) (1.349,-0.013) (1.349,0.020) (1.349,0.052) (1.348,0.085) (1.346,0.117) (1.344,0.150) (1.341,0.182) (1.338,0.215) (1.334,0.247) (1.330,0.280) (1.325,0.312) (1.320,0.345) (1.314,0.377) (1.308,0.409) (1.301,0.442) (1.294,0.475) (1.287,0.505) (1.278,0.540) (1.270,0.572) (1.261,0.605) (1.252,0.637) (1.242,0.670) (1.232,0.702) (1.222,0.734) (1.211,0.767) (1.200,0.800) (1.190,0.831) (1.178,0.865) (1.166,0.897) (1.155,0.930) (1.144,0.959) (1.133,0.988) (1.122,1.021) (1.110,1.053) (1.099,1.085) (1.087,1.118) (1.075,1.151) (1.064,1.183) (1.053,1.216) (1.042,1.248) (1.032,1.281) (1.022,1.313) (1.012,1.346) (1.003,1.378) (0.993,1.411) (0.985,1.443) (0.976,1.476) (0.969,1.507) (0.961,1.541) (0.953,1.573) (0.946,1.606) (0.940,1.638) (0.933,1.671) (0.927,1.703) (0.921,1.736) (0.917,1.765) (0.911,1.800) (0.906,1.833) (0.901,1.866) (0.897,1.898) (0.893,1.930) (0.889,1.963) (0.885,1.996) (0.882,2.028) (0.878,2.061) (0.875,2.093) (0.872,2.126) (0.869,2.158) (0.866,2.191) (0.864,2.223) (0.861,2.256) (0.859,2.288) (0.857,2.321) (0.854,2.353) (0.852,2.385) (0.851,2.418) (0.849,2.451) (0.847,2.483) (0.845,2.516) (0.844,2.548) (0.842,2.581) (0.841,2.600)}
\EQ{(0.423,-2.600) (0.424,-2.567) (0.425,-2.535) (0.426,-2.502) (0.427,-2.470) (0.428,-2.438) (0.429,-2.405) (0.431,-2.373) (0.432,-2.340) (0.433,-2.308) (0.435,-2.275) (0.437,-2.243) (0.438,-2.210) (0.440,-2.178) (0.442,-2.146) (0.444,-2.112) (0.446,-2.080) (0.449,-2.047) (0.451,-2.015) (0.454,-1.982) (0.457,-1.950) (0.460,-1.917) (0.463,-1.885) (0.467,-1.853) (0.471,-1.820) (0.475,-1.788) (0.480,-1.755) (0.485,-1.723) (0.490,-1.690) (0.496,-1.657) (0.502,-1.625) (0.509,-1.593) (0.516,-1.560) (0.525,-1.528) (0.533,-1.496) (0.543,-1.462) (0.553,-1.433) (0.563,-1.404) (0.575,-1.371) (0.588,-1.339) (0.602,-1.306) (0.618,-1.274) (0.634,-1.242) (0.650,-1.212) (0.666,-1.183) (0.681,-1.157) (0.697,-1.131) (0.714,-1.105) (0.730,-1.079) (0.747,-1.053) (0.764,-1.027) (0.782,-1.001) (0.799,-0.975) (0.816,-0.949) (0.833,-0.923) (0.850,-0.897) (0.867,-0.871) (0.884,-0.845) (0.900,-0.819) (0.915,-0.793) (0.931,-0.767) (0.949,-0.735) (0.967,-0.702) (0.982,-0.674) (0.997,-0.643) (1.012,-0.611) (1.027,-0.579) (1.040,-0.548) (1.051,-0.520) (1.064,-0.487) (1.075,-0.455) (1.086,-0.425) (1.096,-0.390) (1.105,-0.358) (1.113,-0.325) (1.121,-0.292) (1.128,-0.260) (1.134,-0.228) (1.139,-0.195) (1.143,-0.163) (1.146,-0.130) (1.149,-0.097) (1.151,-0.065) (1.152,-0.032) (1.153,0.000) (1.152,0.033) (1.151,0.065) (1.149,0.098) (1.146,0.130) (1.143,0.163) (1.139,0.195) (1.134,0.228) (1.128,0.260) (1.121,0.292) (1.113,0.325) (1.105,0.358) (1.096,0.390) (1.086,0.423) (1.075,0.455) (1.064,0.488) (1.053,0.516) (1.041,0.546) (1.027,0.579) (1.014,0.607) (1.001,0.635) (0.987,0.663) (0.970,0.696) (0.956,0.723) (0.938,0.754) (0.923,0.780) (0.908,0.806) (0.892,0.832) (0.875,0.858) (0.859,0.884) (0.842,0.910) (0.825,0.936) (0.808,0.962) (0.790,0.988) (0.774,1.013) (0.756,1.040) (0.739,1.066) (0.722,1.092) (0.705,1.118) (0.689,1.144) (0.674,1.170) (0.659,1.196) (0.644,1.224) (0.627,1.255) (0.611,1.287) (0.598,1.316) (0.585,1.346) (0.572,1.378) (0.560,1.411) (0.549,1.443) (0.540,1.474) (0.530,1.508) (0.521,1.540) (0.514,1.572) (0.506,1.606) (0.499,1.638) (0.493,1.671) (0.488,1.703) (0.483,1.736) (0.478,1.768) (0.473,1.800) (0.469,1.833) (0.465,1.866) (0.462,1.898) (0.459,1.930) (0.456,1.963) (0.453,1.996) (0.450,2.028) (0.448,2.061) (0.445,2.093) (0.443,2.126) (0.441,2.158) (0.439,2.191) (0.438,2.223) (0.436,2.256) (0.434,2.288) (0.433,2.321) (0.432,2.353) (0.430,2.385) (0.429,2.417) (0.428,2.451) (0.427,2.483) (0.426,2.516) (0.425,2.548) (0.424,2.581) (0.423,2.600)}
\EQ{(0.000,-2.600) (0.000,-2.567) (0.000,-2.535) (0.000,-2.502) (0.000,-2.470) (0.000,-2.438) (0.000,-2.405) (0.000,-2.373) (0.000,-2.340) (0.000,-2.307) (0.000,-2.275) (0.000,-2.243) (0.000,-2.210) (0.000,-2.178) (0.000,-2.145) (0.000,-2.113) (0.000,-2.080) (0.000,-2.047) (0.000,-2.015) (0.000,-1.982) (0.000,-1.950) (0.000,-1.918) (0.000,-1.885) (0.000,-1.853) (0.000,-1.820) (0.000,-1.788) (0.000,-1.755) (0.000,-1.723) (0.000,-1.690) (0.000,-1.657) (0.000,-1.625) (0.000,-1.593) (0.000,-1.560) (0.000,-1.528) (0.000,-1.495) (0.000,-1.462) (0.000,-1.430) (0.000,-1.398) (0.000,-1.365) (0.000,-1.333) (0.000,-1.300) (0.000,-1.268) (0.000,-1.235) (0.000,-1.202) (0.000,-1.170) (0.000,-1.138) (0.000,-1.105) (0.000,-1.073) (0.000,-1.040) (0.000,-1.008) (0.000,-1.001)}
\EQ{(0.000,1.001) (0.000,1.034) (0.000,1.066) (0.000,1.099) (0.000,1.131) (0.000,1.164) (0.000,1.196) (0.000,1.229) (0.000,1.261) (0.000,1.294) (0.000,1.326) (0.000,1.359) (0.000,1.391) (0.000,1.424) (0.000,1.456) (0.000,1.489) (0.000,1.521) (0.000,1.554) (0.000,1.586) (0.000,1.619) (0.000,1.651) (0.000,1.683) (0.000,1.716) (0.000,1.749) (0.000,1.781) (0.000,1.813) (0.000,1.846) (0.000,1.879) (0.000,1.911) (0.000,1.944) (0.000,1.976) (0.000,2.009) (0.000,2.041) (0.000,2.074) (0.000,2.106) (0.000,2.139) (0.000,2.171) (0.000,2.204) (0.000,2.236) (0.000,2.268) (0.000,2.301) (0.000,2.334) (0.000,2.366) (0.000,2.399) (0.000,2.431) (0.000,2.464) (0.000,2.496) (0.000,2.529) (0.000,2.561) (0.000,2.594) (0.000,2.600)}
\EQ{(-0.423,-2.600) (-0.424,-2.567) (-0.425,-2.535) (-0.426,-2.502) (-0.427,-2.470) (-0.428,-2.438) (-0.429,-2.405) (-0.431,-2.373) (-0.432,-2.340) (-0.433,-2.308) (-0.435,-2.275) (-0.437,-2.243) (-0.438,-2.210) (-0.440,-2.178) (-0.442,-2.146) (-0.444,-2.112) (-0.446,-2.080) (-0.449,-2.047) (-0.451,-2.015) (-0.454,-1.982) (-0.457,-1.950) (-0.460,-1.918) (-0.463,-1.885) (-0.467,-1.853) (-0.471,-1.820) (-0.475,-1.788) (-0.480,-1.755) (-0.485,-1.723) (-0.490,-1.690) (-0.496,-1.658) (-0.502,-1.625) (-0.509,-1.593) (-0.516,-1.560) (-0.525,-1.528) (-0.533,-1.496) (-0.543,-1.462) (-0.552,-1.433) (-0.563,-1.404) (-0.575,-1.371) (-0.588,-1.339) (-0.602,-1.306) (-0.617,-1.274) (-0.634,-1.241) (-0.650,-1.212) (-0.666,-1.183) (-0.681,-1.157) (-0.697,-1.131) (-0.714,-1.105) (-0.730,-1.079) (-0.747,-1.053) (-0.764,-1.027) (-0.782,-1.001) (-0.799,-0.975) (-0.816,-0.949) (-0.833,-0.923) (-0.850,-0.897) (-0.867,-0.871) (-0.884,-0.845) (-0.900,-0.819) (-0.915,-0.793) (-0.931,-0.767) (-0.949,-0.735) (-0.967,-0.702) (-0.982,-0.674) (-0.997,-0.643) (-1.012,-0.611) (-1.027,-0.579) (-1.040,-0.548) (-1.051,-0.520) (-1.064,-0.487) (-1.075,-0.455) (-1.085,-0.425) (-1.096,-0.390) (-1.105,-0.358) (-1.113,-0.325) (-1.121,-0.292) (-1.128,-0.260) (-1.134,-0.228) (-1.139,-0.195) (-1.143,-0.163) (-1.146,-0.130) (-1.149,-0.097) (-1.151,-0.065) (-1.152,-0.032) (-1.153,0.000) (-1.152,0.033) (-1.151,0.065) (-1.149,0.098) (-1.146,0.130) (-1.143,0.163) (-1.139,0.195) (-1.134,0.228) (-1.128,0.260) (-1.121,0.292) (-1.113,0.325) (-1.105,0.358) (-1.096,0.390) (-1.086,0.423) (-1.075,0.455) (-1.064,0.488) (-1.053,0.516) (-1.041,0.546) (-1.027,0.579) (-1.014,0.607) (-1.001,0.635) (-0.987,0.663) (-0.970,0.696) (-0.956,0.723) (-0.938,0.754) (-0.923,0.780) (-0.908,0.806) (-0.892,0.832) (-0.875,0.858) (-0.859,0.884) (-0.842,0.910) (-0.825,0.936) (-0.808,0.962) (-0.790,0.988) (-0.773,1.013) (-0.756,1.040) (-0.739,1.066) (-0.722,1.092) (-0.705,1.118) (-0.689,1.144) (-0.674,1.170) (-0.659,1.196) (-0.643,1.224) (-0.627,1.255) (-0.611,1.287) (-0.598,1.316) (-0.585,1.346) (-0.572,1.378) (-0.560,1.411) (-0.549,1.443) (-0.539,1.474) (-0.530,1.508) (-0.521,1.541) (-0.514,1.572) (-0.506,1.606) (-0.499,1.638) (-0.493,1.671) (-0.488,1.703) (-0.483,1.736) (-0.478,1.768) (-0.473,1.800) (-0.469,1.833) (-0.465,1.866) (-0.462,1.898) (-0.459,1.930) (-0.456,1.963) (-0.453,1.996) (-0.450,2.028) (-0.448,2.061) (-0.445,2.093) (-0.443,2.126) (-0.441,2.158) (-0.439,2.191) (-0.438,2.223) (-0.436,2.256) (-0.434,2.288) (-0.433,2.321) (-0.432,2.353) (-0.430,2.385) (-0.429,2.417) (-0.428,2.451) (-0.427,2.483) (-0.426,2.516) (-0.425,2.548) (-0.424,2.581) (-0.423,2.600)}
\EQ{(-0.841,-2.600) (-0.843,-2.567) (-0.844,-2.535) (-0.846,-2.502) (-0.848,-2.470) (-0.849,-2.438) (-0.851,-2.405) (-0.853,-2.373) (-0.855,-2.340) (-0.857,-2.308) (-0.860,-2.275) (-0.862,-2.243) (-0.865,-2.212) (-0.867,-2.178) (-0.870,-2.145) (-0.873,-2.112) (-0.876,-2.080) (-0.879,-2.047) (-0.883,-2.015) (-0.887,-1.982) (-0.890,-1.950) (-0.894,-1.917) (-0.899,-1.885) (-0.903,-1.853) (-0.908,-1.820) (-0.913,-1.788) (-0.918,-1.755) (-0.924,-1.723) (-0.929,-1.690) (-0.936,-1.657) (-0.942,-1.625) (-0.949,-1.593) (-0.956,-1.563) (-0.964,-1.527) (-0.971,-1.495) (-0.980,-1.462) (-0.988,-1.431) (-0.997,-1.397) (-1.006,-1.365) (-1.016,-1.333) (-1.026,-1.300) (-1.036,-1.268) (-1.046,-1.236) (-1.058,-1.202) (-1.069,-1.170) (-1.079,-1.140) (-1.089,-1.111) (-1.101,-1.079) (-1.111,-1.049) (-1.124,-1.014) (-1.136,-0.982) (-1.148,-0.949) (-1.159,-0.916) (-1.170,-0.886) (-1.182,-0.851) (-1.194,-0.819) (-1.205,-0.786) (-1.215,-0.754) (-1.226,-0.722) (-1.235,-0.693) (-1.244,-0.663) (-1.254,-0.630) (-1.263,-0.598) (-1.272,-0.566) (-1.280,-0.533) (-1.288,-0.500) (-1.296,-0.468) (-1.303,-0.435) (-1.309,-0.403) (-1.315,-0.370) (-1.321,-0.338) (-1.326,-0.307) (-1.331,-0.273) (-1.335,-0.240) (-1.339,-0.208) (-1.342,-0.175) (-1.344,-0.143) (-1.346,-0.111) (-1.348,-0.078) (-1.349,-0.045) (-1.349,-0.013) (-1.349,0.020) (-1.349,0.052) (-1.348,0.085) (-1.346,0.117) (-1.344,0.150) (-1.341,0.182) (-1.338,0.215) (-1.334,0.247) (-1.330,0.280) (-1.325,0.312) (-1.320,0.345) (-1.314,0.377) (-1.308,0.409) (-1.301,0.442) (-1.294,0.475) (-1.287,0.505) (-1.278,0.540) (-1.270,0.572) (-1.261,0.605) (-1.252,0.637) (-1.242,0.670) (-1.232,0.702) (-1.222,0.734) (-1.211,0.767) (-1.200,0.800) (-1.189,0.831) (-1.178,0.865) (-1.166,0.897) (-1.155,0.930) (-1.144,0.959) (-1.133,0.988) (-1.122,1.021) (-1.110,1.053) (-1.099,1.085) (-1.087,1.118) (-1.075,1.151) (-1.064,1.183) (-1.053,1.216) (-1.042,1.248) (-1.032,1.281) (-1.022,1.313) (-1.012,1.346) (-1.003,1.378) (-0.993,1.411) (-0.985,1.443) (-0.976,1.476) (-0.968,1.507) (-0.961,1.541) (-0.953,1.573) (-0.946,1.606) (-0.940,1.638) (-0.933,1.671) (-0.927,1.703) (-0.921,1.736) (-0.916,1.765) (-0.911,1.800) (-0.906,1.833) (-0.901,1.866) (-0.897,1.898) (-0.893,1.930) (-0.889,1.963) (-0.885,1.996) (-0.882,2.028) (-0.878,2.061) (-0.875,2.093) (-0.872,2.126) (-0.869,2.158) (-0.866,2.191) (-0.864,2.223) (-0.861,2.256) (-0.859,2.288) (-0.857,2.321) (-0.854,2.353) (-0.852,2.385) (-0.851,2.418) (-0.849,2.450) (-0.847,2.483) (-0.845,2.516) (-0.844,2.548) (-0.842,2.581) (-0.841,2.600)}
\EQ{(-1.252,-2.600) (-1.254,-2.567) (-1.255,-2.535) (-1.257,-2.502) (-1.259,-2.470) (-1.261,-2.438) (-1.263,-2.405) (-1.265,-2.373) (-1.267,-2.340) (-1.269,-2.308) (-1.272,-2.275) (-1.274,-2.243) (-1.277,-2.210) (-1.279,-2.178) (-1.282,-2.145) (-1.285,-2.112) (-1.288,-2.080) (-1.291,-2.047) (-1.294,-2.015) (-1.298,-1.983) (-1.301,-1.950) (-1.305,-1.917) (-1.308,-1.885) (-1.312,-1.853) (-1.316,-1.820) (-1.320,-1.788) (-1.325,-1.755) (-1.329,-1.723) (-1.334,-1.690) (-1.338,-1.657) (-1.343,-1.625) (-1.348,-1.593) (-1.354,-1.560) (-1.359,-1.528) (-1.365,-1.495) (-1.370,-1.462) (-1.376,-1.430) (-1.382,-1.397) (-1.388,-1.365) (-1.394,-1.333) (-1.401,-1.300) (-1.407,-1.268) (-1.414,-1.235) (-1.420,-1.202) (-1.427,-1.170) (-1.434,-1.138) (-1.441,-1.105) (-1.448,-1.073) (-1.454,-1.040) (-1.461,-1.007) (-1.468,-0.975) (-1.475,-0.942) (-1.482,-0.910) (-1.488,-0.878) (-1.495,-0.846) (-1.502,-0.813) (-1.508,-0.781) (-1.514,-0.747) (-1.521,-0.715) (-1.527,-0.682) (-1.532,-0.650) (-1.538,-0.617) (-1.543,-0.585) (-1.548,-0.552) (-1.553,-0.520) (-1.558,-0.487) (-1.563,-0.455) (-1.567,-0.424) (-1.571,-0.390) (-1.574,-0.357) (-1.578,-0.325) (-1.581,-0.292) (-1.583,-0.260) (-1.586,-0.228) (-1.588,-0.195) (-1.590,-0.163) (-1.591,-0.130) (-1.592,-0.097) (-1.593,-0.065) (-1.594,-0.032) (-1.594,0.000) (-1.594,0.033) (-1.593,0.065) (-1.592,0.098) (-1.591,0.130) (-1.590,0.163) (-1.588,0.195) (-1.586,0.228) (-1.583,0.260) (-1.581,0.292) (-1.578,0.325) (-1.574,0.358) (-1.571,0.390) (-1.567,0.423) (-1.563,0.455) (-1.558,0.488) (-1.554,0.519) (-1.548,0.553) (-1.543,0.585) (-1.538,0.618) (-1.532,0.650) (-1.527,0.683) (-1.521,0.712) (-1.514,0.747) (-1.508,0.780) (-1.502,0.812) (-1.495,0.845) (-1.489,0.878) (-1.482,0.909) (-1.476,0.940) (-1.469,0.971) (-1.463,1.001) (-1.456,1.033) (-1.450,1.063) (-1.443,1.094) (-1.437,1.125) (-1.430,1.156) (-1.423,1.187) (-1.417,1.219) (-1.411,1.251) (-1.405,1.281) (-1.398,1.313) (-1.392,1.346) (-1.386,1.378) (-1.380,1.411) (-1.374,1.443) (-1.368,1.476) (-1.362,1.508) (-1.357,1.541) (-1.352,1.573) (-1.346,1.606) (-1.341,1.638) (-1.336,1.671) (-1.332,1.703) (-1.327,1.736) (-1.323,1.768) (-1.319,1.800) (-1.315,1.833) (-1.311,1.866) (-1.307,1.898) (-1.303,1.930) (-1.300,1.963) (-1.296,1.996) (-1.293,2.028) (-1.290,2.061) (-1.287,2.090) (-1.284,2.126) (-1.281,2.158) (-1.278,2.191) (-1.276,2.223) (-1.273,2.256) (-1.271,2.288) (-1.269,2.321) (-1.266,2.353) (-1.264,2.385) (-1.262,2.418) (-1.260,2.451) (-1.258,2.483) (-1.257,2.516) (-1.255,2.548) (-1.253,2.581) (-1.252,2.600)}
\EQ{(-1.657,-2.600) (-1.658,-2.567) (-1.660,-2.535) (-1.661,-2.502) (-1.663,-2.470) (-1.665,-2.438) (-1.667,-2.405) (-1.668,-2.373) (-1.670,-2.340) (-1.672,-2.308) (-1.674,-2.275) (-1.676,-2.243) (-1.679,-2.210) (-1.681,-2.178) (-1.683,-2.145) (-1.685,-2.112) (-1.688,-2.080) (-1.690,-2.047) (-1.693,-2.015) (-1.696,-1.982) (-1.698,-1.950) (-1.701,-1.917) (-1.704,-1.885) (-1.707,-1.853) (-1.710,-1.820) (-1.713,-1.788) (-1.716,-1.755) (-1.719,-1.723) (-1.723,-1.691) (-1.726,-1.657) (-1.729,-1.625) (-1.733,-1.592) (-1.737,-1.560) (-1.740,-1.528) (-1.744,-1.495) (-1.748,-1.462) (-1.752,-1.430) (-1.755,-1.397) (-1.759,-1.365) (-1.763,-1.333) (-1.767,-1.300) (-1.771,-1.268) (-1.776,-1.235) (-1.780,-1.202) (-1.784,-1.170) (-1.788,-1.137) (-1.792,-1.105) (-1.796,-1.073) (-1.800,-1.040) (-1.804,-1.007) (-1.809,-0.975) (-1.813,-0.942) (-1.817,-0.910) (-1.821,-0.877) (-1.824,-0.845) (-1.828,-0.812) (-1.832,-0.780) (-1.836,-0.747) (-1.839,-0.715) (-1.843,-0.682) (-1.846,-0.652) (-1.849,-0.617) (-1.853,-0.586) (-1.856,-0.552) (-1.859,-0.520) (-1.861,-0.487) (-1.864,-0.455) (-1.866,-0.422) (-1.869,-0.390) (-1.871,-0.357) (-1.873,-0.325) (-1.875,-0.292) (-1.876,-0.260) (-1.878,-0.228) (-1.879,-0.195) (-1.880,-0.163) (-1.881,-0.130) (-1.881,-0.097) (-1.882,-0.065) (-1.882,-0.032) (-1.882,0.000) (-1.882,0.033) (-1.882,0.065) (-1.881,0.098) (-1.881,0.130) (-1.880,0.163) (-1.879,0.195) (-1.878,0.228) (-1.876,0.260) (-1.875,0.292) (-1.873,0.325) (-1.871,0.358) (-1.869,0.390) (-1.866,0.423) (-1.864,0.455) (-1.861,0.488) (-1.859,0.520) (-1.856,0.553) (-1.853,0.585) (-1.849,0.618) (-1.846,0.650) (-1.843,0.683) (-1.840,0.713) (-1.836,0.748) (-1.832,0.780) (-1.828,0.812) (-1.824,0.845) (-1.821,0.878) (-1.817,0.910) (-1.813,0.943) (-1.809,0.975) (-1.804,1.008) (-1.801,1.039) (-1.796,1.073) (-1.792,1.105) (-1.788,1.138) (-1.784,1.170) (-1.780,1.203) (-1.776,1.235) (-1.771,1.268) (-1.767,1.300) (-1.763,1.333) (-1.759,1.365) (-1.755,1.398) (-1.752,1.430) (-1.748,1.462) (-1.744,1.495) (-1.740,1.528) (-1.737,1.560) (-1.733,1.593) (-1.729,1.625) (-1.726,1.658) (-1.723,1.690) (-1.719,1.723) (-1.716,1.755) (-1.713,1.788) (-1.710,1.820) (-1.707,1.853) (-1.704,1.885) (-1.701,1.917) (-1.698,1.950) (-1.696,1.983) (-1.693,2.015) (-1.690,2.047) (-1.688,2.080) (-1.685,2.113) (-1.683,2.145) (-1.681,2.178) (-1.679,2.210) (-1.676,2.243) (-1.674,2.275) (-1.672,2.308) (-1.670,2.340) (-1.668,2.373) (-1.667,2.405) (-1.665,2.438) (-1.663,2.470) (-1.661,2.502) (-1.660,2.535) (-1.658,2.568) (-1.657,2.600)}
\EQ{(-2.056,-2.600) (-2.058,-2.567) (-2.059,-2.535) (-2.060,-2.502) (-2.062,-2.470) (-2.063,-2.438) (-2.065,-2.405) (-2.066,-2.372) (-2.068,-2.340) (-2.069,-2.308) (-2.071,-2.275) (-2.073,-2.243) (-2.074,-2.210) (-2.076,-2.178) (-2.078,-2.145) (-2.080,-2.112) (-2.082,-2.080) (-2.084,-2.047) (-2.086,-2.015) (-2.087,-1.982) (-2.090,-1.950) (-2.092,-1.917) (-2.094,-1.885) (-2.096,-1.853) (-2.098,-1.820) (-2.100,-1.788) (-2.102,-1.755) (-2.105,-1.723) (-2.107,-1.690) (-2.109,-1.657) (-2.112,-1.625) (-2.114,-1.593) (-2.116,-1.560) (-2.119,-1.528) (-2.121,-1.495) (-2.124,-1.462) (-2.126,-1.430) (-2.129,-1.397) (-2.131,-1.365) (-2.134,-1.333) (-2.137,-1.300) (-2.139,-1.268) (-2.142,-1.235) (-2.144,-1.202) (-2.147,-1.170) (-2.149,-1.137) (-2.152,-1.105) (-2.155,-1.073) (-2.157,-1.040) (-2.160,-1.007) (-2.162,-0.975) (-2.164,-0.943) (-2.167,-0.910) (-2.169,-0.878) (-2.172,-0.845) (-2.174,-0.812) (-2.176,-0.780) (-2.178,-0.747) (-2.180,-0.715) (-2.183,-0.682) (-2.185,-0.650) (-2.186,-0.617) (-2.188,-0.585) (-2.190,-0.552) (-2.192,-0.520) (-2.193,-0.487) (-2.195,-0.455) (-2.196,-0.422) (-2.198,-0.390) (-2.199,-0.357) (-2.200,-0.325) (-2.201,-0.292) (-2.202,-0.260) (-2.203,-0.228) (-2.204,-0.197) (-2.204,-0.163) (-2.205,-0.130) (-2.205,-0.097) (-2.205,-0.065) (-2.206,-0.032) (-2.206,0.000) (-2.206,0.033) (-2.205,0.065) (-2.205,0.098) (-2.205,0.130) (-2.204,0.163) (-2.204,0.195) (-2.203,0.228) (-2.202,0.260) (-2.201,0.292) (-2.200,0.325) (-2.199,0.358) (-2.198,0.390) (-2.196,0.423) (-2.195,0.455) (-2.193,0.488) (-2.192,0.520) (-2.190,0.553) (-2.188,0.585) (-2.186,0.618) (-2.185,0.650) (-2.183,0.683) (-2.180,0.715) (-2.178,0.748) (-2.176,0.780) (-2.174,0.812) (-2.172,0.845) (-2.169,0.878) (-2.167,0.910) (-2.165,0.943) (-2.162,0.975) (-2.160,1.008) (-2.157,1.040) (-2.155,1.073) (-2.152,1.105) (-2.149,1.138) (-2.147,1.170) (-2.144,1.203) (-2.142,1.235) (-2.139,1.268) (-2.137,1.300) (-2.134,1.333) (-2.131,1.365) (-2.129,1.398) (-2.126,1.430) (-2.124,1.462) (-2.121,1.495) (-2.119,1.526) (-2.116,1.560) (-2.114,1.593) (-2.112,1.625) (-2.109,1.658) (-2.107,1.690) (-2.105,1.723) (-2.102,1.755) (-2.100,1.788) (-2.098,1.820) (-2.096,1.853) (-2.094,1.885) (-2.092,1.917) (-2.090,1.950) (-2.087,1.983) (-2.086,2.015) (-2.084,2.047) (-2.082,2.080) (-2.080,2.113) (-2.078,2.145) (-2.076,2.178) (-2.074,2.210) (-2.073,2.243) (-2.071,2.275) (-2.069,2.308) (-2.068,2.340) (-2.066,2.372) (-2.065,2.405) (-2.063,2.438) (-2.062,2.470) (-2.061,2.502) (-2.059,2.535) (-2.058,2.568) (-2.056,2.600)}
\EQ{(-2.454,-2.600) (-2.455,-2.567) (-2.456,-2.535) (-2.457,-2.502) (-2.458,-2.470) (-2.459,-2.438) (-2.460,-2.405) (-2.462,-2.373) (-2.463,-2.340) (-2.464,-2.308) (-2.465,-2.275) (-2.467,-2.243) (-2.468,-2.210) (-2.469,-2.178) (-2.471,-2.145) (-2.472,-2.112) (-2.473,-2.080) (-2.475,-2.047) (-2.476,-2.015) (-2.478,-1.982) (-2.479,-1.950) (-2.480,-1.917) (-2.482,-1.885) (-2.484,-1.853) (-2.485,-1.820) (-2.487,-1.788) (-2.488,-1.755) (-2.490,-1.723) (-2.491,-1.690) (-2.493,-1.657) (-2.495,-1.625) (-2.496,-1.593) (-2.498,-1.560) (-2.499,-1.527) (-2.501,-1.495) (-2.503,-1.462) (-2.504,-1.430) (-2.506,-1.397) (-2.508,-1.365) (-2.509,-1.333) (-2.511,-1.300) (-2.513,-1.268) (-2.514,-1.235) (-2.516,-1.202) (-2.518,-1.170) (-2.519,-1.137) (-2.521,-1.105) (-2.522,-1.073) (-2.524,-1.040) (-2.526,-1.007) (-2.527,-0.975) (-2.529,-0.942) (-2.530,-0.910) (-2.532,-0.877) (-2.533,-0.845) (-2.534,-0.812) (-2.536,-0.780) (-2.537,-0.747) (-2.538,-0.715) (-2.540,-0.682) (-2.541,-0.650) (-2.542,-0.617) (-2.543,-0.585) (-2.544,-0.552) (-2.545,-0.520) (-2.546,-0.487) (-2.547,-0.455) (-2.548,-0.422) (-2.549,-0.390) (-2.549,-0.357) (-2.550,-0.325) (-2.551,-0.292) (-2.551,-0.260) (-2.552,-0.228) (-2.552,-0.195) (-2.553,-0.163) (-2.553,-0.130) (-2.553,-0.097) (-2.553,-0.065) (-2.553,-0.032) (-2.553,0.000) (-2.553,0.033) (-2.553,0.065) (-2.553,0.098) (-2.553,0.130) (-2.553,0.163) (-2.552,0.195) (-2.552,0.228) (-2.551,0.260) (-2.551,0.292) (-2.550,0.325) (-2.549,0.358) (-2.549,0.390) (-2.548,0.423) (-2.547,0.455) (-2.546,0.488) (-2.545,0.520) (-2.544,0.553) (-2.543,0.585) (-2.542,0.618) (-2.541,0.650) (-2.540,0.683) (-2.538,0.715) (-2.537,0.748) (-2.536,0.780) (-2.534,0.812) (-2.533,0.845) (-2.532,0.878) (-2.530,0.910) (-2.529,0.943) (-2.527,0.975) (-2.526,1.008) (-2.524,1.040) (-2.522,1.073) (-2.521,1.105) (-2.519,1.138) (-2.518,1.170) (-2.516,1.203) (-2.514,1.235) (-2.513,1.268) (-2.511,1.300) (-2.509,1.333) (-2.508,1.365) (-2.506,1.398) (-2.504,1.430) (-2.503,1.462) (-2.501,1.495) (-2.499,1.528) (-2.498,1.560) (-2.496,1.593) (-2.495,1.625) (-2.493,1.658) (-2.491,1.690) (-2.490,1.723) (-2.488,1.755) (-2.487,1.788) (-2.485,1.820) (-2.484,1.853) (-2.482,1.885) (-2.480,1.917) (-2.479,1.950) (-2.478,1.983) (-2.476,2.015) (-2.475,2.047) (-2.473,2.080) (-2.472,2.113) (-2.471,2.145) (-2.469,2.178) (-2.468,2.210) (-2.467,2.243) (-2.465,2.275) (-2.464,2.308) (-2.463,2.340) (-2.462,2.373) (-2.460,2.405) (-2.459,2.438) (-2.458,2.470) (-2.457,2.501) (-2.456,2.535) (-2.455,2.568) (-2.454,2.600)}

      % Esfera en campo uniforme - líneas de campo
% Generado por scripts/generar_datos.py (no editar a mano).
\LC{(1.000,-0.105) (1.034,-0.108) (1.066,-0.111) (1.099,-0.114) (1.129,-0.117) (1.164,-0.120) (1.196,-0.123) (1.229,-0.125) (1.261,-0.128) (1.294,-0.130) (1.326,-0.132) (1.359,-0.135) (1.391,-0.137) (1.424,-0.139) (1.456,-0.141) (1.489,-0.142) (1.521,-0.144) (1.554,-0.146) (1.586,-0.147) (1.619,-0.149) (1.651,-0.150) (1.683,-0.151) (1.716,-0.153) (1.749,-0.154) (1.781,-0.155) (1.813,-0.156) (1.846,-0.157) (1.879,-0.158) (1.911,-0.159) (1.944,-0.160) (1.976,-0.161) (2.009,-0.161) (2.041,-0.162) (2.074,-0.163) (2.106,-0.163) (2.139,-0.164) (2.171,-0.165) (2.204,-0.165) (2.236,-0.166) (2.268,-0.166) (2.301,-0.167) (2.334,-0.167) (2.366,-0.168) (2.399,-0.168) (2.431,-0.169) (2.464,-0.169) (2.496,-0.169) (2.529,-0.170) (2.561,-0.170) (2.594,-0.171) (2.600,-0.171)}
\LC{(-2.600,-0.171) (-2.567,-0.170) (-2.535,-0.170) (-2.502,-0.170) (-2.470,-0.169) (-2.438,-0.169) (-2.405,-0.168) (-2.373,-0.168) (-2.340,-0.167) (-2.308,-0.167) (-2.275,-0.166) (-2.242,-0.166) (-2.210,-0.165) (-2.178,-0.165) (-2.145,-0.164) (-2.112,-0.164) (-2.080,-0.163) (-2.047,-0.162) (-2.015,-0.161) (-1.982,-0.161) (-1.950,-0.160) (-1.917,-0.159) (-1.885,-0.158) (-1.853,-0.157) (-1.820,-0.156) (-1.788,-0.155) (-1.755,-0.154) (-1.723,-0.153) (-1.690,-0.152) (-1.657,-0.150) (-1.625,-0.149) (-1.593,-0.147) (-1.560,-0.146) (-1.527,-0.144) (-1.495,-0.143) (-1.462,-0.141) (-1.430,-0.139) (-1.397,-0.137) (-1.365,-0.135) (-1.333,-0.133) (-1.300,-0.131) (-1.268,-0.128) (-1.235,-0.126) (-1.202,-0.123) (-1.170,-0.121) (-1.137,-0.118) (-1.105,-0.115) (-1.073,-0.112) (-1.040,-0.109) (-1.007,-0.105) (-1.000,-0.105)}
\LC{(1.000,0.105) (1.034,0.108) (1.066,0.111) (1.099,0.114) (1.129,0.117) (1.164,0.120) (1.196,0.123) (1.229,0.125) (1.261,0.128) (1.294,0.130) (1.326,0.132) (1.359,0.135) (1.391,0.137) (1.424,0.139) (1.456,0.141) (1.489,0.142) (1.521,0.144) (1.554,0.146) (1.586,0.147) (1.619,0.149) (1.651,0.150) (1.683,0.151) (1.716,0.153) (1.749,0.154) (1.781,0.155) (1.814,0.156) (1.846,0.157) (1.879,0.158) (1.911,0.159) (1.944,0.160) (1.976,0.161) (2.009,0.161) (2.041,0.162) (2.074,0.163) (2.106,0.163) (2.139,0.164) (2.171,0.165) (2.204,0.165) (2.236,0.166) (2.268,0.166) (2.301,0.167) (2.334,0.167) (2.366,0.168) (2.399,0.168) (2.431,0.169) (2.464,0.169) (2.496,0.169) (2.529,0.170) (2.561,0.170) (2.594,0.171) (2.600,0.171)}
\LC{(-2.600,0.171) (-2.567,0.170) (-2.535,0.170) (-2.502,0.170) (-2.470,0.169) (-2.438,0.169) (-2.405,0.168) (-2.373,0.168) (-2.340,0.167) (-2.308,0.167) (-2.275,0.166) (-2.242,0.166) (-2.210,0.165) (-2.178,0.165) (-2.145,0.164) (-2.112,0.164) (-2.080,0.163) (-2.047,0.162) (-2.015,0.161) (-1.982,0.161) (-1.950,0.160) (-1.917,0.159) (-1.885,0.158) (-1.853,0.157) (-1.820,0.156) (-1.788,0.155) (-1.755,0.154) (-1.723,0.153) (-1.690,0.152) (-1.657,0.150) (-1.625,0.149) (-1.593,0.147) (-1.560,0.146) (-1.528,0.144) (-1.495,0.143) (-1.462,0.141) (-1.430,0.139) (-1.397,0.137) (-1.365,0.135) (-1.333,0.133) (-1.300,0.131) (-1.268,0.128) (-1.235,0.126) (-1.202,0.123) (-1.170,0.121) (-1.137,0.118) (-1.105,0.115) (-1.073,0.112) (-1.040,0.109) (-1.007,0.105) (-1.000,0.105)}
\LC{(0.964,-0.278) (0.995,-0.287) (1.027,-0.295) (1.060,-0.304) (1.092,-0.312) (1.125,-0.320) (1.157,-0.327) (1.190,-0.334) (1.222,-0.341) (1.255,-0.348) (1.287,-0.354) (1.320,-0.360) (1.352,-0.365) (1.385,-0.370) (1.417,-0.375) (1.450,-0.380) (1.482,-0.385) (1.515,-0.389) (1.547,-0.393) (1.579,-0.397) (1.612,-0.400) (1.645,-0.404) (1.677,-0.407) (1.710,-0.410) (1.742,-0.413) (1.773,-0.416) (1.807,-0.419) (1.840,-0.421) (1.872,-0.424) (1.905,-0.426) (1.937,-0.428) (1.970,-0.430) (2.002,-0.432) (2.034,-0.434) (2.067,-0.436) (2.100,-0.438) (2.132,-0.439) (2.165,-0.441) (2.197,-0.442) (2.230,-0.444) (2.262,-0.445) (2.295,-0.446) (2.327,-0.448) (2.360,-0.449) (2.392,-0.450) (2.425,-0.451) (2.457,-0.452) (2.490,-0.453) (2.522,-0.454) (2.555,-0.455) (2.587,-0.456) (2.600,-0.456)}
\LC{(-2.600,-0.456) (-2.567,-0.455) (-2.535,-0.454) (-2.502,-0.453) (-2.470,-0.452) (-2.438,-0.451) (-2.405,-0.450) (-2.373,-0.449) (-2.340,-0.448) (-2.308,-0.447) (-2.275,-0.446) (-2.243,-0.444) (-2.210,-0.443) (-2.178,-0.441) (-2.145,-0.440) (-2.112,-0.438) (-2.080,-0.437) (-2.047,-0.435) (-2.015,-0.433) (-1.982,-0.431) (-1.950,-0.429) (-1.917,-0.427) (-1.885,-0.425) (-1.854,-0.422) (-1.820,-0.420) (-1.788,-0.417) (-1.755,-0.414) (-1.723,-0.412) (-1.690,-0.409) (-1.657,-0.405) (-1.625,-0.402) (-1.593,-0.398) (-1.560,-0.395) (-1.528,-0.391) (-1.495,-0.386) (-1.462,-0.382) (-1.430,-0.377) (-1.397,-0.372) (-1.365,-0.367) (-1.333,-0.362) (-1.300,-0.356) (-1.267,-0.350) (-1.235,-0.344) (-1.202,-0.337) (-1.170,-0.330) (-1.137,-0.323) (-1.105,-0.315) (-1.073,-0.307) (-1.040,-0.299) (-1.007,-0.290) (-0.975,-0.281) (-0.964,-0.278)}
\LC{(0.964,0.278) (0.995,0.287) (1.027,0.295) (1.060,0.304) (1.092,0.312) (1.125,0.320) (1.157,0.327) (1.190,0.334) (1.222,0.341) (1.255,0.348) (1.287,0.354) (1.320,0.360) (1.352,0.365) (1.385,0.370) (1.417,0.375) (1.450,0.380) (1.482,0.385) (1.515,0.389) (1.547,0.393) (1.579,0.397) (1.612,0.400) (1.645,0.404) (1.677,0.407) (1.710,0.410) (1.742,0.413) (1.773,0.416) (1.807,0.419) (1.840,0.421) (1.872,0.424) (1.905,0.426) (1.937,0.428) (1.970,0.430) (2.002,0.432) (2.034,0.434) (2.067,0.436) (2.100,0.438) (2.132,0.439) (2.165,0.441) (2.197,0.442) (2.230,0.444) (2.262,0.445) (2.295,0.446) (2.327,0.448) (2.360,0.449) (2.392,0.450) (2.425,0.451) (2.457,0.452) (2.490,0.453) (2.522,0.454) (2.555,0.455) (2.587,0.456) (2.600,0.456)}
\LC{(-2.600,0.456) (-2.567,0.455) (-2.535,0.454) (-2.502,0.453) (-2.470,0.452) (-2.438,0.451) (-2.405,0.450) (-2.373,0.449) (-2.340,0.448) (-2.308,0.447) (-2.275,0.446) (-2.243,0.444) (-2.210,0.443) (-2.178,0.441) (-2.145,0.440) (-2.112,0.438) (-2.080,0.437) (-2.047,0.435) (-2.015,0.433) (-1.982,0.431) (-1.950,0.429) (-1.917,0.427) (-1.885,0.425) (-1.854,0.423) (-1.820,0.420) (-1.788,0.417) (-1.755,0.414) (-1.723,0.412) (-1.690,0.409) (-1.657,0.405) (-1.625,0.402) (-1.593,0.398) (-1.560,0.395) (-1.528,0.391) (-1.495,0.386) (-1.462,0.382) (-1.430,0.377) (-1.397,0.372) (-1.365,0.367) (-1.333,0.362) (-1.300,0.356) (-1.267,0.350) (-1.235,0.344) (-1.202,0.337) (-1.170,0.330) (-1.137,0.323) (-1.105,0.315) (-1.073,0.307) (-1.040,0.299) (-1.007,0.290) (-0.975,0.281) (-0.964,0.278)}
\LC{(0.894,-0.451) (0.923,-0.465) (0.956,-0.481) (0.985,-0.494) (1.014,-0.506) (1.047,-0.520) (1.079,-0.532) (1.112,-0.544) (1.144,-0.555) (1.175,-0.566) (1.209,-0.576) (1.239,-0.585) (1.274,-0.595) (1.307,-0.603) (1.337,-0.611) (1.372,-0.619) (1.404,-0.626) (1.437,-0.633) (1.469,-0.640) (1.502,-0.646) (1.534,-0.652) (1.567,-0.657) (1.599,-0.663) (1.632,-0.668) (1.664,-0.672) (1.696,-0.677) (1.729,-0.681) (1.762,-0.685) (1.794,-0.689) (1.827,-0.693) (1.859,-0.696) (1.892,-0.699) (1.924,-0.703) (1.957,-0.706) (1.989,-0.708) (2.022,-0.711) (2.054,-0.714) (2.087,-0.716) (2.119,-0.718) (2.151,-0.721) (2.184,-0.723) (2.216,-0.725) (2.249,-0.727) (2.282,-0.729) (2.314,-0.731) (2.346,-0.732) (2.379,-0.734) (2.412,-0.735) (2.444,-0.737) (2.477,-0.738) (2.509,-0.740) (2.540,-0.741) (2.574,-0.742) (2.600,-0.743)}
\LC{(-2.600,-0.743) (-2.567,-0.742) (-2.535,-0.741) (-2.502,-0.739) (-2.470,-0.738) (-2.438,-0.737) (-2.405,-0.735) (-2.373,-0.734) (-2.340,-0.732) (-2.308,-0.730) (-2.275,-0.728) (-2.243,-0.726) (-2.210,-0.725) (-2.178,-0.722) (-2.145,-0.720) (-2.112,-0.718) (-2.080,-0.716) (-2.047,-0.713) (-2.015,-0.711) (-1.982,-0.708) (-1.950,-0.705) (-1.919,-0.702) (-1.885,-0.699) (-1.854,-0.696) (-1.820,-0.692) (-1.788,-0.688) (-1.755,-0.684) (-1.723,-0.680) (-1.691,-0.676) (-1.657,-0.671) (-1.625,-0.667) (-1.593,-0.662) (-1.561,-0.656) (-1.528,-0.651) (-1.495,-0.645) (-1.462,-0.639) (-1.430,-0.632) (-1.397,-0.625) (-1.365,-0.618) (-1.333,-0.610) (-1.300,-0.602) (-1.268,-0.593) (-1.235,-0.584) (-1.202,-0.574) (-1.170,-0.564) (-1.137,-0.553) (-1.105,-0.542) (-1.073,-0.530) (-1.040,-0.517) (-1.007,-0.504) (-0.975,-0.489) (-0.943,-0.474) (-0.915,-0.462) (-0.894,-0.451)}
\LC{(0.894,0.451) (0.923,0.465) (0.956,0.481) (0.985,0.494) (1.014,0.506) (1.047,0.520) (1.079,0.532) (1.112,0.544) (1.144,0.555) (1.175,0.566) (1.209,0.576) (1.239,0.585) (1.274,0.595) (1.307,0.603) (1.337,0.611) (1.372,0.619) (1.404,0.626) (1.437,0.633) (1.469,0.640) (1.502,0.646) (1.534,0.652) (1.567,0.657) (1.599,0.663) (1.632,0.668) (1.664,0.672) (1.696,0.677) (1.729,0.681) (1.762,0.685) (1.794,0.689) (1.827,0.693) (1.859,0.696) (1.892,0.699) (1.924,0.703) (1.957,0.706) (1.989,0.708) (2.022,0.711) (2.054,0.714) (2.087,0.716) (2.119,0.718) (2.151,0.721) (2.184,0.723) (2.217,0.725) (2.249,0.727) (2.282,0.729) (2.314,0.731) (2.347,0.732) (2.379,0.734) (2.412,0.735) (2.444,0.737) (2.476,0.738) (2.509,0.740) (2.540,0.741) (2.574,0.742) (2.600,0.743)}
\LC{(-2.600,0.743) (-2.567,0.742) (-2.535,0.741) (-2.503,0.739) (-2.470,0.738) (-2.438,0.737) (-2.405,0.735) (-2.373,0.734) (-2.340,0.732) (-2.308,0.730) (-2.275,0.728) (-2.243,0.726) (-2.210,0.725) (-2.178,0.722) (-2.145,0.720) (-2.112,0.718) (-2.080,0.716) (-2.047,0.713) (-2.015,0.711) (-1.983,0.708) (-1.950,0.705) (-1.919,0.702) (-1.885,0.699) (-1.854,0.696) (-1.820,0.692) (-1.788,0.688) (-1.755,0.684) (-1.723,0.680) (-1.691,0.676) (-1.658,0.671) (-1.625,0.667) (-1.592,0.662) (-1.561,0.657) (-1.528,0.651) (-1.495,0.645) (-1.462,0.639) (-1.430,0.632) (-1.397,0.625) (-1.365,0.618) (-1.333,0.610) (-1.300,0.602) (-1.268,0.593) (-1.235,0.584) (-1.202,0.574) (-1.170,0.564) (-1.137,0.553) (-1.105,0.542) (-1.073,0.530) (-1.040,0.517) (-1.007,0.504) (-0.975,0.489) (-0.943,0.475) (-0.915,0.462) (-0.894,0.451)}
\LC{(0.785,-0.626) (0.812,-0.647) (0.839,-0.666) (0.865,-0.684) (0.891,-0.701) (0.917,-0.717) (0.943,-0.732) (0.971,-0.747) (1.001,-0.763) (1.034,-0.779) (1.065,-0.793) (1.095,-0.806) (1.125,-0.818) (1.157,-0.830) (1.190,-0.842) (1.219,-0.851) (1.248,-0.861) (1.281,-0.871) (1.313,-0.880) (1.346,-0.889) (1.378,-0.897) (1.411,-0.905) (1.443,-0.912) (1.476,-0.919) (1.508,-0.926) (1.541,-0.932) (1.573,-0.938) (1.606,-0.944) (1.637,-0.949) (1.671,-0.954) (1.703,-0.959) (1.736,-0.964) (1.768,-0.968) (1.800,-0.972) (1.833,-0.976) (1.866,-0.980) (1.898,-0.984) (1.930,-0.987) (1.963,-0.991) (1.996,-0.994) (2.028,-0.997) (2.061,-1.000) (2.093,-1.003) (2.126,-1.005) (2.158,-1.008) (2.191,-1.010) (2.223,-1.013) (2.256,-1.015) (2.288,-1.017) (2.321,-1.019) (2.353,-1.021) (2.385,-1.023) (2.418,-1.025) (2.451,-1.027) (2.483,-1.028) (2.516,-1.030) (2.548,-1.031) (2.581,-1.033) (2.600,-1.034)}
\LC{(-2.600,-1.034) (-2.567,-1.032) (-2.535,-1.031) (-2.502,-1.029) (-2.470,-1.028) (-2.438,-1.026) (-2.405,-1.024) (-2.373,-1.022) (-2.342,-1.020) (-2.308,-1.018) (-2.275,-1.016) (-2.243,-1.014) (-2.210,-1.012) (-2.178,-1.009) (-2.145,-1.007) (-2.112,-1.004) (-2.080,-1.002) (-2.047,-0.999) (-2.015,-0.996) (-1.982,-0.993) (-1.950,-0.989) (-1.917,-0.986) (-1.885,-0.983) (-1.853,-0.979) (-1.821,-0.975) (-1.788,-0.971) (-1.755,-0.966) (-1.723,-0.962) (-1.690,-0.957) (-1.657,-0.952) (-1.625,-0.947) (-1.593,-0.941) (-1.562,-0.936) (-1.528,-0.930) (-1.495,-0.923) (-1.463,-0.916) (-1.434,-0.910) (-1.404,-0.903) (-1.371,-0.895) (-1.339,-0.887) (-1.306,-0.878) (-1.274,-0.869) (-1.242,-0.859) (-1.209,-0.848) (-1.180,-0.839) (-1.150,-0.828) (-1.118,-0.815) (-1.085,-0.802) (-1.053,-0.788) (-1.022,-0.773) (-0.994,-0.760) (-0.962,-0.743) (-0.935,-0.728) (-0.903,-0.709) (-0.877,-0.693) (-0.851,-0.675) (-0.825,-0.657) (-0.799,-0.637) (-0.785,-0.626)}
\LC{(0.785,0.626) (0.812,0.647) (0.839,0.666) (0.865,0.684) (0.891,0.701) (0.917,0.717) (0.943,0.732) (0.971,0.748) (1.001,0.763) (1.034,0.779) (1.065,0.793) (1.095,0.806) (1.125,0.818) (1.157,0.830) (1.190,0.842) (1.219,0.852) (1.248,0.861) (1.281,0.871) (1.313,0.880) (1.346,0.889) (1.378,0.897) (1.411,0.905) (1.443,0.912) (1.476,0.919) (1.508,0.926) (1.541,0.932) (1.573,0.938) (1.606,0.944) (1.637,0.949) (1.671,0.954) (1.703,0.959) (1.736,0.964) (1.768,0.968) (1.800,0.972) (1.833,0.976) (1.866,0.980) (1.898,0.984) (1.930,0.987) (1.963,0.991) (1.996,0.994) (2.028,0.997) (2.061,1.000) (2.093,1.003) (2.126,1.005) (2.158,1.008) (2.191,1.010) (2.223,1.013) (2.256,1.015) (2.288,1.017) (2.321,1.019) (2.353,1.021) (2.385,1.023) (2.418,1.025) (2.451,1.027) (2.483,1.028) (2.516,1.030) (2.548,1.031) (2.581,1.033) (2.600,1.034)}
\LC{(-2.600,1.034) (-2.568,1.032) (-2.535,1.031) (-2.502,1.029) (-2.470,1.028) (-2.438,1.026) (-2.405,1.024) (-2.373,1.022) (-2.342,1.021) (-2.308,1.018) (-2.275,1.016) (-2.243,1.014) (-2.210,1.012) (-2.178,1.009) (-2.145,1.007) (-2.112,1.004) (-2.080,1.002) (-2.047,0.999) (-2.015,0.996) (-1.982,0.993) (-1.950,0.989) (-1.917,0.986) (-1.885,0.983) (-1.853,0.979) (-1.821,0.975) (-1.788,0.971) (-1.755,0.966) (-1.723,0.962) (-1.690,0.957) (-1.657,0.952) (-1.625,0.947) (-1.593,0.941) (-1.562,0.936) (-1.528,0.930) (-1.495,0.923) (-1.463,0.917) (-1.434,0.910) (-1.404,0.903) (-1.371,0.895) (-1.339,0.887) (-1.306,0.878) (-1.274,0.869) (-1.242,0.859) (-1.209,0.848) (-1.180,0.839) (-1.150,0.828) (-1.118,0.815) (-1.085,0.802) (-1.053,0.788) (-1.022,0.774) (-0.994,0.760) (-0.962,0.743) (-0.935,0.728) (-0.903,0.709) (-0.877,0.693) (-0.852,0.675) (-0.825,0.657) (-0.799,0.637) (-0.785,0.626)}
\LC{(0.606,-0.798) (0.627,-0.825) (0.650,-0.851) (0.675,-0.877) (0.702,-0.903) (0.726,-0.923) (0.750,-0.942) (0.777,-0.962) (0.806,-0.982) (0.832,-0.998) (0.858,-1.013) (0.891,-1.031) (0.920,-1.046) (0.947,-1.059) (0.975,-1.072) (1.007,-1.085) (1.039,-1.099) (1.073,-1.111) (1.105,-1.122) (1.138,-1.133) (1.170,-1.143) (1.203,-1.153) (1.235,-1.162) (1.264,-1.170) (1.300,-1.179) (1.333,-1.187) (1.365,-1.194) (1.398,-1.201) (1.430,-1.208) (1.462,-1.214) (1.495,-1.221) (1.528,-1.226) (1.560,-1.232) (1.593,-1.237) (1.625,-1.242) (1.658,-1.247) (1.690,-1.252) (1.723,-1.256) (1.755,-1.261) (1.788,-1.265) (1.820,-1.269) (1.853,-1.272) (1.885,-1.276) (1.917,-1.279) (1.950,-1.283) (1.983,-1.286) (2.015,-1.289) (2.047,-1.292) (2.080,-1.295) (2.113,-1.297) (2.144,-1.300) (2.178,-1.303) (2.210,-1.305) (2.243,-1.307) (2.275,-1.309) (2.308,-1.312) (2.340,-1.314) (2.373,-1.316) (2.405,-1.318) (2.438,-1.319) (2.470,-1.321) (2.502,-1.323) (2.535,-1.325) (2.568,-1.326) (2.600,-1.328)}
\LC{(-2.600,-1.328) (-2.567,-1.326) (-2.535,-1.325) (-2.502,-1.323) (-2.470,-1.321) (-2.439,-1.319) (-2.405,-1.318) (-2.373,-1.316) (-2.340,-1.314) (-2.308,-1.312) (-2.275,-1.309) (-2.243,-1.307) (-2.210,-1.305) (-2.178,-1.303) (-2.145,-1.300) (-2.112,-1.297) (-2.080,-1.295) (-2.047,-1.292) (-2.015,-1.289) (-1.982,-1.286) (-1.950,-1.283) (-1.917,-1.279) (-1.885,-1.276) (-1.853,-1.272) (-1.820,-1.269) (-1.788,-1.265) (-1.757,-1.261) (-1.723,-1.256) (-1.690,-1.252) (-1.657,-1.247) (-1.625,-1.242) (-1.592,-1.237) (-1.560,-1.232) (-1.528,-1.226) (-1.495,-1.221) (-1.462,-1.214) (-1.430,-1.208) (-1.397,-1.201) (-1.365,-1.194) (-1.333,-1.187) (-1.300,-1.179) (-1.268,-1.171) (-1.235,-1.162) (-1.202,-1.153) (-1.172,-1.144) (-1.137,-1.133) (-1.105,-1.122) (-1.074,-1.111) (-1.040,-1.099) (-1.007,-1.086) (-0.976,-1.073) (-0.947,-1.059) (-0.920,-1.046) (-0.890,-1.031) (-0.859,-1.014) (-0.832,-0.998) (-0.806,-0.982) (-0.780,-0.964) (-0.754,-0.945) (-0.728,-0.925) (-0.703,-0.903) (-0.676,-0.878) (-0.650,-0.851) (-0.627,-0.825) (-0.606,-0.799) (-0.606,-0.798)}
\LC{(0.606,0.798) (0.627,0.826) (0.650,0.851) (0.675,0.878) (0.702,0.903) (0.726,0.923) (0.750,0.943) (0.777,0.962) (0.806,0.982) (0.832,0.998) (0.858,1.013) (0.891,1.031) (0.920,1.047) (0.947,1.060) (0.975,1.072) (1.007,1.086) (1.039,1.099) (1.073,1.111) (1.105,1.122) (1.138,1.133) (1.170,1.143) (1.203,1.153) (1.235,1.162) (1.264,1.170) (1.300,1.179) (1.333,1.187) (1.365,1.194) (1.398,1.201) (1.430,1.208) (1.462,1.214) (1.495,1.221) (1.528,1.226) (1.560,1.232) (1.593,1.237) (1.625,1.242) (1.657,1.247) (1.690,1.252) (1.723,1.256) (1.755,1.261) (1.788,1.265) (1.820,1.269) (1.853,1.272) (1.885,1.276) (1.917,1.279) (1.950,1.283) (1.983,1.286) (2.015,1.289) (2.047,1.292) (2.080,1.295) (2.113,1.297) (2.144,1.300) (2.178,1.303) (2.210,1.305) (2.243,1.307) (2.275,1.309) (2.308,1.312) (2.340,1.314) (2.373,1.316) (2.405,1.318) (2.438,1.319) (2.470,1.321) (2.502,1.323) (2.535,1.325) (2.568,1.326) (2.600,1.328)}
\LC{(-2.600,1.328) (-2.567,1.326) (-2.535,1.325) (-2.502,1.323) (-2.470,1.321) (-2.439,1.320) (-2.405,1.318) (-2.373,1.316) (-2.340,1.314) (-2.308,1.312) (-2.275,1.309) (-2.243,1.307) (-2.210,1.305) (-2.178,1.303) (-2.145,1.300) (-2.112,1.297) (-2.080,1.295) (-2.047,1.292) (-2.015,1.289) (-1.982,1.286) (-1.950,1.283) (-1.917,1.279) (-1.885,1.276) (-1.853,1.272) (-1.820,1.269) (-1.788,1.265) (-1.757,1.261) (-1.723,1.256) (-1.690,1.252) (-1.657,1.247) (-1.625,1.242) (-1.593,1.237) (-1.560,1.232) (-1.528,1.226) (-1.495,1.221) (-1.462,1.214) (-1.430,1.208) (-1.397,1.201) (-1.365,1.194) (-1.333,1.187) (-1.300,1.179) (-1.268,1.171) (-1.235,1.162) (-1.202,1.153) (-1.172,1.144) (-1.137,1.133) (-1.105,1.122) (-1.074,1.112) (-1.040,1.099) (-1.007,1.086) (-0.976,1.073) (-0.947,1.060) (-0.920,1.047) (-0.890,1.031) (-0.859,1.014) (-0.832,0.998) (-0.806,0.982) (-0.780,0.964) (-0.754,0.945) (-0.728,0.925) (-0.703,0.904) (-0.676,0.878) (-0.650,0.852) (-0.627,0.826) (-0.606,0.799) (-0.606,0.798)}
\LC{(0.245,-0.975) (0.256,-1.007) (0.272,-1.040) (0.290,-1.066) (0.310,-1.092) (0.332,-1.115) (0.356,-1.137) (0.379,-1.157) (0.403,-1.176) (0.429,-1.194) (0.455,-1.212) (0.481,-1.228) (0.507,-1.243) (0.540,-1.261) (0.572,-1.278) (0.605,-1.293) (0.633,-1.307) (0.663,-1.319) (0.694,-1.333) (0.727,-1.345) (0.761,-1.358) (0.793,-1.370) (0.826,-1.381) (0.857,-1.391) (0.891,-1.402) (0.920,-1.411) (0.956,-1.421) (0.988,-1.430) (1.021,-1.438) (1.053,-1.446) (1.086,-1.454) (1.118,-1.462) (1.150,-1.469) (1.181,-1.476) (1.212,-1.482) (1.242,-1.488) (1.274,-1.494) (1.306,-1.500) (1.339,-1.506) (1.372,-1.512) (1.404,-1.517) (1.437,-1.522) (1.469,-1.527) (1.502,-1.532) (1.534,-1.537) (1.567,-1.541) (1.599,-1.545) (1.632,-1.549) (1.664,-1.553) (1.696,-1.557) (1.729,-1.561) (1.762,-1.565) (1.794,-1.568) (1.827,-1.571) (1.859,-1.575) (1.892,-1.578) (1.924,-1.581) (1.957,-1.584) (1.989,-1.586) (2.022,-1.589) (2.054,-1.592) (2.087,-1.594) (2.119,-1.597) (2.151,-1.599) (2.184,-1.601) (2.217,-1.604) (2.247,-1.606) (2.282,-1.608) (2.314,-1.610) (2.347,-1.612) (2.379,-1.613) (2.412,-1.615) (2.444,-1.617) (2.477,-1.619) (2.509,-1.620) (2.542,-1.622) (2.574,-1.623) (2.600,-1.625)}
\LC{(-2.600,-1.625) (-2.567,-1.623) (-2.535,-1.622) (-2.502,-1.620) (-2.470,-1.618) (-2.438,-1.617) (-2.405,-1.615) (-2.373,-1.613) (-2.340,-1.611) (-2.308,-1.609) (-2.275,-1.607) (-2.243,-1.605) (-2.210,-1.603) (-2.178,-1.601) (-2.145,-1.599) (-2.112,-1.596) (-2.080,-1.594) (-2.047,-1.591) (-2.015,-1.589) (-1.984,-1.586) (-1.950,-1.583) (-1.917,-1.580) (-1.885,-1.577) (-1.853,-1.574) (-1.820,-1.571) (-1.788,-1.567) (-1.755,-1.564) (-1.723,-1.560) (-1.690,-1.557) (-1.657,-1.553) (-1.625,-1.549) (-1.592,-1.544) (-1.563,-1.540) (-1.528,-1.536) (-1.495,-1.531) (-1.462,-1.526) (-1.430,-1.521) (-1.397,-1.516) (-1.365,-1.511) (-1.333,-1.505) (-1.300,-1.499) (-1.268,-1.493) (-1.235,-1.487) (-1.202,-1.480) (-1.170,-1.473) (-1.137,-1.466) (-1.105,-1.459) (-1.073,-1.451) (-1.040,-1.443) (-1.007,-1.435) (-0.975,-1.426) (-0.943,-1.417) (-0.910,-1.407) (-0.877,-1.398) (-0.845,-1.387) (-0.812,-1.376) (-0.780,-1.365) (-0.747,-1.353) (-0.715,-1.341) (-0.682,-1.328) (-0.650,-1.314) (-0.619,-1.300) (-0.591,-1.287) (-0.559,-1.271) (-0.528,-1.254) (-0.500,-1.239) (-0.474,-1.224) (-0.449,-1.207) (-0.422,-1.190) (-0.397,-1.171) (-0.371,-1.150) (-0.345,-1.127) (-0.322,-1.105) (-0.300,-1.079) (-0.281,-1.053) (-0.265,-1.027) (-0.251,-0.994) (-0.245,-0.975)}
\LC{(0.245,0.975) (0.256,1.008) (0.272,1.040) (0.290,1.066) (0.310,1.092) (0.332,1.115) (0.356,1.138) (0.379,1.157) (0.403,1.176) (0.429,1.194) (0.455,1.212) (0.481,1.228) (0.507,1.243) (0.540,1.261) (0.572,1.278) (0.605,1.293) (0.633,1.307) (0.663,1.319) (0.694,1.333) (0.727,1.346) (0.761,1.358) (0.793,1.370) (0.826,1.381) (0.857,1.391) (0.891,1.402) (0.920,1.411) (0.956,1.421) (0.988,1.430) (1.021,1.438) (1.053,1.446) (1.086,1.454) (1.118,1.462) (1.150,1.469) (1.181,1.476) (1.212,1.482) (1.242,1.488) (1.274,1.494) (1.307,1.500) (1.339,1.506) (1.372,1.512) (1.404,1.517) (1.437,1.522) (1.469,1.527) (1.502,1.532) (1.534,1.537) (1.567,1.541) (1.599,1.545) (1.632,1.549) (1.664,1.553) (1.696,1.557) (1.729,1.561) (1.762,1.565) (1.794,1.568) (1.827,1.571) (1.859,1.575) (1.892,1.578) (1.924,1.581) (1.957,1.584) (1.989,1.586) (2.022,1.589) (2.054,1.592) (2.087,1.594) (2.119,1.597) (2.151,1.599) (2.184,1.601) (2.216,1.604) (2.247,1.606) (2.282,1.608) (2.314,1.610) (2.347,1.612) (2.379,1.613) (2.412,1.615) (2.444,1.617) (2.477,1.619) (2.509,1.620) (2.542,1.622) (2.574,1.623) (2.600,1.625)}
\LC{(-2.600,1.625) (-2.567,1.623) (-2.535,1.622) (-2.502,1.620) (-2.470,1.618) (-2.438,1.617) (-2.405,1.615) (-2.373,1.613) (-2.340,1.611) (-2.308,1.609) (-2.275,1.607) (-2.243,1.605) (-2.210,1.603) (-2.178,1.601) (-2.145,1.599) (-2.112,1.596) (-2.080,1.594) (-2.047,1.591) (-2.015,1.589) (-1.984,1.586) (-1.950,1.583) (-1.917,1.580) (-1.885,1.577) (-1.853,1.574) (-1.820,1.571) (-1.788,1.567) (-1.755,1.564) (-1.723,1.560) (-1.690,1.557) (-1.657,1.553) (-1.625,1.549) (-1.593,1.544) (-1.563,1.541) (-1.528,1.536) (-1.495,1.531) (-1.462,1.526) (-1.430,1.521) (-1.397,1.516) (-1.365,1.511) (-1.333,1.505) (-1.300,1.499) (-1.268,1.493) (-1.235,1.487) (-1.202,1.480) (-1.170,1.473) (-1.137,1.466) (-1.105,1.459) (-1.073,1.451) (-1.040,1.443) (-1.007,1.435) (-0.975,1.426) (-0.943,1.417) (-0.910,1.407) (-0.877,1.398) (-0.845,1.387) (-0.812,1.376) (-0.780,1.365) (-0.747,1.353) (-0.715,1.341) (-0.682,1.328) (-0.650,1.314) (-0.619,1.300) (-0.591,1.287) (-0.559,1.271) (-0.528,1.255) (-0.500,1.239) (-0.474,1.224) (-0.449,1.207) (-0.422,1.190) (-0.397,1.171) (-0.371,1.151) (-0.345,1.127) (-0.322,1.105) (-0.300,1.079) (-0.281,1.053) (-0.265,1.027) (-0.251,0.995) (-0.245,0.975)}
\LC{(-2.600,-1.924) (-2.567,-1.923) (-2.535,-1.921) (-2.502,-1.920) (-2.470,-1.918) (-2.438,-1.917) (-2.405,-1.915) (-2.373,-1.914) (-2.340,-1.912) (-2.308,-1.910) (-2.275,-1.908) (-2.243,-1.907) (-2.210,-1.905) (-2.178,-1.903) (-2.145,-1.901) (-2.112,-1.899) (-2.080,-1.897) (-2.047,-1.895) (-2.015,-1.892) (-1.982,-1.890) (-1.950,-1.888) (-1.917,-1.885) (-1.885,-1.883) (-1.853,-1.880) (-1.820,-1.878) (-1.788,-1.875) (-1.755,-1.872) (-1.723,-1.869) (-1.690,-1.867) (-1.658,-1.864) (-1.625,-1.860) (-1.593,-1.857) (-1.560,-1.854) (-1.528,-1.851) (-1.495,-1.847) (-1.462,-1.844) (-1.430,-1.840) (-1.397,-1.836) (-1.365,-1.832) (-1.333,-1.828) (-1.300,-1.824) (-1.268,-1.820) (-1.235,-1.816) (-1.202,-1.812) (-1.170,-1.807) (-1.137,-1.803) (-1.105,-1.798) (-1.073,-1.793) (-1.040,-1.788) (-1.007,-1.783) (-0.975,-1.778) (-0.942,-1.773) (-0.911,-1.768) (-0.877,-1.762) (-0.845,-1.757) (-0.812,-1.752) (-0.780,-1.746) (-0.747,-1.740) (-0.715,-1.735) (-0.682,-1.729) (-0.650,-1.723) (-0.617,-1.718) (-0.585,-1.712) (-0.552,-1.706) (-0.520,-1.701) (-0.487,-1.695) (-0.458,-1.690) (-0.422,-1.684) (-0.390,-1.679) (-0.357,-1.674) (-0.325,-1.670) (-0.292,-1.665) (-0.260,-1.661) (-0.228,-1.657) (-0.195,-1.654) (-0.163,-1.651) (-0.130,-1.649) (-0.097,-1.647) (-0.065,-1.646) (-0.032,-1.645) (0.000,-1.644) (0.033,-1.645) (0.065,-1.646) (0.098,-1.647) (0.130,-1.649) (0.163,-1.651) (0.195,-1.654) (0.228,-1.657) (0.260,-1.661) (0.292,-1.665) (0.325,-1.670) (0.358,-1.674) (0.390,-1.679) (0.423,-1.684) (0.455,-1.690) (0.488,-1.695) (0.520,-1.701) (0.553,-1.706) (0.585,-1.712) (0.618,-1.718) (0.650,-1.723) (0.682,-1.729) (0.715,-1.735) (0.748,-1.740) (0.780,-1.746) (0.812,-1.752) (0.845,-1.757) (0.878,-1.762) (0.910,-1.768) (0.943,-1.773) (0.975,-1.778) (1.008,-1.783) (1.040,-1.788) (1.073,-1.793) (1.105,-1.798) (1.138,-1.803) (1.169,-1.807) (1.203,-1.812) (1.235,-1.816) (1.266,-1.820) (1.300,-1.824) (1.333,-1.828) (1.365,-1.832) (1.398,-1.836) (1.430,-1.840) (1.462,-1.844) (1.495,-1.847) (1.528,-1.851) (1.560,-1.854) (1.593,-1.857) (1.625,-1.860) (1.658,-1.864) (1.690,-1.867) (1.723,-1.869) (1.755,-1.872) (1.788,-1.875) (1.820,-1.878) (1.853,-1.880) (1.885,-1.883) (1.917,-1.885) (1.950,-1.888) (1.983,-1.890) (2.015,-1.892) (2.047,-1.895) (2.080,-1.897) (2.113,-1.899) (2.145,-1.901) (2.178,-1.903) (2.210,-1.905) (2.243,-1.907) (2.275,-1.908) (2.308,-1.910) (2.340,-1.912) (2.373,-1.914) (2.405,-1.915) (2.438,-1.917) (2.470,-1.918) (2.502,-1.920) (2.535,-1.921) (2.568,-1.923) (2.600,-1.924)}
\LC{(-2.600,1.924) (-2.567,1.923) (-2.535,1.921) (-2.502,1.920) (-2.470,1.918) (-2.438,1.917) (-2.405,1.915) (-2.373,1.914) (-2.340,1.912) (-2.308,1.910) (-2.275,1.908) (-2.243,1.907) (-2.210,1.905) (-2.178,1.903) (-2.145,1.901) (-2.112,1.899) (-2.080,1.897) (-2.047,1.895) (-2.015,1.892) (-1.982,1.890) (-1.950,1.888) (-1.917,1.885) (-1.885,1.883) (-1.853,1.880) (-1.820,1.878) (-1.788,1.875) (-1.755,1.872) (-1.723,1.869) (-1.690,1.867) (-1.657,1.864) (-1.625,1.860) (-1.593,1.857) (-1.560,1.854) (-1.528,1.851) (-1.495,1.847) (-1.462,1.844) (-1.430,1.840) (-1.397,1.836) (-1.365,1.832) (-1.333,1.828) (-1.300,1.824) (-1.268,1.820) (-1.235,1.816) (-1.202,1.812) (-1.170,1.807) (-1.137,1.803) (-1.105,1.798) (-1.073,1.793) (-1.040,1.788) (-1.007,1.783) (-0.975,1.778) (-0.942,1.773) (-0.911,1.768) (-0.877,1.762) (-0.845,1.757) (-0.812,1.752) (-0.780,1.746) (-0.747,1.740) (-0.715,1.735) (-0.682,1.729) (-0.650,1.723) (-0.617,1.718) (-0.585,1.712) (-0.552,1.706) (-0.520,1.701) (-0.487,1.695) (-0.458,1.690) (-0.422,1.684) (-0.390,1.679) (-0.357,1.674) (-0.325,1.670) (-0.292,1.665) (-0.260,1.661) (-0.228,1.658) (-0.195,1.654) (-0.163,1.651) (-0.130,1.649) (-0.097,1.647) (-0.065,1.646) (-0.032,1.645) (0.000,1.644) (0.033,1.645) (0.065,1.646) (0.098,1.647) (0.130,1.649) (0.163,1.651) (0.195,1.654) (0.228,1.657) (0.260,1.661) (0.292,1.665) (0.325,1.670) (0.358,1.674) (0.390,1.679) (0.423,1.684) (0.455,1.690) (0.488,1.695) (0.520,1.701) (0.553,1.706) (0.585,1.712) (0.618,1.718) (0.650,1.723) (0.682,1.729) (0.715,1.735) (0.748,1.740) (0.780,1.746) (0.812,1.752) (0.845,1.757) (0.878,1.762) (0.910,1.768) (0.943,1.773) (0.975,1.778) (1.008,1.783) (1.040,1.788) (1.073,1.793) (1.105,1.798) (1.138,1.803) (1.169,1.807) (1.203,1.812) (1.235,1.816) (1.266,1.820) (1.300,1.824) (1.333,1.828) (1.365,1.832) (1.398,1.836) (1.430,1.840) (1.462,1.844) (1.495,1.847) (1.528,1.851) (1.560,1.854) (1.593,1.857) (1.625,1.860) (1.658,1.864) (1.690,1.867) (1.723,1.869) (1.755,1.872) (1.788,1.875) (1.820,1.878) (1.853,1.880) (1.885,1.883) (1.917,1.885) (1.950,1.888) (1.983,1.890) (2.015,1.892) (2.047,1.895) (2.080,1.897) (2.113,1.899) (2.145,1.901) (2.178,1.903) (2.210,1.905) (2.243,1.907) (2.275,1.908) (2.308,1.910) (2.340,1.912) (2.373,1.914) (2.405,1.915) (2.438,1.917) (2.470,1.918) (2.502,1.920) (2.535,1.921) (2.568,1.923) (2.600,1.924)}
\LC{(-2.600,-2.225) (-2.567,-2.224) (-2.535,-2.223) (-2.502,-2.221) (-2.470,-2.220) (-2.438,-2.219) (-2.405,-2.218) (-2.373,-2.216) (-2.340,-2.215) (-2.308,-2.213) (-2.275,-2.212) (-2.243,-2.210) (-2.210,-2.209) (-2.178,-2.207) (-2.145,-2.206) (-2.112,-2.204) (-2.080,-2.202) (-2.048,-2.200) (-2.015,-2.199) (-1.982,-2.197) (-1.950,-2.195) (-1.917,-2.193) (-1.885,-2.191) (-1.853,-2.189) (-1.820,-2.187) (-1.788,-2.185) (-1.755,-2.183) (-1.723,-2.181) (-1.690,-2.178) (-1.657,-2.176) (-1.625,-2.174) (-1.593,-2.172) (-1.560,-2.169) (-1.528,-2.167) (-1.495,-2.164) (-1.462,-2.162) (-1.430,-2.159) (-1.397,-2.156) (-1.365,-2.154) (-1.333,-2.151) (-1.300,-2.148) (-1.268,-2.146) (-1.235,-2.143) (-1.202,-2.140) (-1.170,-2.137) (-1.137,-2.134) (-1.105,-2.131) (-1.073,-2.128) (-1.040,-2.125) (-1.007,-2.122) (-0.975,-2.119) (-0.942,-2.116) (-0.910,-2.113) (-0.877,-2.110) (-0.845,-2.107) (-0.812,-2.104) (-0.780,-2.101) (-0.747,-2.098) (-0.715,-2.095) (-0.682,-2.092) (-0.650,-2.089) (-0.618,-2.087) (-0.585,-2.084) (-0.552,-2.081) (-0.520,-2.078) (-0.487,-2.076) (-0.455,-2.074) (-0.422,-2.071) (-0.390,-2.069) (-0.357,-2.067) (-0.325,-2.065) (-0.292,-2.063) (-0.260,-2.062) (-0.228,-2.060) (-0.195,-2.059) (-0.163,-2.058) (-0.130,-2.057) (-0.097,-2.056) (-0.065,-2.056) (-0.032,-2.056) (0.000,-2.056) (0.033,-2.056) (0.065,-2.056) (0.098,-2.056) (0.130,-2.057) (0.163,-2.058) (0.195,-2.059) (0.228,-2.060) (0.260,-2.062) (0.292,-2.063) (0.325,-2.065) (0.357,-2.067) (0.390,-2.069) (0.423,-2.071) (0.454,-2.074) (0.488,-2.076) (0.520,-2.078) (0.553,-2.081) (0.585,-2.084) (0.618,-2.086) (0.650,-2.089) (0.683,-2.092) (0.715,-2.095) (0.748,-2.098) (0.780,-2.101) (0.812,-2.104) (0.845,-2.107) (0.878,-2.110) (0.910,-2.113) (0.943,-2.116) (0.973,-2.119) (1.008,-2.122) (1.040,-2.125) (1.073,-2.128) (1.105,-2.131) (1.138,-2.134) (1.170,-2.137) (1.203,-2.140) (1.235,-2.143) (1.268,-2.146) (1.300,-2.148) (1.333,-2.151) (1.365,-2.154) (1.398,-2.156) (1.430,-2.159) (1.462,-2.162) (1.495,-2.164) (1.528,-2.167) (1.560,-2.169) (1.593,-2.172) (1.625,-2.174) (1.658,-2.176) (1.690,-2.178) (1.723,-2.181) (1.755,-2.183) (1.788,-2.185) (1.820,-2.187) (1.853,-2.189) (1.885,-2.191) (1.917,-2.193) (1.950,-2.195) (1.983,-2.197) (2.015,-2.199) (2.048,-2.200) (2.080,-2.202) (2.113,-2.204) (2.145,-2.206) (2.178,-2.207) (2.210,-2.209) (2.243,-2.210) (2.275,-2.212) (2.308,-2.213) (2.340,-2.215) (2.373,-2.216) (2.405,-2.218) (2.438,-2.219) (2.470,-2.220) (2.502,-2.221) (2.535,-2.223) (2.568,-2.224) (2.600,-2.225)}
\LC{(-2.600,2.225) (-2.567,2.224) (-2.535,2.223) (-2.502,2.221) (-2.470,2.220) (-2.438,2.219) (-2.405,2.218) (-2.372,2.216) (-2.340,2.215) (-2.308,2.213) (-2.275,2.212) (-2.243,2.210) (-2.210,2.209) (-2.178,2.207) (-2.145,2.206) (-2.112,2.204) (-2.080,2.202) (-2.047,2.200) (-2.015,2.199) (-1.982,2.197) (-1.950,2.195) (-1.917,2.193) (-1.885,2.191) (-1.853,2.189) (-1.820,2.187) (-1.788,2.185) (-1.755,2.183) (-1.723,2.181) (-1.690,2.178) (-1.658,2.176) (-1.625,2.174) (-1.593,2.172) (-1.560,2.169) (-1.527,2.167) (-1.495,2.164) (-1.462,2.162) (-1.430,2.159) (-1.397,2.156) (-1.365,2.154) (-1.333,2.151) (-1.300,2.148) (-1.268,2.146) (-1.235,2.143) (-1.202,2.140) (-1.170,2.137) (-1.137,2.134) (-1.105,2.131) (-1.073,2.128) (-1.040,2.125) (-1.007,2.122) (-0.975,2.119) (-0.942,2.116) (-0.910,2.113) (-0.877,2.110) (-0.845,2.107) (-0.812,2.104) (-0.780,2.101) (-0.747,2.098) (-0.715,2.095) (-0.682,2.092) (-0.650,2.089) (-0.618,2.087) (-0.585,2.084) (-0.552,2.081) (-0.520,2.078) (-0.487,2.076) (-0.455,2.074) (-0.422,2.071) (-0.390,2.069) (-0.357,2.067) (-0.325,2.065) (-0.292,2.063) (-0.260,2.062) (-0.228,2.060) (-0.195,2.059) (-0.163,2.058) (-0.130,2.057) (-0.097,2.056) (-0.065,2.056) (-0.032,2.056) (0.000,2.056) (0.033,2.056) (0.065,2.056) (0.098,2.056) (0.130,2.057) (0.163,2.058) (0.195,2.059) (0.228,2.060) (0.260,2.062) (0.292,2.063) (0.325,2.065) (0.357,2.067) (0.390,2.069) (0.423,2.071) (0.454,2.074) (0.488,2.076) (0.520,2.078) (0.553,2.081) (0.585,2.084) (0.618,2.086) (0.650,2.089) (0.683,2.092) (0.715,2.095) (0.748,2.098) (0.780,2.101) (0.812,2.104) (0.845,2.107) (0.878,2.110) (0.910,2.113) (0.943,2.116) (0.973,2.119) (1.008,2.122) (1.040,2.125) (1.073,2.128) (1.105,2.131) (1.138,2.134) (1.170,2.137) (1.203,2.140) (1.235,2.143) (1.268,2.146) (1.300,2.148) (1.333,2.151) (1.365,2.154) (1.398,2.156) (1.430,2.159) (1.462,2.162) (1.495,2.164) (1.528,2.167) (1.560,2.169) (1.593,2.172) (1.625,2.174) (1.658,2.176) (1.690,2.178) (1.723,2.181) (1.755,2.183) (1.788,2.185) (1.820,2.187) (1.853,2.189) (1.885,2.191) (1.917,2.193) (1.950,2.195) (1.983,2.197) (2.015,2.199) (2.047,2.200) (2.080,2.202) (2.113,2.204) (2.145,2.206) (2.178,2.207) (2.210,2.209) (2.243,2.210) (2.275,2.212) (2.308,2.213) (2.340,2.215) (2.372,2.216) (2.405,2.218) (2.438,2.219) (2.470,2.220) (2.502,2.221) (2.535,2.223) (2.568,2.224) (2.600,2.225)}
\LC{(-2.600,-2.528) (-2.567,-2.527) (-2.535,-2.525) (-2.502,-2.524) (-2.470,-2.523) (-2.438,-2.522) (-2.405,-2.521) (-2.373,-2.520) (-2.340,-2.519) (-2.308,-2.518) (-2.275,-2.516) (-2.243,-2.515) (-2.210,-2.514) (-2.178,-2.513) (-2.145,-2.511) (-2.112,-2.510) (-2.080,-2.509) (-2.047,-2.507) (-2.015,-2.506) (-1.982,-2.504) (-1.950,-2.503) (-1.918,-2.501) (-1.885,-2.500) (-1.853,-2.498) (-1.820,-2.497) (-1.788,-2.495) (-1.755,-2.494) (-1.723,-2.492) (-1.690,-2.490) (-1.657,-2.489) (-1.625,-2.487) (-1.592,-2.485) (-1.560,-2.483) (-1.528,-2.482) (-1.495,-2.480) (-1.462,-2.478) (-1.430,-2.476) (-1.397,-2.474) (-1.365,-2.473) (-1.333,-2.471) (-1.300,-2.469) (-1.268,-2.467) (-1.235,-2.465) (-1.203,-2.463) (-1.170,-2.461) (-1.137,-2.459) (-1.105,-2.457) (-1.073,-2.456) (-1.040,-2.454) (-1.007,-2.452) (-0.975,-2.450) (-0.942,-2.448) (-0.910,-2.446) (-0.877,-2.444) (-0.845,-2.442) (-0.812,-2.441) (-0.780,-2.439) (-0.747,-2.437) (-0.715,-2.435) (-0.682,-2.434) (-0.650,-2.432) (-0.617,-2.430) (-0.585,-2.429) (-0.552,-2.427) (-0.520,-2.426) (-0.487,-2.425) (-0.455,-2.423) (-0.422,-2.422) (-0.390,-2.421) (-0.357,-2.420) (-0.325,-2.419) (-0.292,-2.418) (-0.260,-2.417) (-0.228,-2.417) (-0.195,-2.416) (-0.163,-2.415) (-0.130,-2.415) (-0.097,-2.415) (-0.065,-2.414) (-0.032,-2.414) (0.000,-2.414) (0.033,-2.414) (0.065,-2.414) (0.098,-2.415) (0.130,-2.415) (0.163,-2.415) (0.195,-2.416) (0.228,-2.417) (0.260,-2.417) (0.292,-2.418) (0.325,-2.419) (0.358,-2.420) (0.390,-2.421) (0.423,-2.422) (0.455,-2.423) (0.488,-2.425) (0.520,-2.426) (0.553,-2.427) (0.585,-2.429) (0.618,-2.430) (0.650,-2.432) (0.683,-2.434) (0.715,-2.435) (0.748,-2.437) (0.780,-2.439) (0.812,-2.441) (0.845,-2.442) (0.878,-2.444) (0.910,-2.446) (0.943,-2.448) (0.975,-2.450) (1.008,-2.452) (1.040,-2.454) (1.073,-2.456) (1.105,-2.457) (1.138,-2.459) (1.170,-2.461) (1.203,-2.463) (1.235,-2.465) (1.268,-2.467) (1.300,-2.469) (1.333,-2.471) (1.365,-2.473) (1.398,-2.474) (1.430,-2.476) (1.462,-2.478) (1.495,-2.480) (1.528,-2.482) (1.560,-2.483) (1.593,-2.485) (1.625,-2.487) (1.657,-2.489) (1.690,-2.490) (1.723,-2.492) (1.755,-2.494) (1.788,-2.495) (1.820,-2.497) (1.853,-2.498) (1.885,-2.500) (1.918,-2.501) (1.950,-2.503) (1.983,-2.504) (2.015,-2.506) (2.047,-2.507) (2.080,-2.509) (2.113,-2.510) (2.145,-2.511) (2.178,-2.513) (2.210,-2.514) (2.243,-2.515) (2.275,-2.516) (2.308,-2.518) (2.340,-2.519) (2.373,-2.520) (2.405,-2.521) (2.438,-2.522) (2.470,-2.523) (2.502,-2.524) (2.535,-2.525) (2.568,-2.527) (2.600,-2.528)}
\LC{(-2.600,2.528) (-2.567,2.527) (-2.535,2.525) (-2.502,2.524) (-2.470,2.523) (-2.438,2.522) (-2.405,2.521) (-2.373,2.520) (-2.340,2.519) (-2.308,2.518) (-2.275,2.516) (-2.243,2.515) (-2.210,2.514) (-2.178,2.513) (-2.145,2.511) (-2.112,2.510) (-2.080,2.509) (-2.047,2.507) (-2.015,2.506) (-1.982,2.504) (-1.950,2.503) (-1.918,2.501) (-1.885,2.500) (-1.853,2.498) (-1.820,2.497) (-1.788,2.495) (-1.755,2.494) (-1.723,2.492) (-1.690,2.490) (-1.657,2.489) (-1.625,2.487) (-1.592,2.485) (-1.560,2.483) (-1.528,2.482) (-1.495,2.480) (-1.462,2.478) (-1.430,2.476) (-1.397,2.474) (-1.365,2.473) (-1.333,2.471) (-1.300,2.469) (-1.268,2.467) (-1.235,2.465) (-1.202,2.463) (-1.170,2.461) (-1.137,2.459) (-1.105,2.457) (-1.073,2.456) (-1.040,2.454) (-1.007,2.452) (-0.975,2.450) (-0.942,2.448) (-0.910,2.446) (-0.877,2.444) (-0.845,2.442) (-0.812,2.441) (-0.780,2.439) (-0.747,2.437) (-0.715,2.435) (-0.682,2.434) (-0.650,2.432) (-0.617,2.430) (-0.585,2.429) (-0.552,2.427) (-0.520,2.426) (-0.487,2.425) (-0.455,2.423) (-0.422,2.422) (-0.390,2.421) (-0.357,2.420) (-0.325,2.419) (-0.292,2.418) (-0.260,2.417) (-0.228,2.417) (-0.195,2.416) (-0.163,2.415) (-0.130,2.415) (-0.097,2.415) (-0.065,2.414) (-0.032,2.414) (0.000,2.414) (0.033,2.414) (0.065,2.414) (0.098,2.415) (0.130,2.415) (0.163,2.415) (0.195,2.416) (0.228,2.417) (0.260,2.417) (0.292,2.418) (0.325,2.419) (0.358,2.420) (0.390,2.421) (0.423,2.422) (0.455,2.423) (0.488,2.425) (0.520,2.426) (0.553,2.427) (0.585,2.429) (0.618,2.430) (0.650,2.432) (0.683,2.434) (0.715,2.435) (0.748,2.437) (0.780,2.439) (0.812,2.441) (0.845,2.442) (0.878,2.444) (0.910,2.446) (0.943,2.448) (0.975,2.450) (1.008,2.452) (1.040,2.454) (1.073,2.456) (1.105,2.457) (1.138,2.459) (1.170,2.461) (1.203,2.463) (1.235,2.465) (1.268,2.467) (1.300,2.469) (1.333,2.471) (1.365,2.473) (1.398,2.474) (1.430,2.476) (1.462,2.478) (1.495,2.480) (1.528,2.482) (1.560,2.483) (1.593,2.485) (1.625,2.487) (1.657,2.489) (1.690,2.490) (1.723,2.492) (1.755,2.494) (1.788,2.495) (1.820,2.497) (1.853,2.498) (1.885,2.500) (1.918,2.501) (1.950,2.503) (1.983,2.504) (2.015,2.506) (2.047,2.507) (2.080,2.509) (2.113,2.510) (2.145,2.511) (2.178,2.513) (2.210,2.514) (2.243,2.515) (2.275,2.516) (2.308,2.518) (2.340,2.519) (2.373,2.520) (2.405,2.521) (2.438,2.522) (2.470,2.523) (2.502,2.524) (2.535,2.525) (2.568,2.527) (2.600,2.528)}
\LC{(1.000,0.105) (1.034,0.108) (1.066,0.111) (1.099,0.114) (1.129,0.117) (1.164,0.120) (1.196,0.123) (1.229,0.125) (1.261,0.128) (1.294,0.130) (1.326,0.132) (1.359,0.135) (1.391,0.137) (1.424,0.139) (1.456,0.141) (1.489,0.142) (1.521,0.144) (1.554,0.146) (1.586,0.147) (1.619,0.149) (1.651,0.150) (1.683,0.151) (1.716,0.153) (1.749,0.154) (1.781,0.155) (1.813,0.156) (1.846,0.157) (1.879,0.158) (1.911,0.159) (1.944,0.160) (1.976,0.161) (2.009,0.161) (2.041,0.162) (2.074,0.163) (2.106,0.163) (2.139,0.164) (2.171,0.165) (2.204,0.165) (2.236,0.166) (2.268,0.166) (2.301,0.167) (2.334,0.167) (2.366,0.168) (2.399,0.168) (2.431,0.169) (2.464,0.169) (2.496,0.169) (2.529,0.170) (2.561,0.170) (2.594,0.171) (2.600,0.171)}
\LC{(-2.600,0.171) (-2.567,0.170) (-2.535,0.170) (-2.502,0.170) (-2.470,0.169) (-2.438,0.169) (-2.405,0.168) (-2.373,0.168) (-2.340,0.167) (-2.308,0.167) (-2.275,0.166) (-2.242,0.166) (-2.210,0.165) (-2.178,0.165) (-2.145,0.164) (-2.112,0.164) (-2.080,0.163) (-2.047,0.162) (-2.015,0.161) (-1.982,0.161) (-1.950,0.160) (-1.917,0.159) (-1.885,0.158) (-1.853,0.157) (-1.820,0.156) (-1.788,0.155) (-1.755,0.154) (-1.723,0.153) (-1.690,0.152) (-1.657,0.150) (-1.625,0.149) (-1.593,0.147) (-1.560,0.146) (-1.527,0.144) (-1.495,0.143) (-1.462,0.141) (-1.430,0.139) (-1.397,0.137) (-1.365,0.135) (-1.333,0.133) (-1.300,0.131) (-1.268,0.128) (-1.235,0.126) (-1.202,0.123) (-1.170,0.121) (-1.137,0.118) (-1.105,0.115) (-1.073,0.112) (-1.040,0.109) (-1.007,0.105) (-1.000,0.105)}
\LC{(1.000,-0.105) (1.034,-0.108) (1.066,-0.111) (1.099,-0.114) (1.129,-0.117) (1.164,-0.120) (1.196,-0.123) (1.229,-0.125) (1.261,-0.128) (1.294,-0.130) (1.326,-0.132) (1.359,-0.135) (1.391,-0.137) (1.424,-0.139) (1.456,-0.141) (1.489,-0.142) (1.521,-0.144) (1.554,-0.146) (1.586,-0.147) (1.619,-0.149) (1.651,-0.150) (1.683,-0.151) (1.716,-0.153) (1.749,-0.154) (1.781,-0.155) (1.814,-0.156) (1.846,-0.157) (1.879,-0.158) (1.911,-0.159) (1.944,-0.160) (1.976,-0.161) (2.009,-0.161) (2.041,-0.162) (2.074,-0.163) (2.106,-0.163) (2.139,-0.164) (2.171,-0.165) (2.204,-0.165) (2.236,-0.166) (2.268,-0.166) (2.301,-0.167) (2.334,-0.167) (2.366,-0.168) (2.399,-0.168) (2.431,-0.169) (2.464,-0.169) (2.496,-0.169) (2.529,-0.170) (2.561,-0.170) (2.594,-0.171) (2.600,-0.171)}
\LC{(-2.600,-0.171) (-2.567,-0.170) (-2.535,-0.170) (-2.502,-0.170) (-2.470,-0.169) (-2.438,-0.169) (-2.405,-0.168) (-2.373,-0.168) (-2.340,-0.167) (-2.308,-0.167) (-2.275,-0.166) (-2.242,-0.166) (-2.210,-0.165) (-2.178,-0.165) (-2.145,-0.164) (-2.112,-0.164) (-2.080,-0.163) (-2.047,-0.162) (-2.015,-0.161) (-1.982,-0.161) (-1.950,-0.160) (-1.917,-0.159) (-1.885,-0.158) (-1.853,-0.157) (-1.820,-0.156) (-1.788,-0.155) (-1.755,-0.154) (-1.723,-0.153) (-1.690,-0.152) (-1.657,-0.150) (-1.625,-0.149) (-1.593,-0.147) (-1.560,-0.146) (-1.528,-0.144) (-1.495,-0.143) (-1.462,-0.141) (-1.430,-0.139) (-1.397,-0.137) (-1.365,-0.135) (-1.333,-0.133) (-1.300,-0.131) (-1.268,-0.128) (-1.235,-0.126) (-1.202,-0.123) (-1.170,-0.121) (-1.137,-0.118) (-1.105,-0.115) (-1.073,-0.112) (-1.040,-0.109) (-1.007,-0.105) (-1.000,-0.105)}
\LC{(0.964,0.278) (0.995,0.287) (1.027,0.295) (1.060,0.304) (1.092,0.312) (1.125,0.320) (1.157,0.327) (1.190,0.334) (1.222,0.341) (1.255,0.348) (1.287,0.354) (1.320,0.360) (1.352,0.365) (1.385,0.370) (1.417,0.375) (1.450,0.380) (1.482,0.385) (1.515,0.389) (1.547,0.393) (1.579,0.397) (1.612,0.400) (1.645,0.404) (1.677,0.407) (1.710,0.410) (1.742,0.413) (1.773,0.416) (1.807,0.419) (1.840,0.421) (1.872,0.424) (1.905,0.426) (1.937,0.428) (1.970,0.430) (2.002,0.432) (2.034,0.434) (2.067,0.436) (2.100,0.438) (2.132,0.439) (2.165,0.441) (2.197,0.442) (2.230,0.444) (2.262,0.445) (2.295,0.446) (2.327,0.448) (2.360,0.449) (2.392,0.450) (2.425,0.451) (2.457,0.452) (2.490,0.453) (2.522,0.454) (2.555,0.455) (2.587,0.456) (2.600,0.456)}
\LC{(-2.600,0.456) (-2.567,0.455) (-2.535,0.454) (-2.502,0.453) (-2.470,0.452) (-2.438,0.451) (-2.405,0.450) (-2.373,0.449) (-2.340,0.448) (-2.308,0.447) (-2.275,0.446) (-2.243,0.444) (-2.210,0.443) (-2.178,0.441) (-2.145,0.440) (-2.112,0.438) (-2.080,0.437) (-2.047,0.435) (-2.015,0.433) (-1.982,0.431) (-1.950,0.429) (-1.917,0.427) (-1.885,0.425) (-1.854,0.422) (-1.820,0.420) (-1.788,0.417) (-1.755,0.414) (-1.723,0.412) (-1.690,0.409) (-1.657,0.405) (-1.625,0.402) (-1.593,0.398) (-1.560,0.395) (-1.528,0.391) (-1.495,0.386) (-1.462,0.382) (-1.430,0.377) (-1.397,0.372) (-1.365,0.367) (-1.333,0.362) (-1.300,0.356) (-1.267,0.350) (-1.235,0.344) (-1.202,0.337) (-1.170,0.330) (-1.137,0.323) (-1.105,0.315) (-1.073,0.307) (-1.040,0.299) (-1.007,0.290) (-0.975,0.281) (-0.964,0.278)}
\LC{(0.964,-0.278) (0.995,-0.287) (1.027,-0.295) (1.060,-0.304) (1.092,-0.312) (1.125,-0.320) (1.157,-0.327) (1.190,-0.334) (1.222,-0.341) (1.255,-0.348) (1.287,-0.354) (1.320,-0.360) (1.352,-0.365) (1.385,-0.370) (1.417,-0.375) (1.450,-0.380) (1.482,-0.385) (1.515,-0.389) (1.547,-0.393) (1.579,-0.397) (1.612,-0.400) (1.645,-0.404) (1.677,-0.407) (1.710,-0.410) (1.742,-0.413) (1.773,-0.416) (1.807,-0.419) (1.840,-0.421) (1.872,-0.424) (1.905,-0.426) (1.937,-0.428) (1.970,-0.430) (2.002,-0.432) (2.034,-0.434) (2.067,-0.436) (2.100,-0.438) (2.132,-0.439) (2.165,-0.441) (2.197,-0.442) (2.230,-0.444) (2.262,-0.445) (2.295,-0.446) (2.327,-0.448) (2.360,-0.449) (2.392,-0.450) (2.425,-0.451) (2.457,-0.452) (2.490,-0.453) (2.522,-0.454) (2.555,-0.455) (2.587,-0.456) (2.600,-0.456)}
\LC{(-2.600,-0.456) (-2.567,-0.455) (-2.535,-0.454) (-2.502,-0.453) (-2.470,-0.452) (-2.438,-0.451) (-2.405,-0.450) (-2.373,-0.449) (-2.340,-0.448) (-2.308,-0.447) (-2.275,-0.446) (-2.243,-0.444) (-2.210,-0.443) (-2.178,-0.441) (-2.145,-0.440) (-2.112,-0.438) (-2.080,-0.437) (-2.047,-0.435) (-2.015,-0.433) (-1.982,-0.431) (-1.950,-0.429) (-1.917,-0.427) (-1.885,-0.425) (-1.854,-0.423) (-1.820,-0.420) (-1.788,-0.417) (-1.755,-0.414) (-1.723,-0.412) (-1.690,-0.409) (-1.657,-0.405) (-1.625,-0.402) (-1.593,-0.398) (-1.560,-0.395) (-1.528,-0.391) (-1.495,-0.386) (-1.462,-0.382) (-1.430,-0.377) (-1.397,-0.372) (-1.365,-0.367) (-1.333,-0.362) (-1.300,-0.356) (-1.267,-0.350) (-1.235,-0.344) (-1.202,-0.337) (-1.170,-0.330) (-1.137,-0.323) (-1.105,-0.315) (-1.073,-0.307) (-1.040,-0.299) (-1.007,-0.290) (-0.975,-0.281) (-0.964,-0.278)}
\LC{(0.894,0.451) (0.923,0.465) (0.956,0.481) (0.985,0.494) (1.014,0.506) (1.047,0.520) (1.079,0.532) (1.112,0.544) (1.144,0.555) (1.175,0.566) (1.209,0.576) (1.239,0.585) (1.274,0.595) (1.307,0.603) (1.337,0.611) (1.372,0.619) (1.404,0.626) (1.437,0.633) (1.469,0.640) (1.502,0.646) (1.534,0.652) (1.567,0.657) (1.599,0.663) (1.632,0.668) (1.664,0.672) (1.696,0.677) (1.729,0.681) (1.762,0.685) (1.794,0.689) (1.827,0.693) (1.859,0.696) (1.892,0.699) (1.924,0.703) (1.957,0.706) (1.989,0.708) (2.022,0.711) (2.054,0.714) (2.087,0.716) (2.119,0.718) (2.151,0.721) (2.184,0.723) (2.216,0.725) (2.249,0.727) (2.282,0.729) (2.314,0.731) (2.346,0.732) (2.379,0.734) (2.412,0.735) (2.444,0.737) (2.477,0.738) (2.509,0.740) (2.540,0.741) (2.574,0.742) (2.600,0.743)}
\LC{(-2.600,0.743) (-2.567,0.742) (-2.535,0.741) (-2.502,0.739) (-2.470,0.738) (-2.438,0.737) (-2.405,0.735) (-2.373,0.734) (-2.340,0.732) (-2.308,0.730) (-2.275,0.728) (-2.243,0.726) (-2.210,0.725) (-2.178,0.722) (-2.145,0.720) (-2.112,0.718) (-2.080,0.716) (-2.047,0.713) (-2.015,0.711) (-1.982,0.708) (-1.950,0.705) (-1.919,0.702) (-1.885,0.699) (-1.854,0.696) (-1.820,0.692) (-1.788,0.688) (-1.755,0.684) (-1.723,0.680) (-1.691,0.676) (-1.657,0.671) (-1.625,0.667) (-1.593,0.662) (-1.561,0.656) (-1.528,0.651) (-1.495,0.645) (-1.462,0.639) (-1.430,0.632) (-1.397,0.625) (-1.365,0.618) (-1.333,0.610) (-1.300,0.602) (-1.268,0.593) (-1.235,0.584) (-1.202,0.574) (-1.170,0.564) (-1.137,0.553) (-1.105,0.542) (-1.073,0.530) (-1.040,0.517) (-1.007,0.504) (-0.975,0.489) (-0.943,0.474) (-0.915,0.462) (-0.894,0.451)}
\LC{(0.894,-0.451) (0.923,-0.465) (0.956,-0.481) (0.985,-0.494) (1.014,-0.506) (1.047,-0.520) (1.079,-0.532) (1.112,-0.544) (1.144,-0.555) (1.175,-0.566) (1.209,-0.576) (1.239,-0.585) (1.274,-0.595) (1.307,-0.603) (1.337,-0.611) (1.372,-0.619) (1.404,-0.626) (1.437,-0.633) (1.469,-0.640) (1.502,-0.646) (1.534,-0.652) (1.567,-0.657) (1.599,-0.663) (1.632,-0.668) (1.664,-0.672) (1.696,-0.677) (1.729,-0.681) (1.762,-0.685) (1.794,-0.689) (1.827,-0.693) (1.859,-0.696) (1.892,-0.699) (1.924,-0.703) (1.957,-0.706) (1.989,-0.708) (2.022,-0.711) (2.054,-0.714) (2.087,-0.716) (2.119,-0.718) (2.151,-0.721) (2.184,-0.723) (2.217,-0.725) (2.249,-0.727) (2.282,-0.729) (2.314,-0.731) (2.347,-0.732) (2.379,-0.734) (2.412,-0.735) (2.444,-0.737) (2.476,-0.738) (2.509,-0.740) (2.540,-0.741) (2.574,-0.742) (2.600,-0.743)}
\LC{(-2.600,-0.743) (-2.567,-0.742) (-2.535,-0.741) (-2.503,-0.739) (-2.470,-0.738) (-2.438,-0.737) (-2.405,-0.735) (-2.373,-0.734) (-2.340,-0.732) (-2.308,-0.730) (-2.275,-0.728) (-2.243,-0.726) (-2.210,-0.725) (-2.178,-0.722) (-2.145,-0.720) (-2.112,-0.718) (-2.080,-0.716) (-2.047,-0.713) (-2.015,-0.711) (-1.983,-0.708) (-1.950,-0.705) (-1.919,-0.702) (-1.885,-0.699) (-1.854,-0.696) (-1.820,-0.692) (-1.788,-0.688) (-1.755,-0.684) (-1.723,-0.680) (-1.691,-0.676) (-1.658,-0.671) (-1.625,-0.667) (-1.592,-0.662) (-1.561,-0.657) (-1.528,-0.651) (-1.495,-0.645) (-1.462,-0.639) (-1.430,-0.632) (-1.397,-0.625) (-1.365,-0.618) (-1.333,-0.610) (-1.300,-0.602) (-1.268,-0.593) (-1.235,-0.584) (-1.202,-0.574) (-1.170,-0.564) (-1.137,-0.553) (-1.105,-0.542) (-1.073,-0.530) (-1.040,-0.517) (-1.007,-0.504) (-0.975,-0.489) (-0.943,-0.475) (-0.915,-0.462) (-0.894,-0.451)}
\LC{(0.785,0.626) (0.812,0.647) (0.839,0.666) (0.865,0.684) (0.891,0.701) (0.917,0.717) (0.943,0.732) (0.971,0.747) (1.001,0.763) (1.034,0.779) (1.065,0.793) (1.095,0.806) (1.125,0.818) (1.157,0.830) (1.190,0.842) (1.219,0.851) (1.248,0.861) (1.281,0.871) (1.313,0.880) (1.346,0.889) (1.378,0.897) (1.411,0.905) (1.443,0.912) (1.476,0.919) (1.508,0.926) (1.541,0.932) (1.573,0.938) (1.606,0.944) (1.637,0.949) (1.671,0.954) (1.703,0.959) (1.736,0.964) (1.768,0.968) (1.800,0.972) (1.833,0.976) (1.866,0.980) (1.898,0.984) (1.930,0.987) (1.963,0.991) (1.996,0.994) (2.028,0.997) (2.061,1.000) (2.093,1.003) (2.126,1.005) (2.158,1.008) (2.191,1.010) (2.223,1.013) (2.256,1.015) (2.288,1.017) (2.321,1.019) (2.353,1.021) (2.385,1.023) (2.418,1.025) (2.451,1.027) (2.483,1.028) (2.516,1.030) (2.548,1.031) (2.581,1.033) (2.600,1.034)}
\LC{(-2.600,1.034) (-2.567,1.032) (-2.535,1.031) (-2.502,1.029) (-2.470,1.028) (-2.438,1.026) (-2.405,1.024) (-2.373,1.022) (-2.342,1.020) (-2.308,1.018) (-2.275,1.016) (-2.243,1.014) (-2.210,1.012) (-2.178,1.009) (-2.145,1.007) (-2.112,1.004) (-2.080,1.002) (-2.047,0.999) (-2.015,0.996) (-1.982,0.993) (-1.950,0.989) (-1.917,0.986) (-1.885,0.983) (-1.853,0.979) (-1.821,0.975) (-1.788,0.971) (-1.755,0.966) (-1.723,0.962) (-1.690,0.957) (-1.657,0.952) (-1.625,0.947) (-1.593,0.941) (-1.562,0.936) (-1.528,0.930) (-1.495,0.923) (-1.463,0.916) (-1.434,0.910) (-1.404,0.903) (-1.371,0.895) (-1.339,0.887) (-1.306,0.878) (-1.274,0.869) (-1.242,0.859) (-1.209,0.848) (-1.180,0.839) (-1.150,0.828) (-1.118,0.815) (-1.085,0.802) (-1.053,0.788) (-1.022,0.773) (-0.994,0.760) (-0.962,0.743) (-0.935,0.728) (-0.903,0.709) (-0.877,0.693) (-0.851,0.675) (-0.825,0.657) (-0.799,0.637) (-0.785,0.626)}
\LC{(0.785,-0.626) (0.812,-0.647) (0.839,-0.666) (0.865,-0.684) (0.891,-0.701) (0.917,-0.717) (0.943,-0.732) (0.971,-0.748) (1.001,-0.763) (1.034,-0.779) (1.065,-0.793) (1.095,-0.806) (1.125,-0.818) (1.157,-0.830) (1.190,-0.842) (1.219,-0.852) (1.248,-0.861) (1.281,-0.871) (1.313,-0.880) (1.346,-0.889) (1.378,-0.897) (1.411,-0.905) (1.443,-0.912) (1.476,-0.919) (1.508,-0.926) (1.541,-0.932) (1.573,-0.938) (1.606,-0.944) (1.637,-0.949) (1.671,-0.954) (1.703,-0.959) (1.736,-0.964) (1.768,-0.968) (1.800,-0.972) (1.833,-0.976) (1.866,-0.980) (1.898,-0.984) (1.930,-0.987) (1.963,-0.991) (1.996,-0.994) (2.028,-0.997) (2.061,-1.000) (2.093,-1.003) (2.126,-1.005) (2.158,-1.008) (2.191,-1.010) (2.223,-1.013) (2.256,-1.015) (2.288,-1.017) (2.321,-1.019) (2.353,-1.021) (2.385,-1.023) (2.418,-1.025) (2.451,-1.027) (2.483,-1.028) (2.516,-1.030) (2.548,-1.031) (2.581,-1.033) (2.600,-1.034)}
\LC{(-2.600,-1.034) (-2.568,-1.032) (-2.535,-1.031) (-2.502,-1.029) (-2.470,-1.028) (-2.438,-1.026) (-2.405,-1.024) (-2.373,-1.022) (-2.342,-1.021) (-2.308,-1.018) (-2.275,-1.016) (-2.243,-1.014) (-2.210,-1.012) (-2.178,-1.009) (-2.145,-1.007) (-2.112,-1.004) (-2.080,-1.002) (-2.047,-0.999) (-2.015,-0.996) (-1.982,-0.993) (-1.950,-0.989) (-1.917,-0.986) (-1.885,-0.983) (-1.853,-0.979) (-1.821,-0.975) (-1.788,-0.971) (-1.755,-0.966) (-1.723,-0.962) (-1.690,-0.957) (-1.657,-0.952) (-1.625,-0.947) (-1.593,-0.941) (-1.562,-0.936) (-1.528,-0.930) (-1.495,-0.923) (-1.463,-0.917) (-1.434,-0.910) (-1.404,-0.903) (-1.371,-0.895) (-1.339,-0.887) (-1.306,-0.878) (-1.274,-0.869) (-1.242,-0.859) (-1.209,-0.848) (-1.180,-0.839) (-1.150,-0.828) (-1.118,-0.815) (-1.085,-0.802) (-1.053,-0.788) (-1.022,-0.774) (-0.994,-0.760) (-0.962,-0.743) (-0.935,-0.728) (-0.903,-0.709) (-0.877,-0.693) (-0.852,-0.675) (-0.825,-0.657) (-0.799,-0.637) (-0.785,-0.626)}
\LC{(0.606,0.798) (0.627,0.825) (0.650,0.851) (0.675,0.877) (0.702,0.903) (0.726,0.923) (0.750,0.942) (0.777,0.962) (0.806,0.982) (0.832,0.998) (0.858,1.013) (0.891,1.031) (0.920,1.046) (0.947,1.059) (0.975,1.072) (1.007,1.085) (1.039,1.099) (1.073,1.111) (1.105,1.122) (1.138,1.133) (1.170,1.143) (1.203,1.153) (1.235,1.162) (1.264,1.170) (1.300,1.179) (1.333,1.187) (1.365,1.194) (1.398,1.201) (1.430,1.208) (1.462,1.214) (1.495,1.221) (1.528,1.226) (1.560,1.232) (1.593,1.237) (1.625,1.242) (1.658,1.247) (1.690,1.252) (1.723,1.256) (1.755,1.261) (1.788,1.265) (1.820,1.269) (1.853,1.272) (1.885,1.276) (1.917,1.279) (1.950,1.283) (1.983,1.286) (2.015,1.289) (2.047,1.292) (2.080,1.295) (2.113,1.297) (2.144,1.300) (2.178,1.303) (2.210,1.305) (2.243,1.307) (2.275,1.309) (2.308,1.312) (2.340,1.314) (2.373,1.316) (2.405,1.318) (2.438,1.319) (2.470,1.321) (2.502,1.323) (2.535,1.325) (2.568,1.326) (2.600,1.328)}
\LC{(-2.600,1.328) (-2.567,1.326) (-2.535,1.325) (-2.502,1.323) (-2.470,1.321) (-2.439,1.319) (-2.405,1.318) (-2.373,1.316) (-2.340,1.314) (-2.308,1.312) (-2.275,1.309) (-2.243,1.307) (-2.210,1.305) (-2.178,1.303) (-2.145,1.300) (-2.112,1.297) (-2.080,1.295) (-2.047,1.292) (-2.015,1.289) (-1.982,1.286) (-1.950,1.283) (-1.917,1.279) (-1.885,1.276) (-1.853,1.272) (-1.820,1.269) (-1.788,1.265) (-1.757,1.261) (-1.723,1.256) (-1.690,1.252) (-1.657,1.247) (-1.625,1.242) (-1.592,1.237) (-1.560,1.232) (-1.528,1.226) (-1.495,1.221) (-1.462,1.214) (-1.430,1.208) (-1.397,1.201) (-1.365,1.194) (-1.333,1.187) (-1.300,1.179) (-1.268,1.171) (-1.235,1.162) (-1.202,1.153) (-1.172,1.144) (-1.137,1.133) (-1.105,1.122) (-1.074,1.111) (-1.040,1.099) (-1.007,1.086) (-0.976,1.073) (-0.947,1.059) (-0.920,1.046) (-0.890,1.031) (-0.859,1.014) (-0.832,0.998) (-0.806,0.982) (-0.780,0.964) (-0.754,0.945) (-0.728,0.925) (-0.703,0.903) (-0.676,0.878) (-0.650,0.851) (-0.627,0.825) (-0.606,0.799) (-0.606,0.798)}
\LC{(0.606,-0.798) (0.627,-0.826) (0.650,-0.851) (0.675,-0.878) (0.702,-0.903) (0.726,-0.923) (0.750,-0.943) (0.777,-0.962) (0.806,-0.982) (0.832,-0.998) (0.858,-1.013) (0.891,-1.031) (0.920,-1.047) (0.947,-1.060) (0.975,-1.072) (1.007,-1.086) (1.039,-1.099) (1.073,-1.111) (1.105,-1.122) (1.138,-1.133) (1.170,-1.143) (1.203,-1.153) (1.235,-1.162) (1.264,-1.170) (1.300,-1.179) (1.333,-1.187) (1.365,-1.194) (1.398,-1.201) (1.430,-1.208) (1.462,-1.214) (1.495,-1.221) (1.528,-1.226) (1.560,-1.232) (1.593,-1.237) (1.625,-1.242) (1.657,-1.247) (1.690,-1.252) (1.723,-1.256) (1.755,-1.261) (1.788,-1.265) (1.820,-1.269) (1.853,-1.272) (1.885,-1.276) (1.917,-1.279) (1.950,-1.283) (1.983,-1.286) (2.015,-1.289) (2.047,-1.292) (2.080,-1.295) (2.113,-1.297) (2.144,-1.300) (2.178,-1.303) (2.210,-1.305) (2.243,-1.307) (2.275,-1.309) (2.308,-1.312) (2.340,-1.314) (2.373,-1.316) (2.405,-1.318) (2.438,-1.319) (2.470,-1.321) (2.502,-1.323) (2.535,-1.325) (2.568,-1.326) (2.600,-1.328)}
\LC{(-2.600,-1.328) (-2.567,-1.326) (-2.535,-1.325) (-2.502,-1.323) (-2.470,-1.321) (-2.439,-1.320) (-2.405,-1.318) (-2.373,-1.316) (-2.340,-1.314) (-2.308,-1.312) (-2.275,-1.309) (-2.243,-1.307) (-2.210,-1.305) (-2.178,-1.303) (-2.145,-1.300) (-2.112,-1.297) (-2.080,-1.295) (-2.047,-1.292) (-2.015,-1.289) (-1.982,-1.286) (-1.950,-1.283) (-1.917,-1.279) (-1.885,-1.276) (-1.853,-1.272) (-1.820,-1.269) (-1.788,-1.265) (-1.757,-1.261) (-1.723,-1.256) (-1.690,-1.252) (-1.657,-1.247) (-1.625,-1.242) (-1.593,-1.237) (-1.560,-1.232) (-1.528,-1.226) (-1.495,-1.221) (-1.462,-1.214) (-1.430,-1.208) (-1.397,-1.201) (-1.365,-1.194) (-1.333,-1.187) (-1.300,-1.179) (-1.268,-1.171) (-1.235,-1.162) (-1.202,-1.153) (-1.172,-1.144) (-1.137,-1.133) (-1.105,-1.122) (-1.074,-1.112) (-1.040,-1.099) (-1.007,-1.086) (-0.976,-1.073) (-0.947,-1.060) (-0.920,-1.047) (-0.890,-1.031) (-0.859,-1.014) (-0.832,-0.998) (-0.806,-0.982) (-0.780,-0.964) (-0.754,-0.945) (-0.728,-0.925) (-0.703,-0.904) (-0.676,-0.878) (-0.650,-0.852) (-0.627,-0.826) (-0.606,-0.799) (-0.606,-0.798)}
\LC{(0.245,0.975) (0.256,1.007) (0.272,1.040) (0.290,1.066) (0.310,1.092) (0.332,1.115) (0.356,1.137) (0.379,1.157) (0.403,1.176) (0.429,1.194) (0.455,1.212) (0.481,1.228) (0.507,1.243) (0.540,1.261) (0.572,1.278) (0.605,1.293) (0.633,1.307) (0.663,1.319) (0.694,1.333) (0.727,1.345) (0.761,1.358) (0.793,1.370) (0.826,1.381) (0.857,1.391) (0.891,1.402) (0.920,1.411) (0.956,1.421) (0.988,1.430) (1.021,1.438) (1.053,1.446) (1.086,1.454) (1.118,1.462) (1.150,1.469) (1.181,1.476) (1.212,1.482) (1.242,1.488) (1.274,1.494) (1.306,1.500) (1.339,1.506) (1.372,1.512) (1.404,1.517) (1.437,1.522) (1.469,1.527) (1.502,1.532) (1.534,1.537) (1.567,1.541) (1.599,1.545) (1.632,1.549) (1.664,1.553) (1.696,1.557) (1.729,1.561) (1.762,1.565) (1.794,1.568) (1.827,1.571) (1.859,1.575) (1.892,1.578) (1.924,1.581) (1.957,1.584) (1.989,1.586) (2.022,1.589) (2.054,1.592) (2.087,1.594) (2.119,1.597) (2.151,1.599) (2.184,1.601) (2.217,1.604) (2.247,1.606) (2.282,1.608) (2.314,1.610) (2.347,1.612) (2.379,1.613) (2.412,1.615) (2.444,1.617) (2.477,1.619) (2.509,1.620) (2.542,1.622) (2.574,1.623) (2.600,1.625)}
\LC{(-2.600,1.625) (-2.567,1.623) (-2.535,1.622) (-2.502,1.620) (-2.470,1.618) (-2.438,1.617) (-2.405,1.615) (-2.373,1.613) (-2.340,1.611) (-2.308,1.609) (-2.275,1.607) (-2.243,1.605) (-2.210,1.603) (-2.178,1.601) (-2.145,1.599) (-2.112,1.596) (-2.080,1.594) (-2.047,1.591) (-2.015,1.589) (-1.984,1.586) (-1.950,1.583) (-1.917,1.580) (-1.885,1.577) (-1.853,1.574) (-1.820,1.571) (-1.788,1.567) (-1.755,1.564) (-1.723,1.560) (-1.690,1.557) (-1.657,1.553) (-1.625,1.549) (-1.592,1.544) (-1.563,1.540) (-1.528,1.536) (-1.495,1.531) (-1.462,1.526) (-1.430,1.521) (-1.397,1.516) (-1.365,1.511) (-1.333,1.505) (-1.300,1.499) (-1.268,1.493) (-1.235,1.487) (-1.202,1.480) (-1.170,1.473) (-1.137,1.466) (-1.105,1.459) (-1.073,1.451) (-1.040,1.443) (-1.007,1.435) (-0.975,1.426) (-0.943,1.417) (-0.910,1.407) (-0.877,1.398) (-0.845,1.387) (-0.812,1.376) (-0.780,1.365) (-0.747,1.353) (-0.715,1.341) (-0.682,1.328) (-0.650,1.314) (-0.619,1.300) (-0.591,1.287) (-0.559,1.271) (-0.528,1.254) (-0.500,1.239) (-0.474,1.224) (-0.449,1.207) (-0.422,1.190) (-0.397,1.171) (-0.371,1.150) (-0.345,1.127) (-0.322,1.105) (-0.300,1.079) (-0.281,1.053) (-0.265,1.027) (-0.251,0.994) (-0.245,0.975)}
\LC{(0.245,-0.975) (0.256,-1.008) (0.272,-1.040) (0.290,-1.066) (0.310,-1.092) (0.332,-1.115) (0.356,-1.138) (0.379,-1.157) (0.403,-1.176) (0.429,-1.194) (0.455,-1.212) (0.481,-1.228) (0.507,-1.243) (0.540,-1.261) (0.572,-1.278) (0.605,-1.293) (0.633,-1.307) (0.663,-1.319) (0.694,-1.333) (0.727,-1.346) (0.761,-1.358) (0.793,-1.370) (0.826,-1.381) (0.857,-1.391) (0.891,-1.402) (0.920,-1.411) (0.956,-1.421) (0.988,-1.430) (1.021,-1.438) (1.053,-1.446) (1.086,-1.454) (1.118,-1.462) (1.150,-1.469) (1.181,-1.476) (1.212,-1.482) (1.242,-1.488) (1.274,-1.494) (1.307,-1.500) (1.339,-1.506) (1.372,-1.512) (1.404,-1.517) (1.437,-1.522) (1.469,-1.527) (1.502,-1.532) (1.534,-1.537) (1.567,-1.541) (1.599,-1.545) (1.632,-1.549) (1.664,-1.553) (1.696,-1.557) (1.729,-1.561) (1.762,-1.565) (1.794,-1.568) (1.827,-1.571) (1.859,-1.575) (1.892,-1.578) (1.924,-1.581) (1.957,-1.584) (1.989,-1.586) (2.022,-1.589) (2.054,-1.592) (2.087,-1.594) (2.119,-1.597) (2.151,-1.599) (2.184,-1.601) (2.216,-1.604) (2.247,-1.606) (2.282,-1.608) (2.314,-1.610) (2.347,-1.612) (2.379,-1.613) (2.412,-1.615) (2.444,-1.617) (2.477,-1.619) (2.509,-1.620) (2.542,-1.622) (2.574,-1.623) (2.600,-1.625)}
\LC{(-2.600,-1.625) (-2.567,-1.623) (-2.535,-1.622) (-2.502,-1.620) (-2.470,-1.618) (-2.438,-1.617) (-2.405,-1.615) (-2.373,-1.613) (-2.340,-1.611) (-2.308,-1.609) (-2.275,-1.607) (-2.243,-1.605) (-2.210,-1.603) (-2.178,-1.601) (-2.145,-1.599) (-2.112,-1.596) (-2.080,-1.594) (-2.047,-1.591) (-2.015,-1.589) (-1.984,-1.586) (-1.950,-1.583) (-1.917,-1.580) (-1.885,-1.577) (-1.853,-1.574) (-1.820,-1.571) (-1.788,-1.567) (-1.755,-1.564) (-1.723,-1.560) (-1.690,-1.557) (-1.657,-1.553) (-1.625,-1.549) (-1.593,-1.544) (-1.563,-1.541) (-1.528,-1.536) (-1.495,-1.531) (-1.462,-1.526) (-1.430,-1.521) (-1.397,-1.516) (-1.365,-1.511) (-1.333,-1.505) (-1.300,-1.499) (-1.268,-1.493) (-1.235,-1.487) (-1.202,-1.480) (-1.170,-1.473) (-1.137,-1.466) (-1.105,-1.459) (-1.073,-1.451) (-1.040,-1.443) (-1.007,-1.435) (-0.975,-1.426) (-0.943,-1.417) (-0.910,-1.407) (-0.877,-1.398) (-0.845,-1.387) (-0.812,-1.376) (-0.780,-1.365) (-0.747,-1.353) (-0.715,-1.341) (-0.682,-1.328) (-0.650,-1.314) (-0.619,-1.300) (-0.591,-1.287) (-0.559,-1.271) (-0.528,-1.255) (-0.500,-1.239) (-0.474,-1.224) (-0.449,-1.207) (-0.422,-1.190) (-0.397,-1.171) (-0.371,-1.151) (-0.345,-1.127) (-0.322,-1.105) (-0.300,-1.079) (-0.281,-1.053) (-0.265,-1.027) (-0.251,-0.995) (-0.245,-0.975)}
\LC{(-2.600,1.924) (-2.567,1.923) (-2.535,1.921) (-2.502,1.920) (-2.470,1.918) (-2.438,1.917) (-2.405,1.915) (-2.373,1.914) (-2.340,1.912) (-2.308,1.910) (-2.275,1.908) (-2.243,1.907) (-2.210,1.905) (-2.178,1.903) (-2.145,1.901) (-2.112,1.899) (-2.080,1.897) (-2.047,1.895) (-2.015,1.892) (-1.982,1.890) (-1.950,1.888) (-1.917,1.885) (-1.885,1.883) (-1.853,1.880) (-1.820,1.878) (-1.788,1.875) (-1.755,1.872) (-1.723,1.869) (-1.690,1.867) (-1.658,1.864) (-1.625,1.860) (-1.593,1.857) (-1.560,1.854) (-1.528,1.851) (-1.495,1.847) (-1.462,1.844) (-1.430,1.840) (-1.397,1.836) (-1.365,1.832) (-1.333,1.828) (-1.300,1.824) (-1.268,1.820) (-1.235,1.816) (-1.202,1.812) (-1.170,1.807) (-1.137,1.803) (-1.105,1.798) (-1.073,1.793) (-1.040,1.788) (-1.007,1.783) (-0.975,1.778) (-0.942,1.773) (-0.911,1.768) (-0.877,1.762) (-0.845,1.757) (-0.812,1.752) (-0.780,1.746) (-0.747,1.740) (-0.715,1.735) (-0.682,1.729) (-0.650,1.723) (-0.617,1.718) (-0.585,1.712) (-0.552,1.706) (-0.520,1.701) (-0.487,1.695) (-0.458,1.690) (-0.422,1.684) (-0.390,1.679) (-0.357,1.674) (-0.325,1.670) (-0.292,1.665) (-0.260,1.661) (-0.228,1.657) (-0.195,1.654) (-0.163,1.651) (-0.130,1.649) (-0.097,1.647) (-0.065,1.646) (-0.032,1.645) (0.000,1.644) (0.033,1.645) (0.065,1.646) (0.098,1.647) (0.130,1.649) (0.163,1.651) (0.195,1.654) (0.228,1.657) (0.260,1.661) (0.292,1.665) (0.325,1.670) (0.358,1.674) (0.390,1.679) (0.423,1.684) (0.455,1.690) (0.488,1.695) (0.520,1.701) (0.553,1.706) (0.585,1.712) (0.618,1.718) (0.650,1.723) (0.682,1.729) (0.715,1.735) (0.748,1.740) (0.780,1.746) (0.812,1.752) (0.845,1.757) (0.878,1.762) (0.910,1.768) (0.943,1.773) (0.975,1.778) (1.008,1.783) (1.040,1.788) (1.073,1.793) (1.105,1.798) (1.138,1.803) (1.169,1.807) (1.203,1.812) (1.235,1.816) (1.266,1.820) (1.300,1.824) (1.333,1.828) (1.365,1.832) (1.398,1.836) (1.430,1.840) (1.462,1.844) (1.495,1.847) (1.528,1.851) (1.560,1.854) (1.593,1.857) (1.625,1.860) (1.658,1.864) (1.690,1.867) (1.723,1.869) (1.755,1.872) (1.788,1.875) (1.820,1.878) (1.853,1.880) (1.885,1.883) (1.917,1.885) (1.950,1.888) (1.983,1.890) (2.015,1.892) (2.047,1.895) (2.080,1.897) (2.113,1.899) (2.145,1.901) (2.178,1.903) (2.210,1.905) (2.243,1.907) (2.275,1.908) (2.308,1.910) (2.340,1.912) (2.373,1.914) (2.405,1.915) (2.438,1.917) (2.470,1.918) (2.502,1.920) (2.535,1.921) (2.568,1.923) (2.600,1.924)}
\LC{(-2.600,-1.924) (-2.567,-1.923) (-2.535,-1.921) (-2.502,-1.920) (-2.470,-1.918) (-2.438,-1.917) (-2.405,-1.915) (-2.373,-1.914) (-2.340,-1.912) (-2.308,-1.910) (-2.275,-1.908) (-2.243,-1.907) (-2.210,-1.905) (-2.178,-1.903) (-2.145,-1.901) (-2.112,-1.899) (-2.080,-1.897) (-2.047,-1.895) (-2.015,-1.892) (-1.982,-1.890) (-1.950,-1.888) (-1.917,-1.885) (-1.885,-1.883) (-1.853,-1.880) (-1.820,-1.878) (-1.788,-1.875) (-1.755,-1.872) (-1.723,-1.869) (-1.690,-1.867) (-1.657,-1.864) (-1.625,-1.860) (-1.593,-1.857) (-1.560,-1.854) (-1.528,-1.851) (-1.495,-1.847) (-1.462,-1.844) (-1.430,-1.840) (-1.397,-1.836) (-1.365,-1.832) (-1.333,-1.828) (-1.300,-1.824) (-1.268,-1.820) (-1.235,-1.816) (-1.202,-1.812) (-1.170,-1.807) (-1.137,-1.803) (-1.105,-1.798) (-1.073,-1.793) (-1.040,-1.788) (-1.007,-1.783) (-0.975,-1.778) (-0.942,-1.773) (-0.911,-1.768) (-0.877,-1.762) (-0.845,-1.757) (-0.812,-1.752) (-0.780,-1.746) (-0.747,-1.740) (-0.715,-1.735) (-0.682,-1.729) (-0.650,-1.723) (-0.617,-1.718) (-0.585,-1.712) (-0.552,-1.706) (-0.520,-1.701) (-0.487,-1.695) (-0.458,-1.690) (-0.422,-1.684) (-0.390,-1.679) (-0.357,-1.674) (-0.325,-1.670) (-0.292,-1.665) (-0.260,-1.661) (-0.228,-1.658) (-0.195,-1.654) (-0.163,-1.651) (-0.130,-1.649) (-0.097,-1.647) (-0.065,-1.646) (-0.032,-1.645) (0.000,-1.644) (0.033,-1.645) (0.065,-1.646) (0.098,-1.647) (0.130,-1.649) (0.163,-1.651) (0.195,-1.654) (0.228,-1.657) (0.260,-1.661) (0.292,-1.665) (0.325,-1.670) (0.358,-1.674) (0.390,-1.679) (0.423,-1.684) (0.455,-1.690) (0.488,-1.695) (0.520,-1.701) (0.553,-1.706) (0.585,-1.712) (0.618,-1.718) (0.650,-1.723) (0.682,-1.729) (0.715,-1.735) (0.748,-1.740) (0.780,-1.746) (0.812,-1.752) (0.845,-1.757) (0.878,-1.762) (0.910,-1.768) (0.943,-1.773) (0.975,-1.778) (1.008,-1.783) (1.040,-1.788) (1.073,-1.793) (1.105,-1.798) (1.138,-1.803) (1.169,-1.807) (1.203,-1.812) (1.235,-1.816) (1.266,-1.820) (1.300,-1.824) (1.333,-1.828) (1.365,-1.832) (1.398,-1.836) (1.430,-1.840) (1.462,-1.844) (1.495,-1.847) (1.528,-1.851) (1.560,-1.854) (1.593,-1.857) (1.625,-1.860) (1.658,-1.864) (1.690,-1.867) (1.723,-1.869) (1.755,-1.872) (1.788,-1.875) (1.820,-1.878) (1.853,-1.880) (1.885,-1.883) (1.917,-1.885) (1.950,-1.888) (1.983,-1.890) (2.015,-1.892) (2.047,-1.895) (2.080,-1.897) (2.113,-1.899) (2.145,-1.901) (2.178,-1.903) (2.210,-1.905) (2.243,-1.907) (2.275,-1.908) (2.308,-1.910) (2.340,-1.912) (2.373,-1.914) (2.405,-1.915) (2.438,-1.917) (2.470,-1.918) (2.502,-1.920) (2.535,-1.921) (2.568,-1.923) (2.600,-1.924)}
\LC{(-2.600,2.225) (-2.567,2.224) (-2.535,2.223) (-2.502,2.221) (-2.470,2.220) (-2.438,2.219) (-2.405,2.218) (-2.373,2.216) (-2.340,2.215) (-2.308,2.213) (-2.275,2.212) (-2.243,2.210) (-2.210,2.209) (-2.178,2.207) (-2.145,2.206) (-2.112,2.204) (-2.080,2.202) (-2.048,2.200) (-2.015,2.199) (-1.982,2.197) (-1.950,2.195) (-1.917,2.193) (-1.885,2.191) (-1.853,2.189) (-1.820,2.187) (-1.788,2.185) (-1.755,2.183) (-1.723,2.181) (-1.690,2.178) (-1.657,2.176) (-1.625,2.174) (-1.593,2.172) (-1.560,2.169) (-1.528,2.167) (-1.495,2.164) (-1.462,2.162) (-1.430,2.159) (-1.397,2.156) (-1.365,2.154) (-1.333,2.151) (-1.300,2.148) (-1.268,2.146) (-1.235,2.143) (-1.202,2.140) (-1.170,2.137) (-1.137,2.134) (-1.105,2.131) (-1.073,2.128) (-1.040,2.125) (-1.007,2.122) (-0.975,2.119) (-0.942,2.116) (-0.910,2.113) (-0.877,2.110) (-0.845,2.107) (-0.812,2.104) (-0.780,2.101) (-0.747,2.098) (-0.715,2.095) (-0.682,2.092) (-0.650,2.089) (-0.618,2.087) (-0.585,2.084) (-0.552,2.081) (-0.520,2.078) (-0.487,2.076) (-0.455,2.074) (-0.422,2.071) (-0.390,2.069) (-0.357,2.067) (-0.325,2.065) (-0.292,2.063) (-0.260,2.062) (-0.228,2.060) (-0.195,2.059) (-0.163,2.058) (-0.130,2.057) (-0.097,2.056) (-0.065,2.056) (-0.032,2.056) (0.000,2.056) (0.033,2.056) (0.065,2.056) (0.098,2.056) (0.130,2.057) (0.163,2.058) (0.195,2.059) (0.228,2.060) (0.260,2.062) (0.292,2.063) (0.325,2.065) (0.357,2.067) (0.390,2.069) (0.423,2.071) (0.454,2.074) (0.488,2.076) (0.520,2.078) (0.553,2.081) (0.585,2.084) (0.618,2.086) (0.650,2.089) (0.683,2.092) (0.715,2.095) (0.748,2.098) (0.780,2.101) (0.812,2.104) (0.845,2.107) (0.878,2.110) (0.910,2.113) (0.943,2.116) (0.973,2.119) (1.008,2.122) (1.040,2.125) (1.073,2.128) (1.105,2.131) (1.138,2.134) (1.170,2.137) (1.203,2.140) (1.235,2.143) (1.268,2.146) (1.300,2.148) (1.333,2.151) (1.365,2.154) (1.398,2.156) (1.430,2.159) (1.462,2.162) (1.495,2.164) (1.528,2.167) (1.560,2.169) (1.593,2.172) (1.625,2.174) (1.658,2.176) (1.690,2.178) (1.723,2.181) (1.755,2.183) (1.788,2.185) (1.820,2.187) (1.853,2.189) (1.885,2.191) (1.917,2.193) (1.950,2.195) (1.983,2.197) (2.015,2.199) (2.048,2.200) (2.080,2.202) (2.113,2.204) (2.145,2.206) (2.178,2.207) (2.210,2.209) (2.243,2.210) (2.275,2.212) (2.308,2.213) (2.340,2.215) (2.373,2.216) (2.405,2.218) (2.438,2.219) (2.470,2.220) (2.502,2.221) (2.535,2.223) (2.568,2.224) (2.600,2.225)}
\LC{(-2.600,-2.225) (-2.567,-2.224) (-2.535,-2.223) (-2.502,-2.221) (-2.470,-2.220) (-2.438,-2.219) (-2.405,-2.218) (-2.372,-2.216) (-2.340,-2.215) (-2.308,-2.213) (-2.275,-2.212) (-2.243,-2.210) (-2.210,-2.209) (-2.178,-2.207) (-2.145,-2.206) (-2.112,-2.204) (-2.080,-2.202) (-2.047,-2.200) (-2.015,-2.199) (-1.982,-2.197) (-1.950,-2.195) (-1.917,-2.193) (-1.885,-2.191) (-1.853,-2.189) (-1.820,-2.187) (-1.788,-2.185) (-1.755,-2.183) (-1.723,-2.181) (-1.690,-2.178) (-1.658,-2.176) (-1.625,-2.174) (-1.593,-2.172) (-1.560,-2.169) (-1.527,-2.167) (-1.495,-2.164) (-1.462,-2.162) (-1.430,-2.159) (-1.397,-2.156) (-1.365,-2.154) (-1.333,-2.151) (-1.300,-2.148) (-1.268,-2.146) (-1.235,-2.143) (-1.202,-2.140) (-1.170,-2.137) (-1.137,-2.134) (-1.105,-2.131) (-1.073,-2.128) (-1.040,-2.125) (-1.007,-2.122) (-0.975,-2.119) (-0.942,-2.116) (-0.910,-2.113) (-0.877,-2.110) (-0.845,-2.107) (-0.812,-2.104) (-0.780,-2.101) (-0.747,-2.098) (-0.715,-2.095) (-0.682,-2.092) (-0.650,-2.089) (-0.618,-2.087) (-0.585,-2.084) (-0.552,-2.081) (-0.520,-2.078) (-0.487,-2.076) (-0.455,-2.074) (-0.422,-2.071) (-0.390,-2.069) (-0.357,-2.067) (-0.325,-2.065) (-0.292,-2.063) (-0.260,-2.062) (-0.228,-2.060) (-0.195,-2.059) (-0.163,-2.058) (-0.130,-2.057) (-0.097,-2.056) (-0.065,-2.056) (-0.032,-2.056) (0.000,-2.056) (0.033,-2.056) (0.065,-2.056) (0.098,-2.056) (0.130,-2.057) (0.163,-2.058) (0.195,-2.059) (0.228,-2.060) (0.260,-2.062) (0.292,-2.063) (0.325,-2.065) (0.357,-2.067) (0.390,-2.069) (0.423,-2.071) (0.454,-2.074) (0.488,-2.076) (0.520,-2.078) (0.553,-2.081) (0.585,-2.084) (0.618,-2.086) (0.650,-2.089) (0.683,-2.092) (0.715,-2.095) (0.748,-2.098) (0.780,-2.101) (0.812,-2.104) (0.845,-2.107) (0.878,-2.110) (0.910,-2.113) (0.943,-2.116) (0.973,-2.119) (1.008,-2.122) (1.040,-2.125) (1.073,-2.128) (1.105,-2.131) (1.138,-2.134) (1.170,-2.137) (1.203,-2.140) (1.235,-2.143) (1.268,-2.146) (1.300,-2.148) (1.333,-2.151) (1.365,-2.154) (1.398,-2.156) (1.430,-2.159) (1.462,-2.162) (1.495,-2.164) (1.528,-2.167) (1.560,-2.169) (1.593,-2.172) (1.625,-2.174) (1.658,-2.176) (1.690,-2.178) (1.723,-2.181) (1.755,-2.183) (1.788,-2.185) (1.820,-2.187) (1.853,-2.189) (1.885,-2.191) (1.917,-2.193) (1.950,-2.195) (1.983,-2.197) (2.015,-2.199) (2.047,-2.200) (2.080,-2.202) (2.113,-2.204) (2.145,-2.206) (2.178,-2.207) (2.210,-2.209) (2.243,-2.210) (2.275,-2.212) (2.308,-2.213) (2.340,-2.215) (2.372,-2.216) (2.405,-2.218) (2.438,-2.219) (2.470,-2.220) (2.502,-2.221) (2.535,-2.223) (2.568,-2.224) (2.600,-2.225)}
\LC{(-2.600,2.528) (-2.567,2.527) (-2.535,2.525) (-2.502,2.524) (-2.470,2.523) (-2.438,2.522) (-2.405,2.521) (-2.373,2.520) (-2.340,2.519) (-2.308,2.518) (-2.275,2.516) (-2.243,2.515) (-2.210,2.514) (-2.178,2.513) (-2.145,2.511) (-2.112,2.510) (-2.080,2.509) (-2.047,2.507) (-2.015,2.506) (-1.982,2.504) (-1.950,2.503) (-1.918,2.501) (-1.885,2.500) (-1.853,2.498) (-1.820,2.497) (-1.788,2.495) (-1.755,2.494) (-1.723,2.492) (-1.690,2.490) (-1.657,2.489) (-1.625,2.487) (-1.592,2.485) (-1.560,2.483) (-1.528,2.482) (-1.495,2.480) (-1.462,2.478) (-1.430,2.476) (-1.397,2.474) (-1.365,2.473) (-1.333,2.471) (-1.300,2.469) (-1.268,2.467) (-1.235,2.465) (-1.203,2.463) (-1.170,2.461) (-1.137,2.459) (-1.105,2.457) (-1.073,2.456) (-1.040,2.454) (-1.007,2.452) (-0.975,2.450) (-0.942,2.448) (-0.910,2.446) (-0.877,2.444) (-0.845,2.442) (-0.812,2.441) (-0.780,2.439) (-0.747,2.437) (-0.715,2.435) (-0.682,2.434) (-0.650,2.432) (-0.617,2.430) (-0.585,2.429) (-0.552,2.427) (-0.520,2.426) (-0.487,2.425) (-0.455,2.423) (-0.422,2.422) (-0.390,2.421) (-0.357,2.420) (-0.325,2.419) (-0.292,2.418) (-0.260,2.417) (-0.228,2.417) (-0.195,2.416) (-0.163,2.415) (-0.130,2.415) (-0.097,2.415) (-0.065,2.414) (-0.032,2.414) (0.000,2.414) (0.033,2.414) (0.065,2.414) (0.098,2.415) (0.130,2.415) (0.163,2.415) (0.195,2.416) (0.228,2.417) (0.260,2.417) (0.292,2.418) (0.325,2.419) (0.358,2.420) (0.390,2.421) (0.423,2.422) (0.455,2.423) (0.488,2.425) (0.520,2.426) (0.553,2.427) (0.585,2.429) (0.618,2.430) (0.650,2.432) (0.683,2.434) (0.715,2.435) (0.748,2.437) (0.780,2.439) (0.812,2.441) (0.845,2.442) (0.878,2.444) (0.910,2.446) (0.943,2.448) (0.975,2.450) (1.008,2.452) (1.040,2.454) (1.073,2.456) (1.105,2.457) (1.138,2.459) (1.170,2.461) (1.203,2.463) (1.235,2.465) (1.268,2.467) (1.300,2.469) (1.333,2.471) (1.365,2.473) (1.398,2.474) (1.430,2.476) (1.462,2.478) (1.495,2.480) (1.528,2.482) (1.560,2.483) (1.593,2.485) (1.625,2.487) (1.657,2.489) (1.690,2.490) (1.723,2.492) (1.755,2.494) (1.788,2.495) (1.820,2.497) (1.853,2.498) (1.885,2.500) (1.918,2.501) (1.950,2.503) (1.983,2.504) (2.015,2.506) (2.047,2.507) (2.080,2.509) (2.113,2.510) (2.145,2.511) (2.178,2.513) (2.210,2.514) (2.243,2.515) (2.275,2.516) (2.308,2.518) (2.340,2.519) (2.373,2.520) (2.405,2.521) (2.438,2.522) (2.470,2.523) (2.502,2.524) (2.535,2.525) (2.568,2.527) (2.600,2.528)}
\LC{(-2.600,-2.528) (-2.567,-2.527) (-2.535,-2.525) (-2.502,-2.524) (-2.470,-2.523) (-2.438,-2.522) (-2.405,-2.521) (-2.373,-2.520) (-2.340,-2.519) (-2.308,-2.518) (-2.275,-2.516) (-2.243,-2.515) (-2.210,-2.514) (-2.178,-2.513) (-2.145,-2.511) (-2.112,-2.510) (-2.080,-2.509) (-2.047,-2.507) (-2.015,-2.506) (-1.982,-2.504) (-1.950,-2.503) (-1.918,-2.501) (-1.885,-2.500) (-1.853,-2.498) (-1.820,-2.497) (-1.788,-2.495) (-1.755,-2.494) (-1.723,-2.492) (-1.690,-2.490) (-1.657,-2.489) (-1.625,-2.487) (-1.592,-2.485) (-1.560,-2.483) (-1.528,-2.482) (-1.495,-2.480) (-1.462,-2.478) (-1.430,-2.476) (-1.397,-2.474) (-1.365,-2.473) (-1.333,-2.471) (-1.300,-2.469) (-1.268,-2.467) (-1.235,-2.465) (-1.202,-2.463) (-1.170,-2.461) (-1.137,-2.459) (-1.105,-2.457) (-1.073,-2.456) (-1.040,-2.454) (-1.007,-2.452) (-0.975,-2.450) (-0.942,-2.448) (-0.910,-2.446) (-0.877,-2.444) (-0.845,-2.442) (-0.812,-2.441) (-0.780,-2.439) (-0.747,-2.437) (-0.715,-2.435) (-0.682,-2.434) (-0.650,-2.432) (-0.617,-2.430) (-0.585,-2.429) (-0.552,-2.427) (-0.520,-2.426) (-0.487,-2.425) (-0.455,-2.423) (-0.422,-2.422) (-0.390,-2.421) (-0.357,-2.420) (-0.325,-2.419) (-0.292,-2.418) (-0.260,-2.417) (-0.228,-2.417) (-0.195,-2.416) (-0.163,-2.415) (-0.130,-2.415) (-0.097,-2.415) (-0.065,-2.414) (-0.032,-2.414) (0.000,-2.414) (0.033,-2.414) (0.065,-2.414) (0.098,-2.415) (0.130,-2.415) (0.163,-2.415) (0.195,-2.416) (0.228,-2.417) (0.260,-2.417) (0.292,-2.418) (0.325,-2.419) (0.358,-2.420) (0.390,-2.421) (0.423,-2.422) (0.455,-2.423) (0.488,-2.425) (0.520,-2.426) (0.553,-2.427) (0.585,-2.429) (0.618,-2.430) (0.650,-2.432) (0.683,-2.434) (0.715,-2.435) (0.748,-2.437) (0.780,-2.439) (0.812,-2.441) (0.845,-2.442) (0.878,-2.444) (0.910,-2.446) (0.943,-2.448) (0.975,-2.450) (1.008,-2.452) (1.040,-2.454) (1.073,-2.456) (1.105,-2.457) (1.138,-2.459) (1.170,-2.461) (1.203,-2.463) (1.235,-2.465) (1.268,-2.467) (1.300,-2.469) (1.333,-2.471) (1.365,-2.473) (1.398,-2.474) (1.430,-2.476) (1.462,-2.478) (1.495,-2.480) (1.528,-2.482) (1.560,-2.483) (1.593,-2.485) (1.625,-2.487) (1.657,-2.489) (1.690,-2.490) (1.723,-2.492) (1.755,-2.494) (1.788,-2.495) (1.820,-2.497) (1.853,-2.498) (1.885,-2.500) (1.918,-2.501) (1.950,-2.503) (1.983,-2.504) (2.015,-2.506) (2.047,-2.507) (2.080,-2.509) (2.113,-2.510) (2.145,-2.511) (2.178,-2.513) (2.210,-2.514) (2.243,-2.515) (2.275,-2.516) (2.308,-2.518) (2.340,-2.519) (2.373,-2.520) (2.405,-2.521) (2.438,-2.522) (2.470,-2.523) (2.502,-2.524) (2.535,-2.525) (2.568,-2.527) (2.600,-2.528)}
\LC{(-2.600,0.000) (-1.000,0.000)}
\LC{(1.000,0.000) (2.600,0.000)}

      % Esfera en campo uniforme - flechas
% Generado por scripts/generar_datos.py (no editar a mano).
\FL{1.898}{-0.158}{-1.6}
\FL{-1.898}{-0.158}{1.6}
\FL{1.898}{0.158}{1.6}
\FL{-1.898}{0.158}{-1.6}
\FL{1.898}{-0.426}{-4.0}
\FL{-1.898}{-0.426}{4.0}
\FL{1.898}{0.426}{4.0}
\FL{-1.898}{0.426}{-4.0}
\FL{1.898}{-0.700}{-5.6}
\FL{-1.898}{-0.700}{5.6}
\FL{1.898}{0.700}{5.6}
\FL{-1.898}{0.700}{-5.6}
\FL{1.898}{-0.984}{-6.3}
\FL{-1.898}{-0.984}{6.3}
\FL{1.898}{0.984}{6.3}
\FL{-1.898}{0.984}{-6.3}
\FL{1.898}{-1.277}{-6.1}
\FL{-1.898}{-1.277}{6.1}
\FL{1.898}{1.277}{6.1}
\FL{-1.898}{1.277}{-6.1}
\FL{1.898}{-1.578}{-5.3}
\FL{-1.898}{-1.578}{5.3}
\FL{1.898}{1.578}{5.3}
\FL{-1.898}{1.578}{-5.3}
\FL{-1.898}{-1.884}{4.4}
\FL{1.898}{-1.884}{-4.4}
\FL{-1.898}{1.884}{-4.4}
\FL{1.898}{1.884}{4.4}
\FL{-1.898}{-2.192}{3.4}
\FL{1.898}{-2.192}{-3.4}
\FL{-1.898}{2.192}{-3.4}
\FL{1.898}{2.192}{3.4}
\FL{-1.898}{-2.500}{2.7}
\FL{1.898}{-2.500}{-2.7}
\FL{-1.898}{2.500}{-2.7}
\FL{1.898}{2.500}{2.7}
\FL{1.898}{0.158}{1.6}
\FL{-1.898}{0.158}{-1.6}
\FL{1.898}{-0.158}{-1.6}
\FL{-1.898}{-0.158}{1.6}
\FL{1.898}{0.426}{4.0}
\FL{-1.898}{0.426}{-4.0}
\FL{1.898}{-0.426}{-4.0}
\FL{-1.898}{-0.426}{4.0}
\FL{1.898}{0.700}{5.6}
\FL{-1.898}{0.700}{-5.6}
\FL{1.898}{-0.700}{-5.6}
\FL{-1.898}{-0.700}{5.6}
\FL{1.898}{0.984}{6.3}
\FL{-1.898}{0.984}{-6.3}
\FL{1.898}{-0.984}{-6.3}
\FL{-1.898}{-0.984}{6.3}
\FL{1.898}{1.277}{6.1}
\FL{-1.898}{1.277}{-6.1}
\FL{1.898}{-1.277}{-6.1}
\FL{-1.898}{-1.277}{6.1}
\FL{1.898}{1.578}{5.3}
\FL{-1.898}{1.578}{-5.3}
\FL{1.898}{-1.578}{-5.3}
\FL{-1.898}{-1.578}{5.3}
\FL{-1.898}{1.884}{-4.4}
\FL{1.898}{1.884}{4.4}
\FL{-1.898}{-1.884}{4.4}
\FL{1.898}{-1.884}{-4.4}
\FL{-1.898}{2.192}{-3.4}
\FL{1.898}{2.192}{3.4}
\FL{-1.898}{-2.192}{3.4}
\FL{1.898}{-2.192}{-3.4}
\FL{-1.898}{2.500}{-2.7}
\FL{1.898}{2.500}{2.7}
\FL{-1.898}{-2.500}{2.7}
\FL{1.898}{-2.500}{-2.7}
\FL{-2.600}{0.000}{0.0}
\FL{2.600}{0.000}{0.0}

    \end{scope}
    \shade[ball color=azulvelo, draw=azulprof] (0,0) circle (1);
    \draw[azulprof, line width=0.6pt] (0,0) circle (1);
    % Esfera en campo uniforme - cargas inducidas
% Generado por scripts/generar_datos.py (no editar a mano).
\CP{1.123}{0.128}
\CP{1.064}{0.379}
\CP{0.939}{0.629}
\CP{0.709}{0.880}
\CP{0.060}{1.128}
\CP{0.712}{-0.877}
\CP{0.941}{-0.626}
\CP{1.066}{-0.376}
\CP{1.123}{-0.124}
\CN{-0.519}{1.004}
\CN{-0.845}{0.750}
\CN{-1.013}{0.500}
\CN{-1.103}{0.247}
\CN{-1.130}{-0.004}
\CN{-1.101}{-0.253}
\CN{-1.010}{-0.507}
\CN{-0.841}{-0.755}
\CN{-0.513}{-1.007}

    \node[font=\footnotesize, text=azulprof] at (0,0) {$\vE=\vc{0}$};
    \draw[azulclaro, line width=0.5pt] (-2.6,-2.6) rectangle (2.6,2.6);
    \node[anchor=north] at (0,-2.75) {(a) Esfera conductora};
  \end{scope}
  % ---------------------------------------------------------------- (b)
  \begin{scope}[xshift=6.45cm, yshift=-0.35cm, x=1.17cm, y=1.17cm]
    \fill[azulfondo!60!white] (-2.3,-1.9) rectangle (2.3,1.9);
    \begin{scope}
      \clip (-2.3,-1.9) rectangle (2.3,1.9);
      \tikzset{numEQ/.style={azulcielo, line width=0.35pt, densely dashed}}
      % Conductor neutro en campo uniforme - equipotenciales
% Generado por scripts/generar_datos.py (no editar a mano).
\EQ{(2.165,-1.940) (2.168,-1.900) (2.171,-1.860) (2.174,-1.820) (2.177,-1.780) (2.180,-1.740) (2.183,-1.700) (2.186,-1.660) (2.189,-1.620) (2.192,-1.580) (2.195,-1.540) (2.199,-1.500) (2.202,-1.460) (2.205,-1.420) (2.209,-1.380) (2.212,-1.340) (2.216,-1.300) (2.219,-1.260) (2.223,-1.220) (2.226,-1.180) (2.230,-1.140) (2.234,-1.100) (2.237,-1.060) (2.240,-1.029) (2.245,-0.980) (2.248,-0.940) (2.252,-0.900) (2.256,-0.860) (2.259,-0.820) (2.263,-0.780) (2.266,-0.740) (2.270,-0.700) (2.273,-0.660) (2.277,-0.620) (2.280,-0.584) (2.284,-0.540) (2.287,-0.500) (2.290,-0.460) (2.293,-0.420) (2.296,-0.380) (2.299,-0.340) (2.302,-0.300) (2.305,-0.260) (2.308,-0.220) (2.310,-0.180) (2.313,-0.140) (2.315,-0.100) (2.317,-0.060) (2.319,-0.020) (2.321,0.020) (2.322,0.060) (2.324,0.100) (2.325,0.140) (2.326,0.180) (2.327,0.220) (2.328,0.260) (2.328,0.300) (2.328,0.340) (2.329,0.380) (2.329,0.420) (2.328,0.460) (2.328,0.500) (2.327,0.540) (2.326,0.580) (2.325,0.620) (2.324,0.660) (2.323,0.700) (2.321,0.740) (2.319,0.780) (2.318,0.820) (2.316,0.860) (2.313,0.900) (2.311,0.940) (2.308,0.980) (2.306,1.020) (2.303,1.060) (2.300,1.100) (2.297,1.140) (2.294,1.180) (2.291,1.220) (2.288,1.260) (2.284,1.300) (2.281,1.340) (2.277,1.380) (2.274,1.420) (2.270,1.460) (2.266,1.500) (2.263,1.540) (2.260,1.570) (2.255,1.620) (2.252,1.660) (2.248,1.700) (2.244,1.740) (2.240,1.780) (2.237,1.820) (2.233,1.860) (2.229,1.900) (2.225,1.940)}
\EQ{(1.870,-1.940) (1.873,-1.900) (1.876,-1.860) (1.879,-1.820) (1.883,-1.780) (1.886,-1.740) (1.890,-1.700) (1.893,-1.660) (1.897,-1.620) (1.900,-1.586) (1.904,-1.540) (1.908,-1.500) (1.912,-1.460) (1.916,-1.420) (1.920,-1.386) (1.925,-1.340) (1.929,-1.300) (1.933,-1.260) (1.938,-1.220) (1.942,-1.180) (1.947,-1.140) (1.951,-1.100) (1.956,-1.060) (1.960,-1.026) (1.965,-0.980) (1.970,-0.940) (1.975,-0.900) (1.980,-0.860) (1.984,-0.820) (1.989,-0.780) (1.994,-0.740) (1.999,-0.700) (2.003,-0.660) (2.008,-0.620) (2.013,-0.580) (2.017,-0.540) (2.022,-0.500) (2.026,-0.460) (2.031,-0.420) (2.035,-0.380) (2.039,-0.340) (2.043,-0.300) (2.047,-0.260) (2.051,-0.220) (2.054,-0.180) (2.058,-0.140) (2.060,-0.109) (2.064,-0.060) (2.066,-0.020) (2.069,0.020) (2.071,0.060) (2.073,0.100) (2.075,0.140) (2.077,0.180) (2.078,0.220) (2.079,0.260) (2.080,0.297) (2.081,0.340) (2.081,0.380) (2.081,0.420) (2.080,0.460) (2.080,0.500) (2.079,0.540) (2.078,0.580) (2.076,0.620) (2.075,0.660) (2.073,0.700) (2.070,0.740) (2.068,0.780) (2.065,0.820) (2.062,0.860) (2.059,0.900) (2.056,0.940) (2.052,0.980) (2.048,1.020) (2.045,1.060) (2.040,1.100) (2.036,1.140) (2.032,1.180) (2.028,1.220) (2.023,1.260) (2.018,1.300) (2.014,1.340) (2.009,1.380) (2.004,1.420) (2.000,1.455) (1.994,1.500) (1.990,1.540) (1.985,1.580) (1.980,1.618) (1.975,1.660) (1.970,1.700) (1.965,1.740) (1.960,1.780) (1.956,1.820) (1.951,1.860) (1.946,1.900) (1.942,1.940)}
\EQ{(1.569,-1.940) (1.572,-1.900) (1.576,-1.860) (1.580,-1.820) (1.583,-1.780) (1.587,-1.740) (1.591,-1.700) (1.595,-1.660) (1.600,-1.620) (1.604,-1.580) (1.608,-1.540) (1.613,-1.500) (1.618,-1.460) (1.622,-1.420) (1.627,-1.380) (1.632,-1.340) (1.638,-1.300) (1.643,-1.260) (1.648,-1.220) (1.654,-1.180) (1.659,-1.140) (1.665,-1.100) (1.671,-1.060) (1.677,-1.020) (1.683,-0.980) (1.689,-0.940) (1.695,-0.900) (1.700,-0.868) (1.708,-0.820) (1.714,-0.780) (1.720,-0.741) (1.727,-0.700) (1.733,-0.660) (1.739,-0.620) (1.746,-0.580) (1.752,-0.540) (1.758,-0.500) (1.764,-0.460) (1.770,-0.420) (1.776,-0.380) (1.782,-0.340) (1.788,-0.300) (1.793,-0.260) (1.799,-0.220) (1.804,-0.180) (1.809,-0.140) (1.813,-0.100) (1.818,-0.060) (1.822,-0.020) (1.826,0.020) (1.829,0.060) (1.832,0.100) (1.835,0.140) (1.837,0.180) (1.839,0.220) (1.841,0.260) (1.842,0.300) (1.843,0.340) (1.843,0.380) (1.843,0.420) (1.843,0.460) (1.842,0.500) (1.841,0.540) (1.839,0.580) (1.837,0.620) (1.834,0.660) (1.831,0.700) (1.828,0.740) (1.825,0.780) (1.821,0.820) (1.816,0.860) (1.812,0.900) (1.807,0.940) (1.801,0.980) (1.796,1.020) (1.790,1.060) (1.785,1.100) (1.780,1.131) (1.772,1.180) (1.766,1.220) (1.760,1.259) (1.753,1.300) (1.747,1.340) (1.740,1.380) (1.734,1.420) (1.727,1.460) (1.721,1.500) (1.714,1.540) (1.708,1.580) (1.701,1.620) (1.695,1.660) (1.689,1.700) (1.683,1.740) (1.676,1.780) (1.670,1.820) (1.665,1.860) (1.660,1.892) (1.653,1.940)}
\EQ{(1.262,-1.940) (1.266,-1.900) (1.269,-1.860) (1.273,-1.820) (1.277,-1.780) (1.281,-1.740) (1.286,-1.700) (1.290,-1.660) (1.295,-1.620) (1.300,-1.580) (1.305,-1.540) (1.310,-1.500) (1.315,-1.460) (1.320,-1.424) (1.326,-1.380) (1.332,-1.340) (1.338,-1.300) (1.344,-1.260) (1.351,-1.220) (1.357,-1.180) (1.364,-1.140) (1.371,-1.100) (1.378,-1.060) (1.386,-1.020) (1.393,-0.980) (1.400,-0.945) (1.409,-0.900) (1.417,-0.860) (1.425,-0.820) (1.433,-0.780) (1.440,-0.747) (1.450,-0.700) (1.458,-0.660) (1.467,-0.620) (1.476,-0.580) (1.484,-0.540) (1.493,-0.500) (1.500,-0.467) (1.510,-0.420) (1.518,-0.380) (1.527,-0.340) (1.535,-0.300) (1.543,-0.260) (1.551,-0.220) (1.559,-0.180) (1.566,-0.140) (1.573,-0.100) (1.580,-0.060) (1.586,-0.020) (1.592,0.020) (1.597,0.060) (1.602,0.100) (1.606,0.140) (1.610,0.180) (1.614,0.220) (1.616,0.260) (1.618,0.300) (1.620,0.340) (1.620,0.380) (1.621,0.420) (1.620,0.454) (1.619,0.500) (1.617,0.540) (1.614,0.580) (1.611,0.620) (1.607,0.660) (1.602,0.700) (1.597,0.740) (1.591,0.780) (1.585,0.820) (1.580,0.852) (1.572,0.900) (1.564,0.940) (1.556,0.980) (1.548,1.020) (1.540,1.058) (1.531,1.100) (1.522,1.140) (1.513,1.180) (1.504,1.220) (1.495,1.260) (1.486,1.300) (1.477,1.340) (1.468,1.380) (1.460,1.415) (1.450,1.460) (1.441,1.500) (1.433,1.540) (1.424,1.580) (1.416,1.620) (1.408,1.660) (1.400,1.699) (1.392,1.740) (1.384,1.780) (1.377,1.820) (1.370,1.860) (1.363,1.900) (1.356,1.940)}
\EQ{(0.948,-1.940) (0.951,-1.900) (0.954,-1.860) (0.958,-1.820) (0.962,-1.780) (0.966,-1.740) (0.970,-1.700) (0.975,-1.660) (0.980,-1.620) (0.985,-1.580) (0.990,-1.540) (0.995,-1.500) (1.000,-1.466) (1.007,-1.420) (1.013,-1.380) (1.019,-1.340) (1.026,-1.300) (1.033,-1.260) (1.040,-1.222) (1.048,-1.180) (1.056,-1.140) (1.064,-1.100) (1.072,-1.060) (1.080,-1.026) (1.090,-0.980) (1.100,-0.940) (1.109,-0.900) (1.119,-0.860) (1.130,-0.820) (1.140,-0.781) (1.151,-0.740) (1.160,-0.707) (1.173,-0.660) (1.185,-0.620) (1.196,-0.580) (1.208,-0.540) (1.220,-0.501) (1.232,-0.460) (1.244,-0.420) (1.256,-0.380) (1.269,-0.340) (1.280,-0.302) (1.293,-0.260) (1.304,-0.220) (1.316,-0.180) (1.327,-0.140) (1.338,-0.100) (1.349,-0.060) (1.359,-0.020) (1.368,0.020) (1.377,0.060) (1.385,0.100) (1.392,0.140) (1.399,0.180) (1.404,0.220) (1.409,0.260) (1.413,0.300) (1.415,0.340) (1.417,0.380) (1.417,0.420) (1.416,0.460) (1.414,0.500) (1.411,0.540) (1.407,0.580) (1.402,0.620) (1.395,0.660) (1.388,0.700) (1.380,0.736) (1.370,0.780) (1.360,0.819) (1.349,0.860) (1.340,0.891) (1.331,0.920) (1.320,0.957) (1.307,1.000) (1.294,1.040) (1.280,1.080) (1.267,1.120) (1.254,1.160) (1.241,1.200) (1.227,1.240) (1.214,1.280) (1.201,1.320) (1.189,1.360) (1.176,1.400) (1.164,1.440) (1.153,1.480) (1.142,1.520) (1.131,1.560) (1.120,1.600) (1.110,1.640) (1.100,1.680) (1.091,1.720) (1.082,1.760) (1.073,1.800) (1.065,1.840) (1.057,1.880) (1.049,1.920) (1.046,1.940)}
\EQ{(0.624,-1.940) (0.627,-1.900) (0.630,-1.860) (0.633,-1.820) (0.636,-1.780) (0.639,-1.740) (0.643,-1.700) (0.647,-1.660) (0.651,-1.620) (0.655,-1.580) (0.660,-1.540) (0.665,-1.500) (0.670,-1.460) (0.676,-1.420) (0.680,-1.390) (0.685,-1.360) (0.691,-1.320) (0.698,-1.280) (0.705,-1.240) (0.713,-1.200) (0.720,-1.166) (0.730,-1.120) (0.739,-1.080) (0.749,-1.040) (0.759,-1.000) (0.769,-0.960) (0.780,-0.922) (0.792,-0.880) (0.804,-0.840) (0.817,-0.800) (0.830,-0.760) (0.840,-0.731) (0.851,-0.700) (0.865,-0.660) (0.880,-0.620) (0.895,-0.580) (0.911,-0.540) (0.927,-0.500) (0.940,-0.469) (0.952,-0.440) (0.969,-0.400) (0.986,-0.360) (1.000,-0.329) (1.013,-0.300) (1.030,-0.260) (1.048,-0.220) (1.066,-0.180) (1.080,-0.148) (1.092,-0.120) (1.109,-0.080) (1.126,-0.040) (1.140,-0.006) (1.158,0.040) (1.172,0.080) (1.186,0.120) (1.198,0.160) (1.210,0.200) (1.219,0.240) (1.227,0.280) (1.233,0.320) (1.238,0.360) (1.240,0.400) (1.240,0.440) (1.238,0.480) (1.234,0.520) (1.227,0.560) (1.220,0.594) (1.208,0.640) (1.195,0.680) (1.181,0.720) (1.165,0.760) (1.147,0.800) (1.129,0.840) (1.109,0.880) (1.089,0.920) (1.068,0.960) (1.047,1.000) (1.026,1.040) (1.005,1.080) (0.984,1.120) (0.964,1.160) (0.944,1.200) (0.925,1.240) (0.907,1.280) (0.890,1.320) (0.873,1.360) (0.860,1.394) (0.843,1.440) (0.828,1.480) (0.815,1.520) (0.803,1.560) (0.791,1.600) (0.780,1.639) (0.769,1.680) (0.760,1.717) (0.750,1.760) (0.741,1.800) (0.733,1.840) (0.725,1.880) (0.718,1.920) (0.715,1.940)}
\EQ{(0.292,-1.940) (0.293,-1.900) (0.295,-1.860) (0.296,-1.820) (0.298,-1.780) (0.300,-1.740) (0.302,-1.700) (0.304,-1.660) (0.306,-1.620) (0.309,-1.580) (0.312,-1.540) (0.315,-1.500) (0.318,-1.460) (0.322,-1.420) (0.326,-1.380) (0.330,-1.340) (0.335,-1.300) (0.340,-1.263) (0.346,-1.220) (0.353,-1.180) (0.360,-1.141) (0.368,-1.100) (0.377,-1.060) (0.387,-1.020) (0.397,-0.980) (0.409,-0.940) (0.420,-0.903) (0.434,-0.860) (0.448,-0.820) (0.460,-0.789) (0.471,-0.760) (0.487,-0.720) (0.500,-0.689) (0.512,-0.660) (0.530,-0.620) (0.549,-0.580) (0.568,-0.540) (0.588,-0.500) (0.609,-0.460) (0.631,-0.420) (0.653,-0.380) (0.677,-0.340) (0.700,-0.302) (0.720,-0.270) (0.739,-0.240) (0.760,-0.209) (0.779,-0.180) (0.800,-0.150) (0.820,-0.121) (0.840,-0.092) (0.860,-0.063) (0.880,-0.035) (0.900,-0.006) (0.918,0.020) (0.940,0.052) (0.958,0.080) (0.980,0.114) (0.996,0.140) (1.019,0.180) (1.039,0.220) (1.058,0.260) (1.073,0.300) (1.085,0.340) (1.094,0.380) (1.098,0.420) (1.098,0.460) (1.092,0.500) (1.081,0.540) (1.065,0.580) (1.044,0.620) (1.020,0.659) (1.000,0.686) (0.980,0.712) (0.960,0.736) (0.940,0.759) (0.902,0.800) (0.880,0.823) (0.844,0.860) (0.820,0.885) (0.786,0.920) (0.760,0.947) (0.729,0.980) (0.700,1.013) (0.680,1.037) (0.661,1.060) (0.640,1.087) (0.620,1.114) (0.602,1.140) (0.580,1.173) (0.564,1.200) (0.541,1.240) (0.521,1.280) (0.502,1.320) (0.486,1.360) (0.471,1.400) (0.460,1.432) (0.445,1.480) (0.434,1.520) (0.423,1.560) (0.414,1.600) (0.405,1.640) (0.398,1.680) (0.390,1.720) (0.384,1.760) (0.378,1.800) (0.372,1.840) (0.367,1.880) (0.362,1.920) (0.359,1.940)}
\EQ{(-0.383,-1.940) (-0.385,-1.900) (-0.388,-1.860) (-0.390,-1.820) (-0.393,-1.780) (-0.396,-1.740) (-0.399,-1.700) (-0.402,-1.660) (-0.406,-1.620) (-0.409,-1.580) (-0.413,-1.540) (-0.418,-1.500) (-0.422,-1.460) (-0.427,-1.420) (-0.433,-1.380) (-0.439,-1.340) (-0.445,-1.300) (-0.453,-1.260) (-0.460,-1.222) (-0.469,-1.180) (-0.478,-1.140) (-0.488,-1.100) (-0.499,-1.060) (-0.511,-1.020) (-0.524,-0.980) (-0.538,-0.940) (-0.553,-0.900) (-0.568,-0.860) (-0.580,-0.831) (-0.593,-0.800) (-0.610,-0.760) (-0.628,-0.720) (-0.640,-0.692) (-0.655,-0.660) (-0.673,-0.620) (-0.692,-0.580) (-0.711,-0.540) (-0.731,-0.500) (-0.752,-0.460) (-0.773,-0.420) (-0.795,-0.380) (-0.818,-0.340) (-0.840,-0.301) (-0.860,-0.268) (-0.877,-0.240) (-0.900,-0.203) (-0.920,-0.172) (-0.940,-0.141) (-0.960,-0.110) (-0.979,-0.080) (-1.000,-0.047) (-1.016,-0.020) (-1.040,0.020) (-1.060,0.056) (-1.080,0.095) (-1.100,0.139) (-1.115,0.180) (-1.127,0.220) (-1.135,0.260) (-1.139,0.300) (-1.139,0.340) (-1.135,0.380) (-1.126,0.420) (-1.113,0.460) (-1.100,0.492) (-1.086,0.520) (-1.063,0.560) (-1.040,0.596) (-1.020,0.624) (-1.000,0.651) (-0.980,0.675) (-0.960,0.699) (-0.940,0.722) (-0.920,0.745) (-0.888,0.780) (-0.860,0.811) (-0.833,0.840) (-0.800,0.876) (-0.779,0.900) (-0.744,0.940) (-0.720,0.968) (-0.700,0.993) (-0.680,1.019) (-0.660,1.046) (-0.640,1.074) (-0.622,1.100) (-0.600,1.136) (-0.580,1.171) (-0.564,1.200) (-0.545,1.240) (-0.527,1.280) (-0.511,1.320) (-0.500,1.349) (-0.489,1.380) (-0.480,1.410) (-0.471,1.440) (-0.460,1.480) (-0.451,1.520) (-0.442,1.560) (-0.434,1.600) (-0.427,1.640) (-0.420,1.680) (-0.414,1.720) (-0.408,1.760) (-0.403,1.800) (-0.398,1.840) (-0.394,1.880) (-0.390,1.920) (-0.388,1.940)}
\EQ{(-0.713,-1.940) (-0.717,-1.900) (-0.720,-1.862) (-0.724,-1.820) (-0.728,-1.780) (-0.732,-1.740) (-0.737,-1.700) (-0.741,-1.660) (-0.746,-1.620) (-0.752,-1.580) (-0.757,-1.540) (-0.763,-1.500) (-0.769,-1.460) (-0.776,-1.420) (-0.783,-1.380) (-0.790,-1.340) (-0.798,-1.300) (-0.806,-1.260) (-0.815,-1.220) (-0.824,-1.180) (-0.833,-1.140) (-0.843,-1.100) (-0.854,-1.060) (-0.864,-1.020) (-0.876,-0.980) (-0.887,-0.940) (-0.900,-0.900) (-0.912,-0.860) (-0.925,-0.820) (-0.938,-0.780) (-0.952,-0.740) (-0.966,-0.700) (-0.980,-0.662) (-0.995,-0.620) (-1.010,-0.580) (-1.026,-0.540) (-1.040,-0.503) (-1.057,-0.460) (-1.073,-0.420) (-1.089,-0.380) (-1.105,-0.340) (-1.120,-0.304) (-1.138,-0.260) (-1.154,-0.220) (-1.170,-0.180) (-1.185,-0.140) (-1.200,-0.100) (-1.214,-0.060) (-1.227,-0.020) (-1.240,0.020) (-1.251,0.060) (-1.260,0.096) (-1.269,0.140) (-1.276,0.180) (-1.282,0.220) (-1.285,0.260) (-1.286,0.300) (-1.286,0.340) (-1.283,0.380) (-1.280,0.411) (-1.272,0.460) (-1.264,0.500) (-1.254,0.540) (-1.241,0.580) (-1.228,0.620) (-1.213,0.660) (-1.200,0.691) (-1.187,0.720) (-1.169,0.760) (-1.150,0.800) (-1.131,0.840) (-1.110,0.880) (-1.090,0.920) (-1.069,0.960) (-1.049,1.000) (-1.029,1.040) (-1.009,1.080) (-0.989,1.120) (-0.970,1.160) (-0.952,1.200) (-0.940,1.228) (-0.926,1.260) (-0.910,1.300) (-0.895,1.340) (-0.880,1.380) (-0.866,1.420) (-0.853,1.460) (-0.840,1.500) (-0.829,1.540) (-0.820,1.572) (-0.807,1.620) (-0.800,1.651) (-0.793,1.680) (-0.784,1.720) (-0.776,1.760) (-0.768,1.800) (-0.761,1.840) (-0.754,1.880) (-0.748,1.920) (-0.745,1.940)}
\EQ{(-1.034,-1.940) (-1.038,-1.900) (-1.042,-1.860) (-1.046,-1.820) (-1.050,-1.780) (-1.055,-1.740) (-1.060,-1.700) (-1.065,-1.660) (-1.070,-1.620) (-1.076,-1.580) (-1.082,-1.540) (-1.088,-1.500) (-1.094,-1.460) (-1.100,-1.422) (-1.107,-1.380) (-1.114,-1.340) (-1.120,-1.308) (-1.129,-1.260) (-1.137,-1.220) (-1.145,-1.180) (-1.153,-1.140) (-1.160,-1.108) (-1.171,-1.060) (-1.180,-1.020) (-1.189,-0.980) (-1.199,-0.940) (-1.209,-0.900) (-1.219,-0.860) (-1.229,-0.820) (-1.239,-0.780) (-1.250,-0.740) (-1.260,-0.703) (-1.272,-0.660) (-1.280,-0.630) (-1.288,-0.600) (-1.299,-0.560) (-1.311,-0.520) (-1.320,-0.487) (-1.333,-0.440) (-1.344,-0.400) (-1.355,-0.360) (-1.366,-0.320) (-1.376,-0.280) (-1.387,-0.240) (-1.397,-0.200) (-1.406,-0.160) (-1.415,-0.120) (-1.424,-0.080) (-1.432,-0.040) (-1.439,-0.000) (-1.446,0.040) (-1.452,0.080) (-1.457,0.120) (-1.460,0.153) (-1.464,0.200) (-1.466,0.240) (-1.467,0.280) (-1.467,0.320) (-1.466,0.360) (-1.463,0.400) (-1.460,0.440) (-1.456,0.480) (-1.450,0.520) (-1.444,0.560) (-1.437,0.600) (-1.428,0.640) (-1.420,0.677) (-1.409,0.720) (-1.400,0.756) (-1.388,0.800) (-1.376,0.840) (-1.364,0.880) (-1.351,0.920) (-1.340,0.956) (-1.326,1.000) (-1.313,1.040) (-1.300,1.078) (-1.286,1.120) (-1.273,1.160) (-1.260,1.200) (-1.248,1.240) (-1.235,1.280) (-1.223,1.320) (-1.211,1.360) (-1.200,1.397) (-1.188,1.440) (-1.180,1.469) (-1.172,1.500) (-1.161,1.540) (-1.151,1.580) (-1.142,1.620) (-1.133,1.660) (-1.124,1.700) (-1.115,1.740) (-1.107,1.780) (-1.100,1.818) (-1.092,1.860) (-1.085,1.900) (-1.080,1.930) (-1.078,1.940)}
\EQ{(-1.345,-1.940) (-1.349,-1.900) (-1.353,-1.860) (-1.357,-1.820) (-1.361,-1.780) (-1.366,-1.740) (-1.371,-1.700) (-1.375,-1.660) (-1.380,-1.624) (-1.386,-1.580) (-1.391,-1.540) (-1.396,-1.500) (-1.402,-1.460) (-1.408,-1.420) (-1.414,-1.380) (-1.420,-1.340) (-1.426,-1.300) (-1.433,-1.260) (-1.439,-1.220) (-1.446,-1.180) (-1.453,-1.140) (-1.460,-1.101) (-1.467,-1.060) (-1.475,-1.020) (-1.482,-0.980) (-1.490,-0.940) (-1.498,-0.900) (-1.505,-0.860) (-1.513,-0.820) (-1.520,-0.786) (-1.529,-0.740) (-1.537,-0.700) (-1.545,-0.660) (-1.553,-0.620) (-1.560,-0.586) (-1.569,-0.540) (-1.577,-0.500) (-1.585,-0.460) (-1.593,-0.420) (-1.600,-0.380) (-1.607,-0.340) (-1.614,-0.300) (-1.620,-0.268) (-1.628,-0.220) (-1.634,-0.180) (-1.640,-0.140) (-1.645,-0.100) (-1.651,-0.060) (-1.655,-0.020) (-1.659,0.020) (-1.663,0.060) (-1.666,0.100) (-1.669,0.140) (-1.671,0.180) (-1.672,0.220) (-1.673,0.260) (-1.673,0.300) (-1.672,0.340) (-1.671,0.380) (-1.669,0.420) (-1.667,0.460) (-1.664,0.500) (-1.660,0.540) (-1.656,0.580) (-1.651,0.620) (-1.646,0.660) (-1.640,0.698) (-1.633,0.740) (-1.627,0.780) (-1.620,0.817) (-1.612,0.860) (-1.604,0.900) (-1.596,0.940) (-1.587,0.980) (-1.580,1.014) (-1.570,1.060) (-1.561,1.100) (-1.552,1.140) (-1.543,1.180) (-1.534,1.220) (-1.525,1.260) (-1.516,1.300) (-1.507,1.340) (-1.500,1.372) (-1.489,1.420) (-1.481,1.460) (-1.472,1.500) (-1.464,1.540) (-1.456,1.580) (-1.448,1.620) (-1.440,1.660) (-1.433,1.700) (-1.426,1.740) (-1.420,1.772) (-1.412,1.820) (-1.405,1.860) (-1.400,1.892) (-1.392,1.940)}
\EQ{(-1.648,-1.940) (-1.652,-1.900) (-1.655,-1.860) (-1.659,-1.820) (-1.663,-1.780) (-1.667,-1.740) (-1.672,-1.700) (-1.676,-1.660) (-1.680,-1.624) (-1.685,-1.580) (-1.690,-1.540) (-1.694,-1.500) (-1.699,-1.460) (-1.704,-1.420) (-1.709,-1.380) (-1.714,-1.340) (-1.720,-1.300) (-1.725,-1.260) (-1.730,-1.220) (-1.736,-1.180) (-1.742,-1.140) (-1.747,-1.100) (-1.753,-1.060) (-1.759,-1.020) (-1.765,-0.980) (-1.770,-0.940) (-1.776,-0.900) (-1.782,-0.860) (-1.788,-0.820) (-1.794,-0.780) (-1.800,-0.740) (-1.806,-0.700) (-1.812,-0.660) (-1.818,-0.620) (-1.823,-0.580) (-1.829,-0.540) (-1.835,-0.500) (-1.840,-0.461) (-1.845,-0.420) (-1.851,-0.380) (-1.856,-0.340) (-1.860,-0.303) (-1.865,-0.260) (-1.869,-0.220) (-1.873,-0.180) (-1.877,-0.140) (-1.880,-0.109) (-1.884,-0.060) (-1.887,-0.020) (-1.890,0.020) (-1.892,0.060) (-1.894,0.100) (-1.895,0.140) (-1.896,0.180) (-1.897,0.220) (-1.898,0.260) (-1.898,0.300) (-1.897,0.340) (-1.896,0.380) (-1.895,0.420) (-1.893,0.460) (-1.891,0.500) (-1.889,0.540) (-1.886,0.580) (-1.883,0.620) (-1.880,0.652) (-1.875,0.700) (-1.871,0.740) (-1.867,0.780) (-1.862,0.820) (-1.857,0.860) (-1.852,0.900) (-1.846,0.940) (-1.840,0.980) (-1.835,1.020) (-1.829,1.060) (-1.822,1.100) (-1.816,1.140) (-1.810,1.180) (-1.803,1.220) (-1.797,1.260) (-1.790,1.300) (-1.784,1.340) (-1.777,1.380) (-1.771,1.420) (-1.764,1.460) (-1.758,1.500) (-1.752,1.540) (-1.745,1.580) (-1.740,1.614) (-1.733,1.660) (-1.727,1.700) (-1.721,1.740) (-1.715,1.780) (-1.709,1.820) (-1.704,1.860) (-1.698,1.900) (-1.693,1.940)}
\EQ{(-1.945,-1.940) (-1.949,-1.900) (-1.952,-1.860) (-1.955,-1.820) (-1.959,-1.780) (-1.962,-1.740) (-1.966,-1.700) (-1.970,-1.660) (-1.974,-1.620) (-1.977,-1.580) (-1.981,-1.540) (-1.985,-1.500) (-1.989,-1.460) (-1.993,-1.420) (-1.998,-1.380) (-2.002,-1.340) (-2.006,-1.300) (-2.010,-1.260) (-2.015,-1.220) (-2.019,-1.180) (-2.024,-1.140) (-2.028,-1.100) (-2.033,-1.060) (-2.037,-1.020) (-2.042,-0.980) (-2.046,-0.940) (-2.051,-0.900) (-2.055,-0.860) (-2.059,-0.820) (-2.064,-0.780) (-2.068,-0.740) (-2.073,-0.700) (-2.077,-0.660) (-2.081,-0.620) (-2.085,-0.580) (-2.089,-0.540) (-2.093,-0.500) (-2.097,-0.460) (-2.101,-0.420) (-2.105,-0.380) (-2.108,-0.340) (-2.111,-0.300) (-2.115,-0.260) (-2.118,-0.220) (-2.120,-0.185) (-2.123,-0.140) (-2.125,-0.100) (-2.127,-0.060) (-2.129,-0.020) (-2.131,0.020) (-2.133,0.060) (-2.134,0.100) (-2.135,0.140) (-2.135,0.180) (-2.136,0.220) (-2.136,0.260) (-2.136,0.300) (-2.136,0.340) (-2.135,0.380) (-2.134,0.420) (-2.133,0.460) (-2.132,0.500) (-2.130,0.540) (-2.128,0.580) (-2.126,0.620) (-2.123,0.660) (-2.121,0.700) (-2.118,0.740) (-2.115,0.780) (-2.111,0.820) (-2.108,0.860) (-2.104,0.900) (-2.100,0.940) (-2.096,0.980) (-2.092,1.020) (-2.088,1.060) (-2.084,1.100) (-2.080,1.134) (-2.075,1.180) (-2.070,1.220) (-2.065,1.260) (-2.061,1.300) (-2.056,1.340) (-2.051,1.380) (-2.046,1.420) (-2.041,1.460) (-2.036,1.500) (-2.032,1.540) (-2.027,1.580) (-2.022,1.620) (-2.017,1.660) (-2.012,1.700) (-2.008,1.740) (-2.003,1.780) (-1.998,1.820) (-1.994,1.860) (-1.989,1.900) (-1.985,1.940)}
\EQ{(-2.238,-1.940) (-2.241,-1.900) (-2.244,-1.860) (-2.247,-1.820) (-2.250,-1.780) (-2.253,-1.740) (-2.256,-1.700) (-2.259,-1.660) (-2.262,-1.620) (-2.265,-1.580) (-2.269,-1.540) (-2.272,-1.500) (-2.275,-1.460) (-2.279,-1.420) (-2.282,-1.380) (-2.285,-1.340) (-2.289,-1.300) (-2.292,-1.260) (-2.296,-1.220) (-2.299,-1.180) (-2.303,-1.140) (-2.306,-1.100) (-2.310,-1.060) (-2.313,-1.020) (-2.316,-0.980) (-2.320,-0.940) (-2.323,-0.900) (-2.327,-0.860) (-2.330,-0.820) (-2.333,-0.780) (-2.337,-0.740) (-2.340,-0.700) (-2.340,-0.699)}
\EQ{(-2.340,1.188) (-2.337,1.220) (-2.334,1.260) (-2.330,1.300) (-2.327,1.340) (-2.323,1.380) (-2.320,1.415) (-2.316,1.460) (-2.312,1.500) (-2.308,1.540) (-2.305,1.580) (-2.301,1.620) (-2.297,1.660) (-2.294,1.700) (-2.290,1.740) (-2.286,1.780) (-2.283,1.820) (-2.279,1.860) (-2.275,1.900) (-2.272,1.940)}

      % Conductor neutro en campo uniforme - líneas
% Generado por scripts/generar_datos.py (no editar a mano).
\LC{(-2.300,-1.850) (-2.264,-1.847) (-2.228,-1.845) (-2.192,-1.842) (-2.156,-1.839) (-2.121,-1.837) (-2.085,-1.834) (-2.049,-1.831) (-2.013,-1.828) (-1.977,-1.825) (-1.941,-1.822) (-1.905,-1.818) (-1.869,-1.815) (-1.834,-1.812) (-1.798,-1.809) (-1.762,-1.805) (-1.726,-1.802) (-1.690,-1.798) (-1.654,-1.795) (-1.619,-1.791) (-1.583,-1.787) (-1.547,-1.784) (-1.511,-1.780) (-1.475,-1.776) (-1.440,-1.772) (-1.404,-1.768) (-1.368,-1.764) (-1.332,-1.760) (-1.296,-1.756) (-1.261,-1.752) (-1.225,-1.747) (-1.189,-1.743) (-1.153,-1.739) (-1.118,-1.735) (-1.082,-1.730) (-1.046,-1.726) (-1.011,-1.722) (-0.975,-1.717) (-0.939,-1.713) (-0.903,-1.709) (-0.868,-1.704) (-0.832,-1.700) (-0.796,-1.696) (-0.760,-1.692) (-0.725,-1.687) (-0.689,-1.683) (-0.653,-1.679) (-0.617,-1.675) (-0.581,-1.671) (-0.546,-1.668) (-0.510,-1.664) (-0.474,-1.661) (-0.438,-1.657) (-0.402,-1.654) (-0.366,-1.651) (-0.331,-1.648) (-0.295,-1.646) (-0.259,-1.644) (-0.223,-1.642) (-0.187,-1.640) (-0.151,-1.638) (-0.115,-1.637) (-0.079,-1.636) (-0.043,-1.635) (-0.007,-1.635) (0.029,-1.635) (0.065,-1.635) (0.101,-1.636) (0.137,-1.637) (0.173,-1.638) (0.209,-1.639) (0.245,-1.641) (0.281,-1.643) (0.317,-1.645) (0.353,-1.647) (0.389,-1.650) (0.425,-1.652) (0.461,-1.655) (0.496,-1.658) (0.532,-1.661) (0.568,-1.664) (0.604,-1.668) (0.640,-1.671) (0.676,-1.675) (0.711,-1.678) (0.747,-1.682) (0.783,-1.686) (0.819,-1.689) (0.855,-1.693) (0.891,-1.697) (0.926,-1.701) (0.962,-1.705) (0.998,-1.709) (1.034,-1.713) (1.069,-1.717) (1.105,-1.720) (1.141,-1.724) (1.177,-1.728) (1.213,-1.732) (1.248,-1.736) (1.284,-1.740) (1.320,-1.743) (1.356,-1.747) (1.392,-1.751) (1.427,-1.755) (1.463,-1.758) (1.499,-1.762) (1.535,-1.765) (1.571,-1.769) (1.607,-1.772) (1.642,-1.776) (1.678,-1.779) (1.714,-1.782) (1.750,-1.786) (1.786,-1.789) (1.822,-1.792) (1.858,-1.795) (1.893,-1.798) (1.929,-1.801) (1.965,-1.804) (2.001,-1.807) (2.037,-1.810) (2.073,-1.812) (2.109,-1.815) (2.145,-1.818) (2.181,-1.820) (2.216,-1.823) (2.252,-1.825) (2.288,-1.828) (2.300,-1.829)}
\LC{(-2.300,-1.645) (-2.264,-1.642) (-2.228,-1.639) (-2.192,-1.636) (-2.156,-1.633) (-2.121,-1.630) (-2.085,-1.627) (-2.049,-1.624) (-2.013,-1.621) (-1.977,-1.617) (-1.941,-1.614) (-1.906,-1.610) (-1.870,-1.607) (-1.834,-1.603) (-1.798,-1.599) (-1.762,-1.595) (-1.727,-1.591) (-1.691,-1.587) (-1.655,-1.583) (-1.619,-1.579) (-1.583,-1.575) (-1.548,-1.570) (-1.512,-1.566) (-1.476,-1.561) (-1.441,-1.557) (-1.405,-1.552) (-1.369,-1.547) (-1.334,-1.542) (-1.298,-1.538) (-1.262,-1.532) (-1.227,-1.527) (-1.191,-1.522) (-1.155,-1.517) (-1.120,-1.512) (-1.084,-1.506) (-1.049,-1.501) (-1.013,-1.495) (-0.977,-1.490) (-0.942,-1.484) (-0.906,-1.479) (-0.871,-1.473) (-0.835,-1.467) (-0.800,-1.462) (-0.764,-1.456) (-0.729,-1.450) (-0.693,-1.445) (-0.657,-1.439) (-0.622,-1.433) (-0.586,-1.428) (-0.551,-1.422) (-0.515,-1.417) (-0.480,-1.412) (-0.444,-1.407) (-0.408,-1.402) (-0.373,-1.397) (-0.337,-1.393) (-0.301,-1.389) (-0.265,-1.385) (-0.229,-1.382) (-0.193,-1.379) (-0.158,-1.377) (-0.122,-1.375) (-0.086,-1.373) (-0.050,-1.372) (-0.014,-1.372) (0.022,-1.372) (0.058,-1.372) (0.094,-1.374) (0.130,-1.375) (0.166,-1.377) (0.202,-1.380) (0.238,-1.382) (0.274,-1.386) (0.310,-1.389) (0.346,-1.393) (0.381,-1.397) (0.417,-1.401) (0.453,-1.405) (0.489,-1.410) (0.524,-1.414) (0.560,-1.419) (0.596,-1.424) (0.631,-1.429) (0.667,-1.434) (0.703,-1.439) (0.738,-1.444) (0.774,-1.449) (0.810,-1.454) (0.845,-1.459) (0.881,-1.464) (0.916,-1.469) (0.952,-1.474) (0.988,-1.479) (1.023,-1.484) (1.059,-1.489) (1.095,-1.494) (1.130,-1.499) (1.166,-1.504) (1.202,-1.509) (1.237,-1.513) (1.273,-1.518) (1.309,-1.523) (1.344,-1.527) (1.380,-1.532) (1.416,-1.536) (1.452,-1.541) (1.487,-1.545) (1.523,-1.549) (1.559,-1.553) (1.595,-1.557) (1.630,-1.561) (1.666,-1.565) (1.702,-1.569) (1.738,-1.573) (1.774,-1.576) (1.809,-1.580) (1.845,-1.584) (1.881,-1.587) (1.917,-1.590) (1.953,-1.594) (1.989,-1.597) (2.025,-1.600) (2.060,-1.603) (2.096,-1.606) (2.132,-1.609) (2.168,-1.612) (2.204,-1.615) (2.240,-1.618) (2.276,-1.621) (2.300,-1.623)}
\LC{(-2.300,-1.440) (-2.264,-1.437) (-2.228,-1.434) (-2.192,-1.431) (-2.157,-1.428) (-2.121,-1.424) (-2.085,-1.421) (-2.049,-1.417) (-2.013,-1.414) (-1.977,-1.410) (-1.942,-1.406) (-1.906,-1.402) (-1.870,-1.398) (-1.834,-1.394) (-1.799,-1.390) (-1.763,-1.386) (-1.727,-1.381) (-1.691,-1.377) (-1.656,-1.372) (-1.620,-1.367) (-1.584,-1.362) (-1.549,-1.357) (-1.513,-1.352) (-1.477,-1.347) (-1.442,-1.341) (-1.406,-1.336) (-1.371,-1.330) (-1.335,-1.324) (-1.300,-1.318) (-1.264,-1.312) (-1.229,-1.306) (-1.193,-1.300) (-1.158,-1.294) (-1.122,-1.287) (-1.087,-1.280) (-1.052,-1.273) (-1.016,-1.267) (-0.981,-1.259) (-0.946,-1.252) (-0.911,-1.245) (-0.875,-1.237) (-0.840,-1.230) (-0.805,-1.222) (-0.770,-1.214) (-0.735,-1.207) (-0.700,-1.199) (-0.665,-1.191) (-0.629,-1.182) (-0.594,-1.174) (-0.559,-1.166) (-0.524,-1.157) (-0.489,-1.149) (-0.454,-1.141) (-0.419,-1.132) (-0.384,-1.124) (-0.349,-1.116) (-0.314,-1.108) (-0.279,-1.100) (-0.244,-1.093) (-0.208,-1.086) (-0.173,-1.080) (-0.137,-1.076) (-0.101,-1.072) (-0.065,-1.070) (-0.029,-1.069) (0.007,-1.070) (0.043,-1.072) (0.078,-1.076) (0.114,-1.080) (0.150,-1.085) (0.185,-1.091) (0.221,-1.097) (0.256,-1.103) (0.292,-1.110) (0.327,-1.116) (0.362,-1.123) (0.398,-1.130) (0.433,-1.137) (0.468,-1.145) (0.503,-1.152) (0.539,-1.159) (0.574,-1.166) (0.609,-1.174) (0.645,-1.181) (0.680,-1.188) (0.715,-1.195) (0.750,-1.202) (0.786,-1.209) (0.821,-1.216) (0.856,-1.223) (0.892,-1.230) (0.927,-1.236) (0.963,-1.243) (0.998,-1.249) (1.033,-1.256) (1.069,-1.262) (1.104,-1.268) (1.140,-1.274) (1.175,-1.280) (1.211,-1.286) (1.246,-1.292) (1.282,-1.298) (1.317,-1.303) (1.353,-1.309) (1.389,-1.314) (1.424,-1.319) (1.460,-1.324) (1.495,-1.330) (1.531,-1.334) (1.567,-1.339) (1.603,-1.344) (1.638,-1.348) (1.674,-1.353) (1.710,-1.357) (1.745,-1.362) (1.781,-1.366) (1.817,-1.370) (1.853,-1.374) (1.889,-1.378) (1.924,-1.381) (1.960,-1.385) (1.996,-1.389) (2.032,-1.392) (2.068,-1.396) (2.103,-1.399) (2.139,-1.402) (2.175,-1.405) (2.211,-1.408) (2.247,-1.411) (2.283,-1.414) (2.300,-1.416)}
\LC{(-2.300,-1.235) (-2.264,-1.232) (-2.228,-1.229) (-2.192,-1.225) (-2.157,-1.222) (-2.121,-1.218) (-2.085,-1.215) (-2.049,-1.211) (-2.013,-1.207) (-1.978,-1.203) (-1.942,-1.199) (-1.906,-1.195) (-1.870,-1.190) (-1.835,-1.186) (-1.799,-1.181) (-1.763,-1.176) (-1.728,-1.172) (-1.692,-1.166) (-1.656,-1.161) (-1.621,-1.156) (-1.585,-1.150) (-1.550,-1.145) (-1.514,-1.139) (-1.479,-1.133) (-1.443,-1.126) (-1.408,-1.120) (-1.372,-1.113) (-1.337,-1.107) (-1.302,-1.100) (-1.266,-1.093) (-1.231,-1.085) (-1.196,-1.078) (-1.161,-1.070) (-1.126,-1.062) (-1.091,-1.054) (-1.056,-1.045) (-1.021,-1.037) (-0.986,-1.028) (-0.951,-1.019) (-0.916,-1.009) (-0.882,-1.000) (-0.847,-0.990) (-0.812,-0.980) (-0.778,-0.969) (-0.744,-0.959) (-0.709,-0.948) (-0.675,-0.936) (-0.641,-0.925) (-0.607,-0.913) (-0.573,-0.900) (-0.540,-0.887) (-0.506,-0.874) (-0.473,-0.860) (-0.440,-0.846) (-0.407,-0.830) (-0.375,-0.814) (-0.344,-0.796) (-0.313,-0.778) (-0.284,-0.757) (-0.256,-0.734) (-0.229,-0.710)}
\LC{(-2.300,-1.030) (-2.264,-1.027) (-2.228,-1.023) (-2.192,-1.020) (-2.157,-1.017) (-2.121,-1.013) (-2.085,-1.009) (-2.049,-1.005) (-2.013,-1.001) (-1.978,-0.997) (-1.942,-0.993) (-1.906,-0.988) (-1.871,-0.983) (-1.835,-0.979) (-1.799,-0.974) (-1.764,-0.968) (-1.728,-0.963) (-1.692,-0.958) (-1.657,-0.952) (-1.621,-0.946) (-1.586,-0.940) (-1.550,-0.934) (-1.515,-0.927) (-1.480,-0.920) (-1.444,-0.913) (-1.409,-0.906) (-1.374,-0.899) (-1.339,-0.891) (-1.304,-0.883) (-1.269,-0.874) (-1.234,-0.866) (-1.199,-0.857) (-1.164,-0.848) (-1.129,-0.838) (-1.095,-0.829) (-1.060,-0.818) (-1.026,-0.808) (-0.991,-0.797) (-0.957,-0.786) (-0.923,-0.774) (-0.889,-0.763) (-0.855,-0.750) (-0.821,-0.738) (-0.788,-0.725) (-0.754,-0.711) (-0.721,-0.697) (-0.688,-0.683) (-0.655,-0.668) (-0.623,-0.653) (-0.590,-0.638) (-0.558,-0.622) (-0.526,-0.605) (-0.494,-0.589) (-0.462,-0.572) (-0.430,-0.555) (-0.399,-0.538) (-0.367,-0.520)}
\LC{(-2.300,-0.825) (-2.264,-0.822) (-2.228,-0.819) (-2.192,-0.815) (-2.157,-0.812) (-2.121,-0.808) (-2.085,-0.804) (-2.049,-0.800) (-2.013,-0.796) (-1.978,-0.792) (-1.942,-0.787) (-1.906,-0.783) (-1.871,-0.778) (-1.835,-0.773) (-1.799,-0.768) (-1.764,-0.763) (-1.728,-0.757) (-1.693,-0.751) (-1.657,-0.745) (-1.622,-0.739) (-1.586,-0.732) (-1.551,-0.725) (-1.516,-0.718) (-1.480,-0.711) (-1.445,-0.703) (-1.410,-0.695) (-1.375,-0.687) (-1.340,-0.678) (-1.305,-0.669) (-1.271,-0.660) (-1.236,-0.650) (-1.201,-0.640) (-1.167,-0.630) (-1.133,-0.619) (-1.098,-0.607) (-1.064,-0.596) (-1.031,-0.583) (-0.997,-0.571) (-0.963,-0.557) (-0.930,-0.544) (-0.897,-0.529) (-0.864,-0.515) (-0.832,-0.499) (-0.799,-0.484) (-0.767,-0.468) (-0.735,-0.451) (-0.703,-0.434) (-0.672,-0.416) (-0.641,-0.398) (-0.610,-0.380) (-0.579,-0.362) (-0.548,-0.343) (-0.518,-0.323) (-0.489,-0.302)}
\LC{(-2.300,-0.620) (-2.264,-0.617) (-2.228,-0.614) (-2.192,-0.611) (-2.157,-0.607) (-2.121,-0.604) (-2.085,-0.600) (-2.049,-0.596) (-2.013,-0.593) (-1.978,-0.588) (-1.942,-0.584) (-1.906,-0.580) (-1.870,-0.575) (-1.835,-0.570) (-1.799,-0.565) (-1.764,-0.560) (-1.728,-0.554) (-1.692,-0.548) (-1.657,-0.542) (-1.622,-0.536) (-1.586,-0.529) (-1.551,-0.522) (-1.516,-0.515) (-1.480,-0.507) (-1.445,-0.499) (-1.410,-0.490) (-1.375,-0.482) (-1.341,-0.472) (-1.306,-0.463) (-1.271,-0.452) (-1.237,-0.442) (-1.203,-0.430) (-1.169,-0.419) (-1.135,-0.406) (-1.101,-0.394) (-1.068,-0.380) (-1.035,-0.366) (-1.002,-0.351) (-0.970,-0.335) (-0.938,-0.319) (-0.906,-0.302) (-0.875,-0.284) (-0.844,-0.266) (-0.813,-0.247) (-0.783,-0.226) (-0.754,-0.205) (-0.725,-0.184) (-0.697,-0.161) (-0.670,-0.138) (-0.643,-0.114) (-0.635,-0.105)}
\LC{(-2.300,-0.415) (-2.264,-0.412) (-2.228,-0.410) (-2.192,-0.407) (-2.156,-0.404) (-2.121,-0.401) (-2.085,-0.398) (-2.049,-0.394) (-2.013,-0.391) (-1.977,-0.387) (-1.941,-0.383) (-1.906,-0.379) (-1.870,-0.375) (-1.834,-0.370) (-1.799,-0.365) (-1.763,-0.360) (-1.727,-0.355) (-1.692,-0.350) (-1.656,-0.344) (-1.621,-0.338) (-1.585,-0.332) (-1.550,-0.325) (-1.515,-0.318) (-1.479,-0.310) (-1.444,-0.302) (-1.409,-0.294) (-1.374,-0.285) (-1.340,-0.275) (-1.305,-0.265) (-1.271,-0.255) (-1.236,-0.244) (-1.202,-0.232) (-1.169,-0.219) (-1.135,-0.206) (-1.102,-0.191) (-1.070,-0.176) (-1.037,-0.160) (-1.006,-0.143) (-0.975,-0.125) (-0.944,-0.106) (-0.915,-0.085) (-0.886,-0.063) (-0.859,-0.040) (-0.833,-0.015) (-0.807,0.011) (-0.791,0.028)}
\LC{(-2.300,-0.210) (-2.264,-0.208) (-2.228,-0.206) (-2.192,-0.204) (-2.156,-0.201) (-2.120,-0.199) (-2.084,-0.196) (-2.049,-0.194) (-2.013,-0.191) (-1.977,-0.188) (-1.941,-0.185) (-1.905,-0.181) (-1.869,-0.178) (-1.833,-0.174) (-1.798,-0.170) (-1.762,-0.166) (-1.726,-0.162) (-1.690,-0.157) (-1.655,-0.152) (-1.619,-0.147) (-1.584,-0.142) (-1.548,-0.136) (-1.513,-0.130) (-1.477,-0.123) (-1.442,-0.116) (-1.407,-0.108) (-1.372,-0.100) (-1.337,-0.092) (-1.302,-0.082) (-1.267,-0.072) (-1.233,-0.061) (-1.199,-0.050) (-1.165,-0.037) (-1.132,-0.023) (-1.099,-0.009) (-1.067,0.007) (-1.035,0.025) (-1.005,0.044) (-0.975,0.064) (-0.947,0.087) (-0.920,0.111) (-0.912,0.119)}
\LC{(-2.300,-0.005) (-2.264,-0.004) (-2.228,-0.002) (-2.192,-0.001) (-2.156,0.000) (-2.120,0.002) (-2.084,0.003) (-2.048,0.005) (-2.012,0.007) (-1.976,0.009) (-1.940,0.011) (-1.904,0.013) (-1.868,0.015) (-1.833,0.018) (-1.797,0.020) (-1.761,0.023) (-1.725,0.026) (-1.689,0.029) (-1.653,0.032) (-1.617,0.036) (-1.582,0.039) (-1.546,0.043) (-1.510,0.048) (-1.474,0.053) (-1.439,0.058) (-1.403,0.063) (-1.368,0.069) (-1.332,0.076) (-1.297,0.083) (-1.262,0.091) (-1.227,0.100) (-1.192,0.109) (-1.158,0.120) (-1.124,0.132) (-1.090,0.145) (-1.058,0.160) (-1.026,0.177) (-0.997,0.198) (-0.988,0.206)}
\LC{(-2.300,0.200) (-2.264,0.200) (-2.228,0.201) (-2.192,0.201) (-2.156,0.201) (-2.120,0.202) (-2.084,0.202) (-2.048,0.203) (-2.012,0.203) (-1.976,0.204) (-1.940,0.204) (-1.904,0.205) (-1.868,0.205) (-1.832,0.206) (-1.796,0.207) (-1.760,0.208) (-1.724,0.209) (-1.688,0.210) (-1.652,0.211) (-1.616,0.212) (-1.580,0.213) (-1.544,0.215) (-1.508,0.216) (-1.472,0.218) (-1.436,0.220) (-1.400,0.222) (-1.364,0.224) (-1.329,0.227) (-1.293,0.230) (-1.257,0.234) (-1.221,0.238) (-1.185,0.242) (-1.150,0.248) (-1.114,0.254) (-1.079,0.262) (-1.044,0.271) (-1.021,0.276)}
\LC{(-2.300,0.405) (-2.264,0.404) (-2.228,0.404) (-2.192,0.403) (-2.156,0.402) (-2.120,0.401) (-2.084,0.400) (-2.048,0.400) (-2.012,0.399) (-1.976,0.398) (-1.940,0.397) (-1.904,0.396) (-1.868,0.394) (-1.832,0.393) (-1.796,0.392) (-1.760,0.391) (-1.724,0.390) (-1.688,0.388) (-1.652,0.387) (-1.616,0.385) (-1.580,0.384) (-1.544,0.382) (-1.508,0.380) (-1.472,0.378) (-1.437,0.376) (-1.401,0.374) (-1.365,0.372) (-1.329,0.370) (-1.293,0.367) (-1.257,0.365) (-1.221,0.362) (-1.185,0.359) (-1.149,0.356) (-1.113,0.352) (-1.078,0.349) (-1.042,0.342) (-1.031,0.337)}
\LC{(-2.300,0.610) (-2.264,0.608) (-2.228,0.607) (-2.192,0.605) (-2.156,0.603) (-2.120,0.601) (-2.084,0.599) (-2.048,0.597) (-2.012,0.595) (-1.976,0.592) (-1.941,0.590) (-1.905,0.587) (-1.869,0.585) (-1.833,0.582) (-1.797,0.579) (-1.761,0.575) (-1.725,0.572) (-1.689,0.568) (-1.654,0.565) (-1.618,0.561) (-1.582,0.556) (-1.546,0.552) (-1.511,0.547) (-1.475,0.542) (-1.440,0.536) (-1.404,0.530) (-1.369,0.524) (-1.333,0.517) (-1.298,0.510) (-1.263,0.502) (-1.228,0.493) (-1.193,0.484) (-1.159,0.473) (-1.125,0.462) (-1.091,0.449) (-1.058,0.435) (-1.027,0.417) (-1.017,0.410)}
\LC{(-2.300,0.815) (-2.264,0.813) (-2.228,0.810) (-2.192,0.808) (-2.156,0.805) (-2.120,0.802) (-2.085,0.799) (-2.049,0.796) (-2.013,0.792) (-1.977,0.789) (-1.941,0.785) (-1.905,0.782) (-1.870,0.778) (-1.834,0.773) (-1.798,0.769) (-1.762,0.764) (-1.727,0.759) (-1.691,0.754) (-1.656,0.749) (-1.620,0.743) (-1.585,0.737) (-1.549,0.730) (-1.514,0.723) (-1.479,0.716) (-1.444,0.708) (-1.409,0.699) (-1.374,0.690) (-1.339,0.681) (-1.305,0.670) (-1.270,0.659) (-1.236,0.647) (-1.203,0.635) (-1.169,0.621) (-1.137,0.606) (-1.105,0.590) (-1.073,0.572) (-1.043,0.553) (-1.013,0.532) (-0.985,0.509) (-0.959,0.485) (-0.950,0.477)}
\LC{(-2.300,1.020) (-2.264,1.017) (-2.228,1.014) (-2.192,1.011) (-2.157,1.007) (-2.121,1.004) (-2.085,1.000) (-2.049,0.996) (-2.013,0.992) (-1.978,0.988) (-1.942,0.984) (-1.906,0.979) (-1.870,0.975) (-1.835,0.969) (-1.799,0.964) (-1.764,0.959) (-1.728,0.953) (-1.693,0.947) (-1.657,0.940) (-1.622,0.933) (-1.587,0.926) (-1.551,0.919) (-1.516,0.911) (-1.481,0.902) (-1.446,0.893) (-1.412,0.884) (-1.377,0.874) (-1.343,0.863) (-1.309,0.852) (-1.275,0.840) (-1.241,0.827) (-1.208,0.813) (-1.175,0.799) (-1.142,0.783) (-1.110,0.767) (-1.079,0.749) (-1.048,0.730) (-1.018,0.710) (-0.990,0.688) (-0.962,0.665) (-0.936,0.641) (-0.911,0.614) (-0.889,0.586) (-0.869,0.556) (-0.856,0.536)}
\LC{(-2.300,1.225) (-2.264,1.222) (-2.228,1.218) (-2.192,1.215) (-2.157,1.211) (-2.121,1.207) (-2.085,1.203) (-2.049,1.199) (-2.014,1.195) (-1.978,1.190) (-1.942,1.185) (-1.907,1.180) (-1.871,1.175) (-1.835,1.170) (-1.800,1.164) (-1.764,1.158) (-1.729,1.152) (-1.693,1.146) (-1.658,1.139) (-1.623,1.132) (-1.587,1.125) (-1.552,1.117) (-1.517,1.109) (-1.482,1.100) (-1.447,1.091) (-1.413,1.082) (-1.378,1.072) (-1.344,1.062) (-1.309,1.051) (-1.275,1.039) (-1.241,1.027) (-1.208,1.014) (-1.174,1.000) (-1.141,0.986) (-1.109,0.971) (-1.076,0.955) (-1.045,0.938) (-1.013,0.921) (-0.983,0.902) (-0.952,0.882) (-0.923,0.861) (-0.895,0.839) (-0.868,0.815) (-0.841,0.790) (-0.817,0.764) (-0.793,0.737) (-0.772,0.708) (-0.752,0.678) (-0.735,0.647) (-0.721,0.613) (-0.718,0.602)}
\LC{(-2.300,1.430) (-2.264,1.427) (-2.228,1.423) (-2.193,1.419) (-2.157,1.416) (-2.121,1.412) (-2.085,1.408) (-2.049,1.403) (-2.014,1.399) (-1.978,1.394) (-1.942,1.390) (-1.907,1.385) (-1.871,1.379) (-1.835,1.374) (-1.800,1.369) (-1.764,1.363) (-1.729,1.357) (-1.693,1.350) (-1.658,1.344) (-1.623,1.337) (-1.587,1.330) (-1.552,1.323) (-1.517,1.315) (-1.482,1.307) (-1.447,1.299) (-1.412,1.290) (-1.377,1.281) (-1.342,1.272) (-1.308,1.262) (-1.273,1.252) (-1.239,1.241) (-1.205,1.230) (-1.170,1.218) (-1.137,1.206) (-1.103,1.193) (-1.069,1.180) (-1.036,1.166) (-1.003,1.151) (-0.971,1.136) (-0.938,1.120) (-0.906,1.104) (-0.875,1.086) (-0.844,1.068) (-0.813,1.049) (-0.783,1.029) (-0.754,1.008) (-0.725,0.987) (-0.697,0.964) (-0.670,0.940) (-0.644,0.915) (-0.620,0.889) (-0.596,0.862) (-0.574,0.833) (-0.554,0.803) (-0.536,0.772) (-0.520,0.740) (-0.506,0.707) (-0.500,0.684)}
\LC{(-2.300,1.635) (-2.264,1.632) (-2.228,1.628) (-2.193,1.624) (-2.157,1.621) (-2.121,1.617) (-2.085,1.613) (-2.049,1.609) (-2.014,1.604) (-1.978,1.600) (-1.942,1.595) (-1.907,1.591) (-1.871,1.586) (-1.835,1.581) (-1.800,1.576) (-1.764,1.570) (-1.728,1.565) (-1.693,1.559) (-1.657,1.553) (-1.622,1.547) (-1.587,1.540) (-1.551,1.534) (-1.516,1.527) (-1.480,1.520) (-1.445,1.512) (-1.410,1.505) (-1.375,1.497) (-1.340,1.489) (-1.305,1.480) (-1.270,1.472) (-1.235,1.463) (-1.200,1.453) (-1.166,1.444) (-1.131,1.434) (-1.096,1.424) (-1.062,1.413) (-1.028,1.402) (-0.993,1.391) (-0.959,1.380) (-0.925,1.368) (-0.892,1.355) (-0.858,1.343) (-0.824,1.330) (-0.791,1.316) (-0.758,1.302) (-0.725,1.288) (-0.692,1.274) (-0.659,1.259) (-0.626,1.244) (-0.594,1.228) (-0.562,1.212) (-0.529,1.196) (-0.498,1.179) (-0.466,1.162) (-0.434,1.145) (-0.403,1.127) (-0.371,1.110) (-0.340,1.092) (-0.309,1.073) (-0.278,1.055) (-0.247,1.037) (-0.216,1.019) (-0.184,1.001) (-0.153,0.984) (-0.120,0.968) (-0.087,0.954) (-0.053,0.944) (-0.018,0.938) (0.018,0.938) (0.054,0.945) (0.087,0.957) (0.120,0.973) (0.152,0.990) (0.182,1.009) (0.213,1.028) (0.243,1.047) (0.273,1.067) (0.304,1.086) (0.334,1.106) (0.364,1.125) (0.395,1.144) (0.426,1.163) (0.457,1.181) (0.488,1.199) (0.519,1.217) (0.551,1.234) (0.582,1.251) (0.614,1.268) (0.646,1.284) (0.679,1.300) (0.711,1.315) (0.744,1.330) (0.777,1.345) (0.810,1.359) (0.844,1.372) (0.877,1.385) (0.911,1.398) (0.945,1.411) (0.978,1.423) (1.013,1.434) (1.047,1.446) (1.081,1.456) (1.116,1.467) (1.150,1.477) (1.185,1.487) (1.219,1.496) (1.254,1.505) (1.289,1.514) (1.324,1.523) (1.359,1.531) (1.394,1.539) (1.429,1.547) (1.465,1.554) (1.500,1.561) (1.535,1.568) (1.571,1.575) (1.606,1.581) (1.642,1.587) (1.677,1.593) (1.713,1.599) (1.748,1.605) (1.784,1.610) (1.819,1.615) (1.855,1.620) (1.891,1.625) (1.926,1.630) (1.962,1.634) (1.998,1.638) (2.034,1.643) (2.069,1.647) (2.105,1.651) (2.141,1.654) (2.177,1.658) (2.213,1.662) (2.248,1.665) (2.284,1.668) (2.300,1.671)}
\LC{(-2.300,1.840) (-2.264,1.837) (-2.228,1.833) (-2.192,1.830) (-2.157,1.826) (-2.121,1.823) (-2.085,1.819) (-2.049,1.815) (-2.013,1.811) (-1.978,1.807) (-1.942,1.803) (-1.906,1.798) (-1.870,1.794) (-1.835,1.789) (-1.799,1.784) (-1.763,1.780) (-1.728,1.775) (-1.692,1.769) (-1.657,1.764) (-1.621,1.759) (-1.585,1.753) (-1.550,1.747) (-1.514,1.741) (-1.479,1.735) (-1.443,1.729) (-1.408,1.722) (-1.373,1.716) (-1.337,1.709) (-1.302,1.702) (-1.267,1.695) (-1.231,1.688) (-1.196,1.680) (-1.161,1.673) (-1.126,1.665) (-1.091,1.657) (-1.056,1.649) (-1.021,1.640) (-0.986,1.632) (-0.951,1.623) (-0.916,1.615) (-0.881,1.606) (-0.846,1.597) (-0.811,1.588) (-0.776,1.578) (-0.742,1.569) (-0.707,1.560) (-0.672,1.551) (-0.637,1.541) (-0.603,1.532) (-0.568,1.523) (-0.533,1.514) (-0.498,1.505) (-0.463,1.496) (-0.428,1.487) (-0.393,1.479) (-0.358,1.471) (-0.323,1.464) (-0.287,1.457) (-0.252,1.451) (-0.216,1.445) (-0.181,1.441) (-0.145,1.437) (-0.109,1.434) (-0.073,1.432) (-0.037,1.431) (-0.001,1.431) (0.035,1.433) (0.071,1.435) (0.107,1.439) (0.142,1.443) (0.178,1.448) (0.213,1.455) (0.249,1.461) (0.284,1.469) (0.319,1.477) (0.354,1.486) (0.389,1.495) (0.424,1.504) (0.458,1.514) (0.493,1.523) (0.528,1.533) (0.562,1.543) (0.597,1.554) (0.631,1.564) (0.666,1.574) (0.700,1.584) (0.735,1.594) (0.770,1.604) (0.804,1.614) (0.839,1.623) (0.874,1.633) (0.908,1.642) (0.943,1.651) (0.978,1.660) (1.013,1.669) (1.048,1.678) (1.083,1.686) (1.118,1.695) (1.153,1.703) (1.188,1.711) (1.223,1.718) (1.258,1.726) (1.294,1.733) (1.329,1.740) (1.364,1.747) (1.400,1.754) (1.435,1.761) (1.470,1.767) (1.506,1.773) (1.541,1.779) (1.577,1.785) (1.612,1.791) (1.648,1.796) (1.684,1.802) (1.719,1.807) (1.755,1.812) (1.791,1.817) (1.826,1.822) (1.862,1.826) (1.898,1.831) (1.933,1.835) (1.969,1.839) (2.005,1.844) (2.041,1.848) (2.076,1.851) (2.112,1.855) (2.148,1.859) (2.184,1.862) (2.220,1.866) (2.256,1.869) (2.291,1.872) (2.300,1.873)}
\LC{(0.050,-0.713) (0.077,-0.737) (0.107,-0.757) (0.138,-0.775) (0.170,-0.792) (0.203,-0.807) (0.236,-0.821) (0.269,-0.834) (0.303,-0.846) (0.337,-0.859) (0.371,-0.870) (0.405,-0.882) (0.439,-0.893) (0.474,-0.904) (0.508,-0.915) (0.542,-0.925) (0.577,-0.936) (0.612,-0.946) (0.646,-0.956) (0.681,-0.965) (0.716,-0.975) (0.750,-0.984) (0.785,-0.993) (0.820,-1.002) (0.855,-1.011) (0.890,-1.020) (0.925,-1.028) (0.960,-1.036) (0.995,-1.044) (1.030,-1.052) (1.065,-1.060) (1.100,-1.068) (1.136,-1.075) (1.171,-1.082) (1.206,-1.089) (1.242,-1.096) (1.277,-1.103) (1.312,-1.109) (1.348,-1.116) (1.383,-1.122) (1.419,-1.128) (1.454,-1.134) (1.490,-1.140) (1.525,-1.145) (1.561,-1.151) (1.597,-1.156) (1.632,-1.161) (1.668,-1.166) (1.703,-1.171) (1.739,-1.176) (1.775,-1.180) (1.811,-1.185) (1.846,-1.189) (1.882,-1.193) (1.918,-1.198) (1.954,-1.202) (1.989,-1.205) (2.025,-1.209) (2.061,-1.213) (2.097,-1.217) (2.133,-1.220) (2.169,-1.223) (2.204,-1.227) (2.240,-1.230) (2.276,-1.233) (2.300,-1.235)}
\LC{(0.193,-0.477) (0.224,-0.495) (0.257,-0.511) (0.288,-0.528) (0.320,-0.544) (0.353,-0.561) (0.385,-0.577) (0.417,-0.593) (0.450,-0.608) (0.482,-0.624) (0.515,-0.638) (0.548,-0.653) (0.581,-0.667) (0.614,-0.681) (0.648,-0.694) (0.681,-0.707) (0.715,-0.720) (0.749,-0.732) (0.783,-0.744) (0.817,-0.756) (0.851,-0.768) (0.885,-0.779) (0.919,-0.789) (0.954,-0.800) (0.988,-0.810) (1.023,-0.820) (1.058,-0.829) (1.093,-0.839) (1.127,-0.848) (1.162,-0.856) (1.197,-0.865) (1.232,-0.873) (1.267,-0.881) (1.303,-0.889) (1.338,-0.896) (1.373,-0.904) (1.408,-0.911) (1.444,-0.918) (1.479,-0.924) (1.515,-0.931) (1.550,-0.937) (1.585,-0.943) (1.621,-0.949) (1.657,-0.954) (1.692,-0.960) (1.728,-0.965) (1.763,-0.970) (1.799,-0.975) (1.835,-0.980) (1.870,-0.984) (1.906,-0.989) (1.942,-0.993) (1.978,-0.998) (2.013,-1.002) (2.049,-1.006) (2.085,-1.009) (2.121,-1.013) (2.157,-1.017) (2.192,-1.020) (2.228,-1.024) (2.264,-1.027) (2.300,-1.030)}
\LC{(0.341,-0.231) (0.372,-0.249) (0.401,-0.270) (0.430,-0.291) (0.460,-0.312) (0.490,-0.332) (0.520,-0.352) (0.550,-0.371) (0.581,-0.389) (0.612,-0.407) (0.644,-0.425) (0.675,-0.442) (0.707,-0.458) (0.740,-0.474) (0.772,-0.490) (0.805,-0.505) (0.838,-0.519) (0.871,-0.533) (0.904,-0.547) (0.938,-0.560) (0.972,-0.573) (1.005,-0.585) (1.040,-0.596) (1.074,-0.608) (1.108,-0.619) (1.142,-0.629) (1.177,-0.639) (1.212,-0.649) (1.246,-0.659) (1.281,-0.668) (1.316,-0.676) (1.351,-0.685) (1.386,-0.693) (1.421,-0.701) (1.457,-0.708) (1.492,-0.715) (1.527,-0.722) (1.562,-0.729) (1.598,-0.736) (1.633,-0.742) (1.669,-0.748) (1.704,-0.754) (1.740,-0.759) (1.776,-0.765) (1.811,-0.770) (1.847,-0.775) (1.882,-0.780) (1.918,-0.784) (1.954,-0.789) (1.990,-0.793) (2.025,-0.797) (2.061,-0.801) (2.097,-0.805) (2.133,-0.809) (2.169,-0.813) (2.204,-0.816) (2.240,-0.820) (2.276,-0.823) (2.300,-0.825)}
\LC{(0.525,-0.039) (0.550,-0.064) (0.576,-0.089) (0.603,-0.113) (0.630,-0.136) (0.658,-0.159) (0.687,-0.181) (0.716,-0.202) (0.746,-0.222) (0.776,-0.242) (0.807,-0.260) (0.838,-0.278) (0.870,-0.296) (0.902,-0.312) (0.934,-0.328) (0.967,-0.343) (1.000,-0.358) (1.033,-0.371) (1.066,-0.385) (1.100,-0.398) (1.134,-0.410) (1.168,-0.421) (1.202,-0.433) (1.237,-0.443) (1.271,-0.453) (1.306,-0.463) (1.340,-0.473) (1.375,-0.481) (1.410,-0.490) (1.445,-0.498) (1.480,-0.506) (1.516,-0.514) (1.551,-0.521) (1.586,-0.528) (1.622,-0.534) (1.657,-0.541) (1.693,-0.547) (1.728,-0.553) (1.764,-0.558) (1.799,-0.564) (1.835,-0.569) (1.871,-0.574) (1.906,-0.578) (1.942,-0.583) (1.978,-0.587) (2.013,-0.592) (2.049,-0.596) (2.085,-0.600) (2.121,-0.603) (2.157,-0.607) (2.192,-0.610) (2.228,-0.614) (2.264,-0.617) (2.300,-0.620)}
\LC{(0.694,0.096) (0.717,0.068) (0.741,0.042) (0.766,0.016) (0.793,-0.008) (0.821,-0.031) (0.849,-0.052) (0.879,-0.073) (0.909,-0.093) (0.940,-0.111) (0.971,-0.129) (1.003,-0.145) (1.036,-0.161) (1.068,-0.176) (1.102,-0.190) (1.135,-0.204) (1.169,-0.216) (1.203,-0.229) (1.237,-0.240) (1.271,-0.251) (1.306,-0.261) (1.340,-0.271) (1.375,-0.280) (1.410,-0.289) (1.445,-0.297) (1.480,-0.305) (1.515,-0.313) (1.550,-0.320) (1.586,-0.327) (1.621,-0.334) (1.657,-0.340) (1.692,-0.346) (1.728,-0.352) (1.763,-0.357) (1.799,-0.362) (1.835,-0.367) (1.870,-0.372) (1.906,-0.376) (1.942,-0.381) (1.977,-0.385) (2.013,-0.389) (2.049,-0.393) (2.085,-0.396) (2.121,-0.400) (2.157,-0.403) (2.192,-0.406) (2.228,-0.409) (2.264,-0.412) (2.300,-0.415)}
\LC{(0.814,0.196) (0.840,0.170) (0.866,0.146) (0.894,0.123) (0.922,0.101) (0.952,0.080) (0.982,0.061) (1.014,0.043) (1.046,0.026) (1.078,0.011) (1.111,-0.004) (1.144,-0.018) (1.178,-0.030) (1.212,-0.042) (1.246,-0.054) (1.280,-0.064) (1.315,-0.074) (1.350,-0.084) (1.385,-0.092) (1.420,-0.101) (1.455,-0.108) (1.490,-0.116) (1.525,-0.123) (1.561,-0.129) (1.596,-0.135) (1.632,-0.141) (1.667,-0.147) (1.703,-0.152) (1.739,-0.157) (1.774,-0.162) (1.810,-0.166) (1.846,-0.170) (1.882,-0.175) (1.917,-0.178) (1.953,-0.182) (1.989,-0.186) (2.025,-0.189) (2.061,-0.192) (2.097,-0.195) (2.132,-0.198) (2.168,-0.201) (2.204,-0.203) (2.240,-0.206) (2.276,-0.208) (2.300,-0.210)}
\LC{(0.911,0.282) (0.938,0.258) (0.967,0.236) (0.997,0.217) (1.028,0.199) (1.061,0.183) (1.094,0.169) (1.127,0.156) (1.161,0.144) (1.195,0.132) (1.230,0.122) (1.264,0.113) (1.299,0.104) (1.334,0.096) (1.370,0.088) (1.405,0.081) (1.440,0.075) (1.476,0.069) (1.511,0.063) (1.547,0.058) (1.583,0.053) (1.618,0.048) (1.654,0.044) (1.690,0.039) (1.726,0.036) (1.761,0.032) (1.797,0.028) (1.833,0.025) (1.869,0.022) (1.905,0.019) (1.941,0.016) (1.977,0.014) (2.012,0.011) (2.048,0.009) (2.084,0.006) (2.120,0.004) (2.156,0.002) (2.192,0.000) (2.228,-0.002) (2.264,-0.003) (2.300,-0.005)}
\LC{(0.973,0.353) (1.005,0.337) (1.038,0.323) (1.072,0.311) (1.107,0.300) (1.141,0.291) (1.176,0.283) (1.212,0.276) (1.247,0.269) (1.283,0.263) (1.318,0.258) (1.354,0.253) (1.390,0.249) (1.425,0.245) (1.461,0.241) (1.497,0.238) (1.533,0.235) (1.569,0.232) (1.605,0.229) (1.641,0.227) (1.677,0.224) (1.712,0.222) (1.748,0.220) (1.784,0.218) (1.820,0.216) (1.856,0.215) (1.892,0.213) (1.928,0.212) (1.964,0.210) (2.000,0.209) (2.036,0.207) (2.072,0.206) (2.108,0.205) (2.144,0.204) (2.180,0.203) (2.216,0.202) (2.252,0.201) (2.288,0.200) (2.300,0.200)}
\LC{(0.994,0.428) (1.029,0.419) (1.064,0.415) (1.100,0.412) (1.136,0.410) (1.172,0.408) (1.208,0.407) (1.244,0.406) (1.280,0.406) (1.316,0.405) (1.352,0.405) (1.388,0.404) (1.424,0.404) (1.460,0.404) (1.496,0.404) (1.532,0.404) (1.568,0.404) (1.604,0.404) (1.640,0.404) (1.676,0.404) (1.712,0.404) (1.748,0.404) (1.784,0.404) (1.820,0.404) (1.856,0.404) (1.892,0.404) (1.928,0.404) (1.964,0.404) (2.000,0.404) (2.036,0.404) (2.072,0.405) (2.108,0.405) (2.144,0.405) (2.180,0.405) (2.216,0.405) (2.252,0.405) (2.288,0.405) (2.300,0.405)}
\LC{(0.991,0.477) (1.024,0.492) (1.059,0.502) (1.094,0.511) (1.129,0.519) (1.164,0.526) (1.200,0.532) (1.235,0.538) (1.271,0.543) (1.306,0.548) (1.342,0.553) (1.378,0.557) (1.414,0.561) (1.449,0.565) (1.485,0.568) (1.521,0.571) (1.557,0.574) (1.593,0.577) (1.629,0.580) (1.665,0.582) (1.701,0.585) (1.736,0.587) (1.772,0.589) (1.808,0.591) (1.844,0.593) (1.880,0.595) (1.916,0.596) (1.952,0.598) (1.988,0.599) (2.024,0.601) (2.060,0.602) (2.096,0.604) (2.132,0.605) (2.168,0.606) (2.204,0.607) (2.240,0.608) (2.276,0.609) (2.300,0.610)}
\LC{(0.956,0.539) (0.983,0.563) (1.012,0.585) (1.042,0.603) (1.074,0.620) (1.107,0.636) (1.140,0.649) (1.174,0.662) (1.208,0.673) (1.242,0.684) (1.277,0.693) (1.312,0.702) (1.347,0.711) (1.382,0.718) (1.417,0.726) (1.453,0.732) (1.488,0.739) (1.524,0.744) (1.559,0.750) (1.595,0.755) (1.630,0.760) (1.666,0.764) (1.702,0.769) (1.738,0.773) (1.773,0.777) (1.809,0.780) (1.845,0.784) (1.881,0.787) (1.917,0.790) (1.953,0.793) (1.989,0.796) (2.025,0.798) (2.060,0.801) (2.096,0.803) (2.132,0.806) (2.168,0.808) (2.204,0.810) (2.240,0.812) (2.276,0.814) (2.300,0.815)}
\LC{(0.873,0.589) (0.893,0.619) (0.915,0.648) (0.939,0.674) (0.965,0.699) (0.993,0.722) (1.022,0.743) (1.052,0.762) (1.083,0.780) (1.115,0.797) (1.148,0.813) (1.181,0.827) (1.214,0.840) (1.248,0.853) (1.282,0.864) (1.317,0.875) (1.351,0.885) (1.386,0.895) (1.421,0.904) (1.456,0.912) (1.491,0.920) (1.526,0.927) (1.561,0.934) (1.597,0.941) (1.632,0.947) (1.668,0.953) (1.703,0.959) (1.739,0.964) (1.775,0.969) (1.810,0.974) (1.846,0.978) (1.882,0.982) (1.917,0.987) (1.953,0.990) (1.989,0.994) (2.025,0.998) (2.061,1.001) (2.097,1.004) (2.132,1.007) (2.168,1.010) (2.204,1.013) (2.240,1.016) (2.276,1.018) (2.300,1.020)}
\LC{(0.752,0.644) (0.764,0.678) (0.780,0.710) (0.799,0.741) (0.821,0.770) (0.844,0.798) (0.868,0.824) (0.895,0.848) (0.922,0.871) (0.951,0.893) (0.981,0.914) (1.011,0.933) (1.042,0.951) (1.074,0.968) (1.106,0.984) (1.139,0.999) (1.172,1.013) (1.205,1.027) (1.239,1.039) (1.273,1.051) (1.307,1.062) (1.342,1.073) (1.376,1.083) (1.411,1.092) (1.446,1.101) (1.481,1.110) (1.516,1.118) (1.551,1.126) (1.586,1.133) (1.622,1.140) (1.657,1.146) (1.692,1.153) (1.728,1.159) (1.764,1.164) (1.799,1.170) (1.835,1.175) (1.870,1.180) (1.906,1.185) (1.942,1.189) (1.978,1.193) (2.013,1.198) (2.049,1.202) (2.085,1.205) (2.121,1.209) (2.157,1.212) (2.192,1.216) (2.228,1.219) (2.264,1.222) (2.300,1.225)}
\LC{(0.559,0.710) (0.569,0.745) (0.582,0.778) (0.598,0.811) (0.616,0.842) (0.636,0.872) (0.658,0.900) (0.682,0.927) (0.707,0.954) (0.733,0.978) (0.760,1.002) (0.788,1.024) (0.817,1.046) (0.847,1.066) (0.877,1.086) (0.908,1.104) (0.939,1.122) (0.971,1.138) (1.004,1.154) (1.036,1.169) (1.069,1.184) (1.102,1.197) (1.136,1.211) (1.170,1.223) (1.204,1.235) (1.238,1.246) (1.272,1.257) (1.307,1.267) (1.341,1.277) (1.376,1.287) (1.411,1.295) (1.446,1.304) (1.481,1.312) (1.516,1.320) (1.551,1.328) (1.587,1.335) (1.622,1.342) (1.657,1.348) (1.693,1.354) (1.728,1.360) (1.764,1.366) (1.799,1.372) (1.835,1.377) (1.871,1.382) (1.906,1.387) (1.942,1.392) (1.978,1.396) (2.014,1.401) (2.049,1.405) (2.085,1.409) (2.121,1.413) (2.157,1.416) (2.192,1.420) (2.228,1.424) (2.264,1.427) (2.300,1.430)}

      % Conductor neutro en campo uniforme - flechas
% Generado por scripts/generar_datos.py (no editar a mano).
\FL{-1.463}{-1.775}{6.2}
\FL{1.475}{-1.759}{-5.7}
\FL{-1.464}{-1.560}{7.3}
\FL{1.476}{-1.543}{-6.8}
\FL{-1.466}{-1.345}{8.6}
\FL{1.472}{-1.326}{-8.0}
\FL{-1.906}{-1.195}{6.9}
\FL{-0.562}{-0.896}{20.8}
\FL{-1.930}{-0.991}{7.1}
\FL{-0.688}{-0.683}{23.7}
\FL{-1.954}{-0.789}{7.1}
\FL{-0.778}{-0.473}{26.9}
\FL{-1.978}{-0.588}{6.7}
\FL{-0.885}{-0.290}{29.8}
\FL{-2.001}{-0.389}{5.9}
\FL{-1.016}{-0.149}{28.7}
\FL{-2.037}{-0.193}{4.5}
\FL{-1.121}{-0.019}{23.9}
\FL{-2.048}{0.005}{2.8}
\FL{-1.204}{0.106}{15.7}
\FL{-2.060}{0.202}{0.7}
\FL{-1.245}{0.235}{6.2}
\FL{-2.060}{0.400}{-1.4}
\FL{-1.257}{0.365}{-4.2}
\FL{-2.060}{0.598}{-3.4}
\FL{-1.228}{0.493}{-14.5}
\FL{-2.037}{0.795}{-5.2}
\FL{-1.158}{0.616}{-24.2}
\FL{-2.013}{0.992}{-6.5}
\FL{-1.048}{0.730}{-32.8}
\FL{-1.978}{1.190}{-7.4}
\FL{-0.904}{0.846}{-38.5}
\FL{-1.930}{1.388}{-7.9}
\FL{-0.707}{0.972}{-39.7}
\FL{-1.422}{1.507}{-12.2}
\FL{1.441}{1.549}{11.9}
\FL{-1.455}{1.731}{-10.1}
\FL{1.470}{1.767}{10.1}
\FL{0.428}{-0.889}{-18.1}
\FL{1.894}{-1.195}{-6.6}
\FL{0.548}{-0.653}{-23.4}
\FL{1.918}{-0.990}{-7.0}
\FL{0.665}{-0.436}{-28.1}
\FL{1.930}{-0.786}{-7.2}
\FL{0.797}{-0.254}{-31.1}
\FL{1.966}{-0.586}{-6.9}
\FL{0.930}{-0.105}{-30.7}
\FL{2.001}{-0.387}{-6.3}
\FL{1.046}{0.026}{-26.7}
\FL{2.025}{-0.189}{-5.2}
\FL{1.138}{0.152}{-20.0}
\FL{2.048}{0.009}{-3.7}
\FL{1.212}{0.276}{-11.0}
\FL{2.060}{0.207}{-1.9}
\FL{1.232}{0.407}{-1.3}
\FL{2.072}{0.405}{0.1}
\FL{1.223}{0.536}{9.1}
\FL{2.072}{0.603}{2.1}
\FL{1.174}{0.662}{19.4}
\FL{2.060}{0.801}{3.9}
\FL{1.083}{0.780}{28.7}
\FL{2.025}{0.998}{5.5}
\FL{0.941}{0.886}{36.8}
\FL{2.001}{1.196}{6.5}
\FL{0.760}{1.002}{39.8}
\FL{1.954}{1.393}{7.3}

    \end{scope}
    % Contorno del conductor irregular
% Generado por scripts/generar_datos.py (no editar a mano).
\def\formairregular{(-0.534,0.681) -- (-0.416,0.719) -- (-0.297,0.748) -- (-0.178,0.769) -- (-0.059,0.782) -- (0.059,0.786) -- (0.178,0.782) -- (0.297,0.770) -- (0.416,0.749) -- (0.534,0.721) -- (0.640,0.691) -- (0.732,0.661) -- (0.811,0.631) -- (0.877,0.599) -- (0.929,0.568) -- (0.968,0.535) -- (0.994,0.502) -- (1.006,0.468) -- (1.006,0.431) -- (0.994,0.390) -- (0.969,0.345) -- (0.931,0.297) -- (0.881,0.245) -- (0.819,0.190) -- (0.744,0.131) -- (0.656,0.069) -- (0.575,0.004) -- (0.500,-0.063) -- (0.431,-0.133) -- (0.369,-0.205) -- (0.312,-0.279) -- (0.262,-0.355) -- (0.219,-0.434) -- (0.181,-0.516) -- (0.144,-0.586) -- (0.106,-0.647) -- (0.069,-0.697) -- (0.031,-0.736) -- (-0.006,-0.765) -- (-0.044,-0.783) -- (-0.081,-0.791) -- (-0.119,-0.789) -- (-0.156,-0.777) -- (-0.194,-0.755) -- (-0.231,-0.724) -- (-0.269,-0.683) -- (-0.306,-0.633) -- (-0.344,-0.574) -- (-0.381,-0.504) -- (-0.419,-0.426) -- (-0.460,-0.350) -- (-0.505,-0.279) -- (-0.555,-0.211) -- (-0.608,-0.147) -- (-0.665,-0.086) -- (-0.726,-0.029) -- (-0.791,0.025) -- (-0.859,0.075) -- (-0.917,0.123) -- (-0.964,0.170) -- (-1.000,0.216) -- (-1.025,0.259) -- (-1.039,0.302) -- (-1.042,0.342) -- (-1.034,0.381) -- (-1.016,0.419) -- (-0.984,0.456) -- (-0.941,0.494) -- (-0.884,0.531) -- (-0.816,0.569) -- (-0.734,0.606) -- (-0.641,0.644) -- cycle}

    \fill[top color=azulvelo!80!white, bottom color=metal, draw=azulprof, line width=0.6pt] \formairregular;
    % Conductor neutro en campo uniforme - cargas
% Generado por scripts/generar_datos.py (no editar a mano).
\CP{0.223}{-0.629}
\CP{0.442}{-0.257}
\CP{0.699}{-0.012}
\CP{0.903}{0.145}
\CP{1.034}{0.283}
\CP{1.096}{0.459}
\CP{0.999}{0.626}
\CP{0.820}{0.724}
\CP{0.492}{0.823}
\CN{-0.744}{0.699}
\CN{-0.969}{0.583}
\CN{-1.107}{0.436}
\CN{-1.119}{0.256}
\CN{-1.016}{0.095}
\CN{-0.848}{-0.044}
\CN{-0.635}{-0.255}
\CN{-0.448}{-0.570}
\CN{-0.273}{0.844}

    \node[font=\footnotesize, text=azulprof] at (-0.15,0.18) {$\vE=\vc{0}$};
    \draw[azulclaro, line width=0.5pt] (-2.3,-1.9) rectangle (2.3,1.9);
    \node[anchor=north] at (0,-2.0) {(b) Conductor irregular};
  \end{scope}
  % flecha del campo aplicado
  \draw[-{Stealth[length=5pt]}, azulprof, line width=0.9pt] (4.95,2.95) -- (6.25,2.95)
    node[midway, above, font=\small] {$\vE_0$};
\end{tikzpicture}
\caption{Conductores neutros en un campo uniforme $\vE_0$ (de izquierda a derecha). Líneas de campo (continuas) y equipotenciales (trazos).
(a) Esfera: la solución exacta, $\phi=-E_0\bigl(r-a^{3}/r^{2}\bigr)\cos\theta$, induce $\sigma=3\eo E_0\cos\theta$. (b) Conductor de forma
arbitraria, calculado numéricamente resolviendo $\Lap\phi=0$ con la condición de carga neta nula. En ambos casos las líneas llegan
perpendiculares a la superficie y el campo interior es nulo.}
\label{fig:uniforme}
\end{figure}


% -------------------------------------------------------------------------
\subsection{Notación y convenios}\label{sec:notacion}
% -------------------------------------------------------------------------
\begin{notacion}\label{not:tema3}
A lo largo del tema usamos la notación de los temas anteriores~\cite{TEMA0,TEMA1,TEMA2}, y además:
\begin{itemize}
  \item $\Vc$ es la región (cerrada) ocupada por un conductor y $\Sc=\partial\Vc$ su superficie; $\un$ es la normal unitaria \emph{exterior} al
        conductor, es decir, la que apunta del metal al vacío. Si el conductor tiene cavidades, $\Sc$ incluye sus caras interiores, y en ellas
        $\un$ apunta hacia la cavidad.
  \item $\sigma(\vc{r})$ es la densidad superficial de carga sobre $\Sc$, y $\Qint$, $\Qext$ las cargas de las caras interior y exterior de un
        conductor hueco.
  \item En un sistema de $N$ conductores, $Q_i$ y $\phi_i$ son la carga total y el potencial del conductor $i$; los agrupamos en los vectores
        columna $\vQ=(Q_1,\dots,Q_N)\Tt$ y $\vphi=(\phi_1,\dots,\phi_N)\Tt$. Las matrices de coeficientes de potencial y de capacidad e influencia
        son $\Pm=(p_{ij})$ y $\Cm=(c_{ij})$.
  \item Mientras no se diga lo contrario, el origen de potenciales está en el infinito, $\phi(\infty)=0$, y escribimos
        $\kc\equiv1/(4\pi\eo)\simeq8{,}988\cdot10^{9}\,\mathrm{N\,m^{2}\,C^{-2}}$~\cite{CODATA}.
  \item «A tierra» significa $\phi=0$ (la Tierra es un conductor tan grande que su potencial no cambia al intercambiar carga con el sistema);
        «aislado» significa que la carga del conductor se conserva.
\end{itemize}
\end{notacion}
